\documentclass[11pt,oneside,openany]{book}
\usepackage[T1]{fontenc}
\usepackage[utf8]{inputenc}
\usepackage[sc]{mathpazo}
\renewcommand{\sfdefault}{lmss}
\renewcommand{\ttdefault}{lmtt}
\usepackage{amsmath,amssymb,amsthm,mathtools}
\usepackage[paperwidth=7.5in,paperheight=10in,inner=.85in,outer=.85in,top=.85in,bottom=.85in,headheight=15pt,headsep=20pt]{geometry}
\usepackage{microtype}
\usepackage{xcolor,graphicx,booktabs,array,longtable}
\newcolumntype{L}[1]{>{\raggedright\arraybackslash}p{#1}}
\usepackage{tikz,pgfplots}
\usetikzlibrary{arrows.meta,calc,patterns,decorations.markings}
\pgfplotsset{compat=1.18}
\definecolor{ink}{HTML}{17334B}
\definecolor{teal}{HTML}{137C83}
\definecolor{copper}{HTML}{AE5932}
\definecolor{soft}{HTML}{EEF3F5}
\pgfplotsset{every axis/.append style={axis line style={black!65},tick style={black!60},tick label style={font=\small},label style={font=\small},legend style={font=\small,draw=none,fill=white},grid style={black!10},major grid style={black!10},line width=.6pt},every axis plot/.append style={line width=1.15pt}}
\usepackage{enumitem,needspace,etoolbox}
\setlist[enumerate]{topsep=5pt,itemsep=5pt,parsep=0pt}
\setlist[itemize]{topsep=4pt,itemsep=3pt,parsep=0pt}
\usepackage{titlesec,fancyhdr}
\titleformat{\chapter}[display]{\normalfont\color{ink}}{\large\sffamily\MakeUppercase{Chapter \thechapter}}{12pt}{\huge\bfseries\raggedright\hyphenpenalty=10000\exhyphenpenalty=10000}[\vspace{12pt}\titlerule]
\titlespacing*{\chapter}{0pt}{-10pt}{24pt}
\titleformat{\section}{\large\bfseries\color{ink}}{\thesection}{.65em}{}
\titleformat{\subsection}{\normalsize\bfseries\color{ink}}{\thesubsection}{.6em}{}
\titlespacing*{\section}{0pt}{18pt plus 3pt minus 2pt}{7pt}
\titlespacing*{\subsection}{0pt}{12pt plus 2pt minus 2pt}{5pt}
\pagestyle{fancy}
\fancyhf{}
\fancyhead[L]{\small\nouppercase{\leftmark}}
\fancyfoot[C]{\small\thepage}
\renewcommand{\headrulewidth}{.3pt}
\renewcommand{\chaptermark}[1]{\markboth{\thechapter\quad #1}{}}
\fancypagestyle{plain}{\fancyhf{}\fancyfoot[C]{\small\thepage}\renewcommand{\headrulewidth}{0pt}}
\usepackage[font=small,labelfont=bf,labelsep=period]{caption}
\usepackage{float,placeins}
\floatplacement{figure}{htbp}
\setcounter{topnumber}{3}
\setcounter{bottomnumber}{2}
\renewcommand{\topfraction}{.9}
\renewcommand{\textfraction}{.08}
\renewcommand{\floatpagefraction}{.8}
\usepackage{makeidx}
\usepackage[columns=3,font=small,indentunit=11pt,columnsep=14pt]{idxlayout}
\makeindex
\usepackage[pdfencoding=auto,psdextra,colorlinks=true,linkcolor=ink,citecolor=teal,urlcolor=teal,bookmarksnumbered=true]{hyperref}
\usepackage{bookmark}
\hypersetup{pdftitle={The Calculus of Probability Distributions},pdfsubject={An undergraduate introduction to information geometry}}
\urlstyle{same}
\newtheorem{theorem}{Theorem}[chapter]
\newtheorem{proposition}[theorem]{Proposition}
\newtheorem{lemma}[theorem]{Lemma}
\newtheorem{corollary}[theorem]{Corollary}
\theoremstyle{definition}
\newtheorem{definition}[theorem]{Definition}
\newtheorem{example}{Example}[chapter]
% A worked example may continue naturally across a page break.  Space, rather
% than a heavy colored box, distinguishes it from the surrounding discussion.
\BeforeBeginEnvironment{example}{\par\addvspace{9pt}\needspace{4\baselineskip}}
\AfterEndEnvironment{example}{\par\addvspace{8pt}}
\theoremstyle{remark}
\newtheorem*{remark}{Remark}
\newenvironment{solution}{\par\smallskip\noindent\textit{Solution.}\ }{\par\smallskip}
\newenvironment{exercises}{\begin{enumerate}[label=\thechapter.\arabic*,ref=\thechapter.\arabic*,leftmargin=2.8em,itemsep=6pt]}{\end{enumerate}}
\newenvironment{answers}{\begin{description}[leftmargin=2.9em,labelwidth=2.5em,labelsep=.4em,itemsep=8pt,font=\normalfont\bfseries]}{\end{description}}
\newcommand{\R}{\mathbb{R}}
\newcommand{\E}{\mathbb{E}}
\newcommand{\Var}{\operatorname{Var}}
\newcommand{\Cov}{\operatorname{Cov}}
\newcommand{\KL}[2]{D\!\left(#1\,\middle\|\,#2\right)}
\newcommand{\grad}{\operatorname{grad}}
\newcommand{\arcosh}{\operatorname{arcosh}}
\newcommand{\sech}{\operatorname{sech}}
\newcommand{\softmax}{\operatorname{softmax}}
\newcommand{\diag}{\operatorname{diag}}
\newcommand{\dd}{\,\mathrm{d}}
\newcommand{\ip}[2]{\left\langle #1,#2\right\rangle}
\newcommand{\p}{\mathbf{p}}
\newcommand{\q}{\mathbf{q}}
\newcommand{\f}{\mathbf{f}}
\newcommand{\uvec}{\mathbf{u}}
\newcommand{\vvec}{\mathbf{v}}
\newcommand{\ones}{\mathbf{1}}
\newcommand{\simp}{\Delta^{n-1}}
\newcommand{\simpi}{\Delta^{n-1}_{\!\circ}}
\numberwithin{equation}{chapter}
\setlength{\parindent}{1.2em}
\setlength{\parskip}{0pt}
\linespread{1.115}
\setlength{\emergencystretch}{2em}
\clubpenalty=10000
\widowpenalty=10000
\displaywidowpenalty=10000
\raggedbottom
\setcounter{tocdepth}{1}
\setcounter{secnumdepth}{2}
% Leave enough room for the four-character subsection numbers in Chapter 10.
\makeatletter
\renewcommand*\l@section{\@dottedtocline{1}{1.5em}{3.35em}}
\makeatother

\begin{document}
\begin{titlepage}
\thispagestyle{empty}
\begin{tikzpicture}[remember picture,overlay]
\fill[ink] (current page.north west) rectangle ([yshift=-.33in]current page.north east);
\draw[teal,line width=1.3pt] ([xshift=.85in,yshift=-4.75in]current page.north west)--([xshift=-.85in,yshift=-4.75in]current page.north east);
\end{tikzpicture}
\vspace*{.5in}
{\Large\sffamily\color{teal} INFORMATION GEOMETRY\par}
\vspace{.55in}
{\fontsize{32}{38}\selectfont\bfseries\color{ink}The Calculus of\\Probability\\Distributions\par}
\vspace{.28in}
{\Large An Undergraduate Introduction\par}
\vspace{.05in}
{\large\color{ink} Christopher Lee\par}
\vspace{.66in}
\begin{center}
\begin{tikzpicture}[scale=1.05]
\coordinate (A) at (0,0);\coordinate (B) at (2.7,0);\coordinate (C) at (1.35,2.34);
\fill[soft] (A)--(B)--(C)--cycle;
\draw[ink,line width=.85pt] (A)--(B)--(C)--cycle;
\draw[teal,line width=1.5pt] (.45,.35) to[bend left=20] (1.85,.95);
\fill[teal] (.45,.35) circle(2pt);\fill[copper] (1.85,.95) circle(2pt);
\draw[->,ink!65,line width=.8pt] (3.2,1.1)--(4.2,1.1);
\draw[ink,line width=.9pt] (5.4,1.1) circle(1.18);
\draw[ink!35] (4.22,1.1) arc(180:360:1.18 and .35);
\draw[ink!35,dashed] (6.58,1.1) arc(0:180:1.18 and .35);
\draw[teal,line width=1.6pt] (4.58,.72) to[bend left=22] (5.95,1.95);
\fill[teal] (4.58,.72) circle(2pt);\fill[copper] (5.95,1.95) circle(2pt);
\end{tikzpicture}
\end{center}
\vfill
{\large Dynamics, probability, information, and geometry\par}
\vspace{.18in}
{\normalsize First edition\par}
\end{titlepage}
\frontmatter
\chapter*{Preface}
\addcontentsline{toc}{chapter}{Preface}
\markboth{Preface}{}
A probability distribution can be a state of a dynamical system, a model for observations, and a point in a geometric space. These interpretations meet in information geometry. Its central construction, the Fisher information metric, measures the local sensitivity of a distribution to changes in its parameters. From that metric come lengths, angles, shortest paths, and a notion of gradient adapted to statistical models.

The development begins with a differential equation for a population proportion. Two competing types lead to a scalar replicator equation; a third type leads to motion on a probability triangle. Studying the same states as probability distributions introduces expectation, entropy, and relative entropy. A Taylor expansion of relative entropy then produces the Fisher metric. Square-root coordinates reveal spherical geometry for categorical distributions, and natural gradients give a geometric interpretation of potential-driven replicator systems. The final chapter develops the hyperbolic geometry of the Gaussian family.

Several examples recur as the theory develops. In particular, a two-type coexistence equation first appears as a phase-line problem and later as an entropy calculation, a curve in Fisher geometry, and the gradient flow of a potential. Comparing these descriptions reveals what each mathematical structure contributes. The same idea extends to three-type interactions and to families of continuously varying probability densities.

The prerequisites are multivariable calculus, elementary linear algebra, and a first course in differential equations. Probability is developed from finite distributions, with continuous densities introduced in Chapter 10. The geometric language is introduced through explicit inner products and coordinate calculations. Familiarity with abstract manifolds or measure-theoretic probability is unnecessary.

Each chapter contains derivations, worked examples, figures, and exercises. The exercises include direct calculations, interpretation of models, short proofs, and numerical investigations. Problems labeled \emph{Exploration} invite the use of a graphing tool or the corresponding interactive course demonstration. Their mathematical questions can also be studied independently. Selected solutions appear at the end of the book.

The figures are intended to be read alongside the equations. A phase-line arrow identifies the sign of a differential equation; a path in a probability simplex represents changing category proportions; and a geodesic depicts length measured by the Fisher metric. Returning to a figure after completing its associated calculation can clarify which features belong to a coordinate drawing and which belong to the underlying model.

The chapters follow the ten-part course \emph{Introduction to Information Geometry}. Its interactive demonstrations provide synchronized views of the models and geometric constructions:
\begin{center}
\url{https://bobovski66.github.io/IntroInfoGeom/}
\end{center}
In a ten-week course, one chapter corresponds to one week. A longer course can develop the proof exercises and computational investigations in greater depth. Chapters 1--3 supply the dynamical foundation; Chapters 4--6 develop probability and information; Chapters 7--10 construct and apply the geometry.

Natural logarithms are used throughout, so information is measured in nats unless bits are explicitly requested. The Fisher metric is normalized as the expected product of score functions. With this convention the categorical square-root sphere has radius 2, and the Gaussian half-plane has curvature $-1/2$. Keeping the normalization fixed allows calculations made in different chapters to be compared directly.

\tableofcontents
\chapter*{Notation and Conventions}
\addcontentsline{toc}{chapter}{Notation and Conventions}
\markboth{Notation and Conventions}{}
{\renewcommand{\arraystretch}{.88}
\begin{longtable}{@{}L{.27\textwidth}L{.67\textwidth}@{}}
\toprule
Symbol & Meaning\\\midrule\endhead
$\p=(p_1,\ldots,p_n)$ & Probability vector; column-vector convention in matrix products.\\[5pt]
$\Delta^{n-1}$ & Closed simplex $\{\p\in\R^n:p_i\geq0,\ \sum_i p_i=1\}$.\\[5pt]
$\Delta^{n-1}_{\!\circ}$ & Positive interior of the simplex, relative to its normalization plane.\\[5pt]
$T_\p\Delta^{n-1}_{\!\circ}$ & Tangent space $\{\vvec\in\R^n:\sum_i v_i=0\}$.\\[5pt]
$f_i,\ \bar f$ & Fitness of type $i$ and mean fitness $\sum_i p_i f_i$.\\[5pt]
$M$ & Payoff matrix; $m_{ij}$ is payoff to type $i$ against type $j$.\\[5pt]
$\E_\p[Z]$ & Expectation of a random variable under distribution $\p$.\\[5pt]
$\Var_\p(Z)$ & Variance $\E_\p[(Z-\E_\p Z)^2]$.\\[5pt]
$H(P)$ & Shannon entropy $-\sum_i p_i\ln p_i$.\\[5pt]
$D(Q\|P)$ & KL divergence $\sum_i q_i\ln(q_i/p_i)$; the first distribution supplies the expectation.\\[5pt]
$g_\p(u,v)$ & Fisher inner product of tangent vectors.\\[5pt]
$G(\theta),\ I(\theta)$ & Fisher information matrix; scalar information for one parameter.\\[5pt]
$L_F(\gamma),\ d_F(P,R)$ & Fisher length of a curve and intrinsic Fisher distance between its endpoints.\\[5pt]
$\grad_F\Phi$ & Gradient of a scalar function with respect to the Fisher metric.\\[5pt]
$N(\mu,\sigma^2)$ & Normal distribution with mean $\mu$ and standard deviation $\sigma>0$.\\[5pt]
$O(\|h\|^k)$ & Remainder bounded in magnitude by a constant times $\|h\|^k$ near $h=0$.\\
\bottomrule
\end{longtable}
}
For a two-outcome distribution, $D(q\|p)$ abbreviates $D((q,1-q)\|(p,1-p))$. A Bernoulli parameter $p$ denotes the probability of outcome 1; ordering the vector as $(p,1-p)$ or $(1-p,p)$ leaves every geometric calculation unchanged if that ordering is used consistently. All tangent-space formulas are stated on positive distributions. Boundary states are included explicitly when limits, invariant faces, or metric completion are discussed.
\mainmatter

\chapter{Replicator Equations and Scalar Dynamics}
\label{ch:scalar}

A population can grow while one of its types becomes less common. To describe this possibility, a model must keep track of relative abundance as well as absolute size. The differential equation for relative abundance provides a first example of dynamics on a space of probability distributions.

The distinction is already visible in two observations of a population: its count may increase from one census to the next while its share of the total decreases. Composition is therefore a useful state variable in its own right. Once the equation for composition is known, a single point on an interval describes the population's state, and a moving point describes its evolution.

\section{Population size and population composition}\index{proportion}\index{relative abundance}
Let $N_A(t)$ and $N_B(t)$ denote the sizes of two populations. In a continuous model these variables are nonnegative real numbers, even when actual observations are integer counts. Suppose their per-capita growth rates are $f_A$ and $f_B$:
\begin{equation}\label{eq:populations}
N_A'=f_A N_A,\qquad N_B'=f_B N_B.
\end{equation}
The rate $f_A$ has units of inverse time: a value of $0.2$ per day means an instantaneous growth rate equal to $20\%$ of the current A population per day. The rates can depend on the current composition. A negative rate represents declining abundance.

Whenever $N=N_A+N_B>0$, define
\[
x=\frac{N_A}{N},\qquad 1-x=\frac{N_B}{N}.
\]
The variable $x$ records the proportion of type A. Its state space is the interval $[0,1]$. The endpoints correspond to populations containing a single type; interior points correspond to mixtures.

\begin{example}[Increasing count and decreasing share]
A population changes from 48 A and 72 B individuals to 90 A and 180 B individuals. Determine the change in A's abundance and proportion.
\begin{solution}
The totals are $120$ and $270$. Type A increases from 48 to 90 individuals, while its proportion changes from $48/120=2/5$ to $90/270=1/3$. The total population grows faster, proportionally, than the A population. The count and the proportion therefore describe different features of the same observations.
\end{solution}
\end{example}

Multiplying both population sizes by the same positive factor leaves $x$ unchanged. Thus many absolute population states represent the same composition. Passing from $(N_A,N_B)$ to $x$ removes the total-size variable while retaining the ratio of the two populations.

\section{Deriving the equation for a proportion}
Differentiate $x=N_A/N$ using the quotient rule. Since $N'=f_A N_A+f_B N_B$,
\begin{align*}
x'&=\frac{N_A'N-N_A N'}{N^2}\\
&=\frac{f_A N_A(N_A+N_B)-N_A(f_A N_A+f_B N_B)}{N^2}\\
&=\frac{N_A N_B}{N^2}(f_A-f_B).
\end{align*}
Substituting the proportions gives the \emph{two-strategy replicator equation}\index{replicator equation!two strategies}:
\begin{equation}\label{eq:binaryrep}
\boxed{x'=x(1-x)\bigl(f_A-f_B\bigr).}
\end{equation}
The word \emph{strategy} can refer to an inherited type, a behavior, or another category whose frequency changes according to relative success. The mathematical model specifies the interpretation of its fitness functions.

For $0<x<1$, the factor $x(1-x)$ is positive. Consequently the sign of $x'$ is the sign of the fitness difference. A type gains share when its per-capita rate exceeds its competitor's. Both populations may be growing, or both may be shrinking; their difference still determines the change in composition.

The factor $x(1-x)$ also affects speed. It is largest at $x=1/2$ and becomes small near either endpoint. For a fixed fitness advantage, a proportion changes slowly when one type is rare. The same advantage acts on a smaller population fraction.

\begin{example}[Relative growth in a declining population]
At a particular time, let $x=0.4$, $f_A=-0.1$, and $f_B=-0.4$, with time measured in days. Find $x'$ and the total population's per-capita growth rate.
\begin{solution}
The relative advantage of A is $0.3$, so $x'=0.4(0.6)(0.3)=0.072$ per day. Meanwhile
\[
\frac{N'}{N}=xf_A+(1-x)f_B=-0.04-0.24=-0.28.
\]
The total population is declining, and A's share is increasing. Its absolute population is declining more slowly than B's.
\end{solution}
\end{example}

For a composition-dependent model, write
\[
\Delta(x)=f_A(x)-f_B(x),\qquad F(x)=x(1-x)\Delta(x).
\]
Then the model is the autonomous scalar equation $x'=F(x)$. Its \emph{equilibria}\index{equilibrium} are the states $e$ satisfying $F(e)=0$. Starting exactly at an equilibrium gives the constant solution $x(t)=e$.

\section{Phase lines and solution curves}\index{phase line}
A phase line records the sign of $F$ on intervals between equilibria. Draw an arrow to the right where $F(x)>0$ and an arrow to the left where $F(x)<0$. These arrows describe the motion of all solutions, with different initial conditions, in a single picture.

\begin{definition}[Stability on the population interval]\index{stability}
An equilibrium $e\in[0,1]$ is \emph{stable} if every sufficiently small initial displacement produces a solution that remains close to $e$ for all future time. It is \emph{asymptotically stable} if it is stable and all sufficiently nearby solutions converge to $e$. Stability is understood relative to the state space $[0,1]$.
\end{definition}

An interior equilibrium with arrows pointing toward it from both sides is asymptotically stable. Arrows pointing away from it indicate instability. At an endpoint, only the side lying inside $[0,1]$ is relevant.

Consider the frequency-dependent fitnesses
\[
f_A(x)=3-2x,\qquad f_B(x)=2x.
\]
The resulting equation is
\begin{equation}\label{eq:thread}
x'=x(1-x)(3-4x).
\end{equation}
Its equilibria are $0$, $3/4$, and $1$. For $0<x<3/4$, every factor is positive. For $3/4<x<1$, the last factor is negative. Therefore every interior solution moves toward $3/4$.

\begin{figure}[htbp]
\centering
\begin{tikzpicture}[x=5.5cm,y=1cm,>=Stealth]
\draw[->] (-.03,0)--(1.06,0) node[right] {$x$};
\draw[teal,thick,->] (.15,0)--(.42,0);
\draw[teal,thick,->] (.56,0)--(.68,0);
\draw[teal,thick,->] (.96,0)--(.82,0);
\foreach \x/\lab in {0/0,.75/{3/4},1/1}{\draw[fill=white] (\x,0) circle(2pt);\node[below=7pt] at (\x,0) {$\lab$};}
\fill[teal] (.75,0) circle(2.4pt);
\node[above=8pt] at (.75,0) {attracting};
\end{tikzpicture}
\caption{Phase line for $x'=x(1-x)(3-4x)$. The filled equilibrium attracts all interior initial states.}
\end{figure}

There are three useful graphs associated with this equation. The phase line shows directions of motion. The graph of $F(x)$ shows instantaneous velocity as a function of the state. The graph of $x(t)$ shows the state of one solution as time passes. In the last graph, the slope at time $t$ equals the height $F(x(t))$ in the velocity graph.

\begin{figure}[htbp]
\centering
\begin{tikzpicture}
\begin{axis}[width=.88\textwidth,height=5cm,xmin=0,xmax=1,ymin=-.15,ymax=.43,axis lines=middle,xlabel={$x$},ylabel={$F(x)$},xtick={0,.5,.75,1},xticklabels={0,$1/2$,$3/4$,1},grid=major]
\addplot[teal,domain=0:1,samples=101]{x*(1-x)*(3-4*x)};
\addplot[only marks,mark=*,mark size=1.8pt,ink] coordinates {(0,0)(.75,0)(1,0)};
\end{axis}
\end{tikzpicture}
\caption{Velocity versus state for the coexistence model. Positive heights correspond to increasing solution curves; negative heights correspond to decreasing solution curves.}
\end{figure}

An interior solution beginning below $3/4$ is increasing and bounded above by $3/4$. It has a limit. That limit must be an equilibrium: if $F$ were positive at the limit, continuity would keep the velocity bounded away from zero near it, preventing convergence. The only possible limit is therefore $3/4$. The same argument applies to a solution beginning above $3/4$.

The phase line and the velocity graph carry complementary information. The arrows show where the state moves and which equilibria attract. The height of the velocity graph gives the slope of the solution at the current state. Following one trajectory means reading these two pictures together: as $x(t)$ changes, the corresponding point $(x,F(x))$ moves along the velocity graph.

\begin{example}[Reading a velocity]
Evaluate the coexistence model at $x=0.4$ and $x=0.9$.
\begin{solution}
At $x=0.4$, $F(x)=0.4(0.6)(1.4)=0.336$. The solution is increasing with instantaneous slope $0.336$. At $x=0.9$, $F(x)=0.9(0.1)(-0.6)=-0.054$. That solution is decreasing. Both motions take the state toward $3/4$.
\end{solution}
\end{example}

\section{Constant advantage and the logistic solution}\index{logistic equation}
If $f_A-f_B=k$ is constant, the equation becomes
\[
x'=kx(1-x).
\]
For an interior solution, separate variables and use partial fractions:
\[
\int\frac{dx}{x(1-x)}=\int\left(\frac1x+\frac1{1-x}\right)dx=kt+C.
\]
Thus $\ln(x/(1-x))=kt+C$. Imposing $x(0)=x_0\in(0,1)$ and solving for $x$ yields
\begin{equation}\label{eq:logisticsolution}
\boxed{x(t)=\frac{1}{1+\left(\frac{1-x_0}{x_0}\right)e^{-kt}}.}
\end{equation}
For $k>0$, every interior solution approaches $1$; for $k<0$, it approaches $0$. When $k=0$, every state is stationary. The endpoint initial conditions are constant solutions and can be treated separately from the separation calculation.

\begin{example}[Time to an equal mixture]
Suppose $x'=1.2x(1-x)$ and $x(0)=1/4$. When does the population first reach $x=1/2$?
\begin{solution}
The solution is $x(t)=1/(1+3e^{-1.2t})$. Setting it equal to $1/2$ gives $3e^{-1.2t}=1$, so
\[
t=\frac{\ln3}{1.2}\approx0.916.
\]
The time unit is the one used in the fitness coefficient $1.2$.
\end{solution}
\end{example}

\begin{figure}[htbp]
\centering
\begin{tikzpicture}
\begin{axis}[width=.88\textwidth,height=5cm,xlabel={time $t$},ylabel={$x(t)$},xmin=0,xmax=5,ymin=0,ymax=1,grid=major,legend pos=south east]
\addplot[teal,domain=0:5,samples=90]{1/(1+3*exp(-1.2*x))};\addlegendentry{$x_0=1/4$}
\addplot[copper,domain=0:5,samples=90]{1/(1+exp(-1.2*x)/9)};\addlegendentry{$x_0=9/10$}
\end{axis}
\end{tikzpicture}
\caption{Two solutions with the same positive constant advantage. The initial condition determines the starting composition; the vector field determines their subsequent slopes.}
\end{figure}

The logistic form also occurs in models of population size. In the present derivation, the limiting value $1$ means complete relative dominance. The total population still depends on the two absolute growth equations.

This explicit solution supplies a time scale that the phase line alone does not provide. Two populations with the same constant fitness advantage follow the same arrows, yet their positions at a fixed time depend on the starting proportions. In the following section, linearization recovers a useful approximation to such time scales even when the full equation cannot be solved explicitly.

\section{Linearization and the rate of approach}\index{relaxation rate}
Near an equilibrium $e$, write $x=e+u$. Taylor's theorem gives
\[
u'=F(e+u)=F'(e)u+O(u^2).
\]
The linearized equation is $u'=F'(e)u$, whose solution is $u(t)=u(0)e^{F'(e)t}$. A negative derivative predicts exponential decay of small perturbations; a positive derivative predicts growth.

\begin{proposition}[Scalar linearization test]\index{linearization}
For a continuously differentiable scalar vector field, an equilibrium with $F'(e)<0$ is locally asymptotically stable, and an interior equilibrium with $F'(e)>0$ is unstable. If $F'(e)=0$, the linearization test gives no conclusion.
\end{proposition}
\begin{proof}
For $x$ close to $e$, the quotient $F(x)/(x-e)$ has the sign of $F'(e)$ when that derivative is nonzero. A negative derivative gives arrows toward $e$ on both sides; a positive derivative gives arrows away. The same reasoning applies on the available side at a boundary equilibrium.
\end{proof}

\begin{example}[A local relaxation time]
Find the linearized dynamics near the interior equilibrium of $x'=x(1-x)(3-4x)$.
\begin{solution}
At $e=3/4$, the factor $3-4e$ vanishes. Differentiating the product leaves
\[
F'(e)=e(1-e)(-4)=-\frac34.
\]
Hence $u'\approx-3u/4$. Small displacements decay approximately as $e^{-3t/4}$. Their magnitude is reduced by a factor $e$ in approximately $4/3$ time units. This is a local estimate; farther from equilibrium the nonlinear terms affect the rate.
\end{solution}
\end{example}

\begin{example}[An unstable threshold]
Analyze $x'=x(1-x)(4x-2)$.
\begin{solution}
The equilibria are $0$, $1/2$, and $1$. The velocity is negative below $1/2$ and positive above it. Thus $0$ and $1$ attract from the interior, and $1/2$ is an unstable threshold. Its derivative is $F'(1/2)=(1/2)(1/2)4=1$, agreeing with the phase line. The initial condition determines which endpoint is approached.
\end{solution}
\end{example}

\section{Recovering population sizes from a composition path}
The proportion equation was derived from an absolute population model. When the two fitnesses are specified, the total population can be recovered after the composition has been found. Its equation is
\[
N'=\bar f(x(t))N,\qquad
\bar f(x)=xf_A(x)+(1-x)f_B(x).
\]
For an initial total $N_0>0$, the solution is
\[
N(t)=N_0\exp\left(\int_0^t\bar f(x(\tau))\dd\tau\right).
\]
Then $N_A=xN$ and $N_B=(1-x)N$. To check the reconstruction, differentiate the first product:
\begin{align*}
N_A'&=N\bigl[x'+x\bar f\bigr]\\
&=Nx\bigl[(1-x)(f_A-f_B)+xf_A+(1-x)f_B\bigr]\\
&=Nx f_A=N_Af_A.
\end{align*}
The B equation follows in the same way. Thus the composition equation and the total-size equation together recover the original system.

\begin{example}[Following counts and proportions together]
Initially there are 20 A and 80 B individuals. Suppose $f_A=1$ and $f_B=0$, with time measured in days. Find the counts, composition, and total when the two types first have equal proportions.
\begin{solution}
The absolute equations give $N_A(t)=20e^t$ and $N_B(t)=80$. Consequently
\[
x(t)=\frac{20e^t}{20e^t+80}=\frac1{1+4e^{-t}}.
\]
Equal proportions occur at $e^t=4$, or $t=\ln4$. At that time both counts are 80 and the total is 160. The proportion equation gives the same time without tracking the total, while the absolute equations also give the number of individuals.
\end{solution}
\end{example}

A composition path determines the fitness difference along that path. It leaves a common contribution to both fitnesses unspecified. Such a contribution changes total growth, while canceling from the composition equation. This freedom is useful when deciding which observations a model is intended to explain.

For example, the fitness pairs $(3-4x,0)$ and $(3-2x,2x)$ both produce $\Delta(x)=3-4x$. At the equilibrium composition $x=3/4$, the first pair has both fitnesses zero, so a population already at that composition has constant total size. The second pair has both fitnesses $3/2$, so a population at the same composition grows as $N(t)=N_0e^{3t/2}$. Both retain the same $3:1$ ratio of A to B.

The term equilibrium therefore refers to the variables chosen for a state description. A composition equilibrium means stationary relative abundance. In a model retaining absolute counts as state variables, the same population may continue moving along a ray of fixed composition. Choosing the state variables establishes which features of the observations count as motion.

\section{Invariance and the probability interpretation}
The factor $x(1-x)$ makes both endpoints stationary. Suppose $\Delta$ is continuously differentiable on a neighborhood of $[0,1]$. The vector field is then locally Lipschitz, so the initial-value problem has a unique solution. An interior solution cannot cross an endpoint equilibrium in finite time: crossing would give two solutions through the same state, one constant and one moving.

This proves that $[0,1]$ is \emph{invariant}\index{invariant state space}. The bounded state space and smooth vector field also give existence for all future time. An interior trajectory can tend to an endpoint as $t\to\infty$ while remaining strictly interior at every finite time.

If one individual is sampled uniformly from the population, the probabilities of its type are $(x,1-x)$. Consequently a solution of the replicator equation is a curve of two-outcome probability distributions. The next chapters develop the same description for interaction models and for populations with three or more types.

\FloatBarrier
\section{Exercises}
\begin{exercises}
\item A population changes from 36 A and 54 B individuals to 72 A and 168 B individuals. Compute both totals and both A proportions. Interpret the two changes.
\item Derive equation \eqref{eq:binaryrep} directly by differentiating $x=N_A/(N_A+N_B)$, without introducing a separate total-size symbol.
\item At $x=0.3$, let $f_A=0.8$ and $f_B=0.2$. Compute $x'$ and $N'/N$. Explain the units of each quantity.
\item Draw a phase line and classify the equilibria of $x'=2x(1-x)$ and $x'=-3x(1-x)$ on $[0,1]$.
\item For $x'=x(1-x)(2-5x)$, find the equilibria, construct a sign chart, and calculate the derivative at the interior equilibrium.
\item For $x'=x(1-x)(5x-2)$, determine the long-term limit from initial states $0.2$, $0.4$, and $0.8$.
\item Solve $x'=2x(1-x)$ with $x(0)=1/5$. Find the first time at which $x=4/5$.
\item For positive constant advantage $k$, show that the time to move from $a$ to $b$, where $0<a<b<1$, is $k^{-1}\ln[b(1-a)/(a(1-b))]$.
\item For $F(x)=x(1-x)\Delta(x)$, prove that an interior root $e$ of $\Delta$ satisfies $F'(e)=e(1-e)\Delta'(e)$.
\item Analyze the phase line of $x'=x(1-x)(x-1/2)^2$. Describe attraction or repulsion on each side of $1/2$ and decide whether that equilibrium is stable.
\item A proposed model is $x'=0.1-x$. Decide whether $[0,1]$ is forward invariant. Compare the boundary behavior with that of a smooth replicator field.
\item Show that a smooth autonomous scalar equation on an interval has no nonconstant periodic solutions. Explain the role of uniqueness.
\item If both fitnesses are increased by the same function $c(x)$, what happens to the proportion equation? What happens to $N'/N$?
\item \emph{Exploration.} For the coexistence model \eqref{eq:thread}, compare solutions from $x_0=0.2$ and $x_0=0.9$. At a chosen time, identify the state on the phase line, its height in the velocity graph, and its slope in the time graph.
\item For $x'=x(1-x)(a-bx)$ with $b>0$ and $0<a<b$, find the interior equilibrium and its local exponential decay rate.
\item Suppose $f_A=f_B=c$ is a constant. Solve for both population sizes and the total size. Explain why every initial composition is stationary.
\item Prove that a nonconstant solution approaching a hyperbolic stable equilibrium has eventually decreasing distance from it. Compare this conclusion with the definition of stability.
\item \emph{Exploration.} Use a forward Euler step $x_{m+1}=x_m+hF(x_m)$ for $F(x)=4x(1-x)$, starting at $x_0=0.8$. Compare $h=0.1$ and $h=1$. Explain why a numerical approximation can leave an invariant interval.
\end{exercises}

\chapter{Payoff Matrices and Stability}
\label{ch:payoffs}

The fitness functions in a replicator equation can be obtained from a model of interactions. A payoff matrix records the effect of each possible encounter. Taking a weighted average over opponents converts the matrix into fitness functions and, consequently, into a differential equation.

The state now influences the rates governing its own motion. When the proportion of A changes, the typical opponent encountered by either type changes as well. This feedback can create an interior equilibrium, favor one type throughout the interval, or divide the interval into basins of attraction. The entries of a two-by-two matrix are enough to produce all four possibilities.

\section{Expected fitness from pairwise interactions}\index{payoff matrix}\index{expected fitness}
Consider two strategies, A and B, with frequencies $x$ and $1-x$. A payoff matrix has the form
\begin{equation}\label{eq:matrix2}
M=\begin{pmatrix}a&b\\c&d\end{pmatrix}.
\end{equation}
The entry in row $i$, column $j$ is the payoff received by strategy $i$ when it encounters strategy $j$. Thus $b$ is A's payoff against B, while $c$ is B's payoff against A. The two players in the same encounter can receive different payoffs.

In a well-mixed population, an opponent is A with probability $x$ and B with probability $1-x$. The expected fitnesses are therefore
\begin{equation}\label{eq:expectedfitness}
f_A(x)=ax+b(1-x),\qquad f_B(x)=cx+d(1-x).
\end{equation}
This model assumes random encounters weighted by population frequency and takes expected payoff as the per-capita rate that drives selection. If a particular application uses a different payoff-to-growth rule, that rule must be included before deriving its differential equation.

\begin{example}[Constructing the fitness functions]
Let $M=\left(\begin{smallmatrix}2&4\\3&1\end{smallmatrix}\right)$. Find the fitnesses at a general state and at $x=1/4$.
\begin{solution}
The weighted averages are
\[
f_A=2x+4(1-x)=4-2x,\qquad f_B=3x+(1-x)=1+2x.
\]
At $x=1/4$, they are $7/2$ and $3/2$. Thus A's proportion increases. At $x=3/4$, both fitnesses equal $5/2$, producing an interior equilibrium.
\end{solution}
\end{example}

\begin{figure}[htbp]
\centering
\begin{tikzpicture}
\begin{axis}[width=.88\textwidth,height=5cm,xmin=0,xmax=1,ymin=0,ymax=4.3,xlabel={proportion $x$ of A},ylabel={expected fitness},grid=major,legend pos=north east,xtick={0,.25,.5,.75,1},xticklabels={0,$1/4$,$1/2$,$3/4$,1}]
\addplot[teal,domain=0:1]{4-2*x};\addlegendentry{$f_A=4-2x$}
\addplot[copper,domain=0:1]{1+2*x};\addlegendentry{$f_B=1+2x$}
\addplot[only marks,ink,mark=*] coordinates {(.75,2.5)};
\end{axis}
\end{tikzpicture}
\caption{Expected fitnesses obtained from a payoff matrix. At their intersection the two present strategies have equal fitness, so their proportions are stationary.}
\end{figure}

\section{The advantage line}\index{fitness difference}
Subtracting the fitnesses in \eqref{eq:expectedfitness} gives
\begin{align}
\Delta(x)&=(a-c)x+(b-d)(1-x)\notag\\
&=(b-d)+(a-b-c+d)x.\label{eq:advantageline}
\end{align}
The replicator equation is
\begin{equation}\label{eq:matrixrep}
x'=x(1-x)\bigl[(b-d)+(a-b-c+d)x\bigr].
\end{equation}
The right-hand side is cubic, but its sign in $(0,1)$ is controlled by the affine function $\Delta$. This reduces the qualitative analysis to the position and slope of a line.

Write
\[
\alpha=b-d,\qquad\beta=a-b-c+d.
\]
If $\beta\ne0$, the only possible interior equilibrium is
\begin{equation}\label{eq:interioreq}
x^*=-\frac{\alpha}{\beta}=\frac{d-b}{a-b-c+d}.
\end{equation}
It represents a population state only when $0<x^*<1$. If $\beta=0$, the fitness difference is constant. A nonzero constant gives dominance; if $\alpha=\beta=0$, every composition is stationary.

The endpoint values have direct meanings:
\[
\Delta(0)=b-d,\qquad \Delta(1)=a-c.
\]
A rare A individual encounters almost exclusively B individuals. It increases in frequency when its payoff $b$ against B exceeds the resident payoff $d$. Similarly, a rare B increases near an A population when $c>a$.

\begin{definition}[Invasion from rarity]\index{invasion}
A strategy \emph{invades} a pure resident population if its proportion increases when introduced at a sufficiently small positive frequency. A strictly positive rare-strategy advantage gives invasion in the present model.
\end{definition}

The strictness matters at a parameter value where the leading rare-strategy advantage vanishes. Higher-order terms then determine the direction of motion.

\section{Four qualitative phase portraits}\index{dominance}
The advantage function is a straight line, so knowing its sign at the two endpoints determines its sign pattern throughout the interval. This reduces the classification of a differential equation to four arrangements of two signs. The four pictures below describe every nondegenerate two-strategy interaction model.
Let $\Delta_0=b-d$ and $\Delta_1=a-c$. If both are nonzero, their signs classify the dynamics completely:
\begin{center}
\begin{tabular}{@{}ccL{.31\textwidth}L{.28\textwidth}@{}}
\toprule
$\Delta_0$ & $\Delta_1$ & Behavior & Interior limits\\\midrule
$+$ & $+$ & A dominance & $x(t)\to1$\\
$-$ & $-$ & B dominance & $x(t)\to0$\\
$+$ & $-$ & Stable coexistence & $x(t)\to x^*$\\
$-$ & $+$ & Coordination & $0$ or $1$, according to the initial side of $x^*$\\
\bottomrule
\end{tabular}
\end{center}
The coexistence case allows each strategy to invade when rare. Its advantage decreases as it becomes more common, and the two fitnesses agree at an attracting interior state. The coordination case gives each strategy an advantage when common. Its interior equilibrium separates two basins of attraction.

\begin{figure}[htbp]
\centering
\begin{tikzpicture}[x=4.5cm,y=1cm,>=Stealth]
\foreach \yy/\title in {0/{A dominance},-1.15/{B dominance},-2.3/{Coexistence},-3.45/{Coordination}}{
\node[anchor=east,font=\small] at (-.08,\yy) {\title};\draw (0,\yy)--(1,\yy);\node[below=3pt,font=\small] at (0,\yy){$0$};\node[below=3pt,font=\small] at (1,\yy){$1$};}
\foreach \xx in {.15,.55}{\draw[teal,thick,->] (\xx,0)--(\xx+.22,0);\draw[teal,thick,<-] (\xx,-1.15)--(\xx+.22,-1.15);}
\draw[teal,thick,->] (.12,-2.3)--(.37,-2.3);\draw[teal,thick,<-] (.63,-2.3)--(.88,-2.3);
\draw[teal,thick,<-] (.12,-3.45)--(.37,-3.45);\draw[teal,thick,->] (.63,-3.45)--(.88,-3.45);
\foreach \yy in {0,-1.15,-2.3,-3.45}{\foreach \xx in {0,1}{\draw[fill=white] (\xx,\yy) circle(2pt);}}
\fill[teal] (1,0) circle(2pt);\fill[teal] (0,-1.15) circle(2pt);\fill[teal] (.5,-2.3) circle(2.4pt);
\draw[fill=white] (.5,-3.45) circle(2.4pt);\fill[teal] (0,-3.45) circle(2pt);\fill[teal] (1,-3.45) circle(2pt);
\end{tikzpicture}
\caption{The four nondegenerate phase portraits. Interior equilibria are drawn at the midpoint schematically; their actual positions depend on the payoffs. Filled points indicate attracting equilibria.}
\end{figure}

\begin{example}[Dominance]
For $M=\left(\begin{smallmatrix}4&3\\2&1\end{smallmatrix}\right)$, classify the dynamics.
\begin{solution}
Here $f_A=3+x$ and $f_B=1+x$, so $\Delta=2$ at every composition. The equation is $x'=2x(1-x)$. A dominates from every interior initial state, with the logistic solution derived in Chapter 1.
\end{solution}
\end{example}

\begin{example}[Stable coexistence]
Classify the model with $M=\left(\begin{smallmatrix}2&4\\3&1\end{smallmatrix}\right)$.
\begin{solution}
The endpoint advantages are $3$ and $-1$. Thus both types invade from rarity. The fitness difference is $3-4x$, giving the equilibrium $x^*=3/4$. The limit from any interior state is a population consisting of $75\%$ A and $25\%$ B.
\end{solution}
\end{example}

\begin{example}[Coordination and basins of attraction]\index{basin of attraction}\index{coordination}
For $M=\left(\begin{smallmatrix}4&1\\2&3\end{smallmatrix}\right)$, find the threshold and the limits from $x_0=0.35$ and $x_0=0.65$.
\begin{solution}
The fitnesses are $f_A=1+3x$ and $f_B=3-x$, so $\Delta=-2+4x$. The threshold is $x^*=1/2$. Below it, A loses share; above it, A gains share. Hence the limits are $0$ and $1$, respectively. The basin of $0$ within the state interval is $[0,1/2)$, and the basin of $1$ is $(1/2,1]$. The threshold itself is stationary.
\end{solution}
\end{example}

\section{Linearization and degenerate cases}
For $F(x)=x(1-x)(\alpha+\beta x)$, differentiation gives
\[
F'(x)=(1-2x)(\alpha+\beta x)+\beta x(1-x).
\]
At an interior equilibrium, the first term vanishes. Therefore
\begin{equation}\label{eq:matrixlinear}
F'(x^*)=\beta x^*(1-x^*).
\end{equation}
Because $x^*(1-x^*)>0$, a negative advantage-line slope gives attraction, and a positive slope gives repulsion. At the endpoints,
\[
F'(0)=b-d,\qquad F'(1)=c-a.
\]
These formulas agree with the invasion interpretation. A negative derivative at an endpoint gives attraction from the available interior side.

\begin{example}[A zero derivative at an attracting endpoint]
Let $M=\left(\begin{smallmatrix}3&4\\3&1\end{smallmatrix}\right)$. Determine the behavior near $x=1$.
\begin{solution}
The advantage is $3(1-x)$, so
\[
x'=3x(1-x)^2.
\]
Every interior velocity is positive, and every interior solution approaches 1. Yet $F'(1)=0$. Setting $u=1-x$ gives $u'=-3(1-u)u^2$. Near $u=0$ this is approximately $u'=-3u^2$, with decay on the scale $u(t)\approx1/(3t+C)$. The phase line establishes attraction even though the exponential linearization is inconclusive.
\end{solution}
\end{example}

If the two rows of $M$ agree, the expected fitnesses agree at every state. Conversely, $\Delta\equiv0$ requires $b=d$ and $a=c$, so the rows agree. This gives a continuum of equilibria. Every state is stable in the sense of remaining close to its initial state, while no individual state attracts its distinct neighbors.

\section{Parameters and bifurcation}
An interaction model can change character when a payoff changes. We can see where this happens by tracking the advantage at an endpoint. An interior equilibrium may move toward that endpoint and disappear into it, changing the long-term behavior of nearby trajectories.
Hold three payoffs fixed and vary the fourth:
\[
M(a)=\begin{pmatrix}a&4\\3&1\end{pmatrix},\qquad \Delta(x)=3+(a-6)x.
\]
For $a<3$, there is an attracting interior equilibrium
\[
x^*(a)=\frac{3}{6-a}.
\]
It approaches the boundary value $1$ as $a\to3$ from below. At $a=3$ the model is the degenerate attracting example above. For $a>3$, both endpoint advantages are positive, and A dominates.

\begin{definition}[Bifurcation]\index{bifurcation}
A \emph{bifurcation} occurs at a parameter value where the qualitative equilibrium or trajectory structure changes as the parameter varies.
\end{definition}

In this family the physically admissible coexistence state merges with a boundary equilibrium. The parameter determines both the final population composition and the rate at which that composition is approached. In particular, the linearized attraction rate at $x^*(a)$ tends to zero as the equilibrium approaches the boundary.

\begin{figure}[htbp]
\centering
\begin{tikzpicture}
\begin{axis}[width=.88\textwidth,height=5cm,xmin=0,xmax=4.5,ymin=0,ymax=1.08,xlabel={payoff parameter $a$},ylabel={equilibrium $x$},grid=major,legend pos=south east]
\addplot[teal,domain=0:3,samples=70]{3/(6-x)};\addlegendentry{attracting}
\addplot[copper,dashed,domain=0:3]{1};\addlegendentry{repelling}
\addplot[teal,domain=3:4.5]{1};
\addplot[copper,dashed,domain=0:4.5]{0};
\addplot[only marks,ink] coordinates {(3,1)};
\end{axis}
\end{tikzpicture}
\caption{Equilibrium branches within the population interval. Stable coexistence reaches the boundary at $a=3$; for larger $a$, the boundary state attracts. The critical point is nonhyperbolic.}
\end{figure}

\section{Transformations of payoffs}
Only fitness differences enter the proportion equation. Adding a constant $k$ to every entry of $M$ increases each expected fitness by $k$ and leaves the vector field unchanged. A more general transformation also has this property:
\begin{equation}\label{eq:columnshift}
\widetilde M=M+\ones r^T,
\end{equation}
where $\ones=(1,1)^T$ and $r=(r_A,r_B)^T$. This adds $r_A$ to both entries of the first column and $r_B$ to both entries of the second. Each strategy then gains the same state-dependent payoff $r_Ax+r_B(1-x)$, which cancels in the difference.

Multiplying all payoffs by a positive constant $\lambda$ gives the vector field $\lambda F$. If $x(t)$ solves the original equation, $x(\lambda t)$ solves the transformed equation with the same initial state. The equilibria, orbit directions, and limiting states agree, while the time scale changes.

\begin{example}[Equivalent interaction models]
Subtract the second row of $M=\left(\begin{smallmatrix}2&4\\3&1\end{smallmatrix}\right)$ from both rows, column by column.
\begin{solution}
The transformed matrix is $\left(\begin{smallmatrix}-1&3\\0&0\end{smallmatrix}\right)$. Its B fitness is zero, and its A fitness is $-x+3(1-x)=3-4x$. The replicator equation remains $x'=x(1-x)(3-4x)$. The transformation provides a simpler representation of the same composition dynamics.
\end{solution}
\end{example}

\section{Designing and identifying an interaction model}
The passage from a matrix to a phase portrait can also be used in reverse. Suppose the intended advantage line is $\Delta(x)=\alpha+\beta x$. One matrix realizing it is
\[
M=\begin{pmatrix}\alpha+\beta&\alpha\\0&0\end{pmatrix}.
\]
Against this matrix, the expected B payoff is zero and the expected A payoff is $(\alpha+\beta)x+\alpha(1-x)=\alpha+\beta x$. Column shifts then generate every matrix with the same row differences. This supplies a direct way to build a test model with specified scalar dynamics.

\begin{example}[Choosing an equilibrium and a relaxation rate]
Construct a two-strategy model with attracting composition $q=0.4$ and local linearized equation $u'=-0.6u$.
\begin{solution}
Choose $\Delta(x)=b(q-x)$ with $b>0$. Its equilibrium derivative is $F'(q)=-bq(1-q)$. The desired value requires
\[
b=\frac{0.6}{0.4(0.6)}=2.5.
\]
Thus $\Delta(x)=1-2.5x$, realized by
\[
M=\begin{pmatrix}-1.5&1\\0&0\end{pmatrix}.
\]
Adding 2 to all entries gives the nonnegative matrix $\left(\begin{smallmatrix}0.5&3\\2&2\end{smallmatrix}\right)$ with identical composition dynamics. The equilibrium and decay rate specify the relative selection pattern, while the common payoff level remains free.
\end{solution}
\end{example}

Observed velocities can be used to infer the advantage line. At an interior state,
\[
\Delta(x)=\frac{x'}{x(1-x)}.
\]
Two distinct state-velocity observations determine the intercept and slope of a proposed affine advantage. A third observation can test that fitted relation. This procedure uses instantaneous velocity estimates; a finite difference over a long time interval generally averages several velocities and requires additional approximation.

\begin{example}[Inferring the advantage from two velocities]
Suppose measured instantaneous proportion velocities are $x'=0.16$ at $x=0.2$ and $x'=-0.24$ at $x=0.6$. Fit a two-strategy matrix model.
\begin{solution}
Dividing by $x(1-x)$ gives advantages 1 and $-1$. The line through $(0.2,1)$ and $(0.6,-1)$ has slope $-5$ and intercept 2. Hence the fitted model is
\[
x'=x(1-x)(2-5x).
\]
Its attracting state is $2/5$, and one payoff matrix is $\left(\begin{smallmatrix}-3&2\\0&0\end{smallmatrix}\right)$. The model predicts velocity $-0.32$ at $x=0.8$, providing an additional value against which to compare observations.
\end{solution}
\end{example}

The division also explains a feature of data interpretation near an endpoint. An error $\varepsilon$ in a velocity estimate produces an error $\varepsilon/[x(1-x)]$ in the inferred advantage. When one type is rare, the same absolute velocity error produces a larger error in the fitted advantage. The phase portrait can still be estimated from direction information, while precise fitness estimation requires attention to the scale of the observations.

\FloatBarrier
\section{Exercises}
\begin{exercises}
\item For $M=\left(\begin{smallmatrix}1&5\\3&2\end{smallmatrix}\right)$, calculate both expected fitnesses, the advantage line, and all equilibria in $[0,1]$.
\item Explain the meaning of all four entries of a general payoff matrix. Describe the random-mixing assumption used in the weighted averages.
\item Classify $M=\left(\begin{smallmatrix}5&2\\3&1\end{smallmatrix}\right)$ using endpoint advantages. Confirm by a phase line.
\item Construct one matrix for B dominance and one for coordination with threshold $x^*=1/3$.
\item For $M=\left(\begin{smallmatrix}1&4\\3&2\end{smallmatrix}\right)$, find the coexistence equilibrium and its local linearized decay rate.
\item Prove the four-case classification table from the fact that $\Delta$ is affine.
\item Find all two-strategy payoff matrices for which every composition is an equilibrium.
\item Show that $M$ and $M+\ones r^T$ generate the same replicator equation, where $r$ is any two-component vector.
\item Show that positive payoff scaling preserves stability. Explain what negative scaling does to the phase arrows and to hyperbolic equilibrium stability.
\item For $M(a)$ in the parameter family above, compute $x^*$ at $a=1,2,2.8$. Find the derivative $F'(x^*)$ at each value.
\item Analyze $x'=x^2(1-x)$ and classify its two endpoint equilibria. Which linearization test is inconclusive?
\item Suppose $b=d$ and $a>c$. Determine the entire phase portrait, including the behavior of a rare A type.
\item For $M=\left(\begin{smallmatrix}4&1\\2&3\end{smallmatrix}\right)$, calculate the fitnesses and the velocity at $x=0.35$ and $x=0.65$. Use these values to interpret coordination.
\item Show that if A can invade B and B can invade A, an interior attracting equilibrium exists for a two-strategy matrix model with strict invasion inequalities.
\item \emph{Exploration.} In the payoff-matrix demonstration, vary $a$ through 3 while $b=4,c=3,d=1$. Compare the fitness intersection, phase line, and time series near the critical parameter.
\item Can two different matrices have identical advantage lines? Characterize all such pairs by their row differences.
\item In the coexistence matrix $M=\left(\begin{smallmatrix}2&4\\3&1\end{smallmatrix}\right)$, compute the equilibrium mean fitness. What is it after adding $10$ to every matrix entry?
\item \emph{Exploration.} Compare the coexistence matrix with twice that matrix. Use the same initial state and determine which times should be compared to obtain matching compositions.
\end{exercises}

\chapter{Three Strategies and the Probability Simplex}
\label{ch:simplex}

Three population proportions require two independent coordinates. Their state space is a triangle, and a trajectory can turn as it moves through that triangle. Cyclic interactions supply a concrete planar system in which trajectories can approach equilibrium, circulate periodically, or move toward the boundary.

The interval of the preceding chapters recorded a single proportion; its direction of motion was determined by a sign. On the probability triangle, a velocity has two independent components. The possibility of turning makes planar geometry essential. The edges of the triangle still contain the familiar two-type models, while the interior permits genuinely three-type behavior.

\section{A triangle of normalized states}\index{simplex!dimension}
Let $p_A,p_B,p_C$ be the proportions of three types. They satisfy
\[
p_A+p_B+p_C=1,\qquad p_A,p_B,p_C\geq0.
\]
The set of such vectors is the \emph{standard 2-simplex}\index{simplex}:
\[
\Delta^2=\{(p_A,p_B,p_C)\in\R^3:p_i\geq0,\ p_A+p_B+p_C=1\}.
\]
It is the triangular region cut from the first octant by the normalization plane. The vertices $(1,0,0)$, $(0,1,0)$, and $(0,0,1)$ are pure states. An edge consists of mixtures in which one type is absent. All three types are present in the relative interior.

For a flat equilateral drawing, choose vertices $A=(0,0)$, $B=(1,0)$, and $C=(1/2,\sqrt3/2)$. The displayed position of a state is
\begin{equation}\label{eq:barycentric}
r(\p)=p_AA+p_BB+p_CC=\left(p_B+\frac12p_C,\frac{\sqrt3}{2}p_C\right).
\end{equation}
The probabilities are \emph{barycentric coordinates}\index{barycentric coordinates}: the state is their weighted average of the three vertices. A point halfway along the AB edge has state $(1/2,1/2,0)$.

\begin{figure}[htbp]
\centering
\begin{tikzpicture}[scale=4.2]
\fill[soft] (0,0)--(1,0)--(.5,.866025)--cycle;
\draw[ink] (0,0)--(1,0)--(.5,.866025)--cycle;
\foreach \t in {.2,.4,.6,.8}{\draw[black!15] ({.5*\t},{.866025*\t})--({1-.5*\t},{.866025*\t});}
\node[below left] at (0,0){$A$};\node[below right] at (1,0){$B$};\node[above] at (.5,.866025){$C$};
\fill[teal] (.4,.173205) circle(.014);
\node[above right,font=\small] at (.4,.173205){$(.5,.3,.2)$};
\draw[copper,thick,->] (.4,.173205)--(.30,.26);
\node[font=\small,anchor=east] at (.28,.28){tangent velocity};
\end{tikzpicture}
\caption{A three-type population state and a tangent velocity in an equilateral simplex drawing. The planar arrow is the image of a three-component velocity whose components sum to zero.}
\end{figure}

One can instead use Cartesian coordinates $x=p_A$, $y=p_B$, and $p_C=1-x-y$. Their domain is the right triangle $x\geq0$, $y\geq0$, $x+y\leq1$. These two drawings give different linear pictures of the same states.

\section{The general replicator system}\index{replicator equation!general}\index{tangent space}
Let $N_i'=f_iN_i$ for $i=A,B,C$, and write $N=\sum_iN_i$ and $p_i=N_i/N$. Then
\[
N'=N\sum_i p_i f_i.
\]
The quantity
\[
\bar f=\sum_i p_i f_i
\]
is the \emph{mean fitness}\index{mean fitness}. Differentiating $p_i$ gives
\begin{equation}\label{eq:generalrep}
\boxed{p_i'=p_i(f_i-\bar f).}
\end{equation}
The same derivation applies to any finite number $n$ of types. Each type changes according to its advantage over the current population average.

\begin{example}[Computing a velocity]
At $\p=(0.2,0.3,0.5)$ suppose $\f=(3,1,2)$. Find the mean fitness and the replicator velocity.
\begin{solution}
The mean is $0.2(3)+0.3(1)+0.5(2)=1.9$. Thus
\[
\p'=(0.2(1.1),0.3(-0.9),0.5(0.1))=(0.22,-0.27,0.05).
\]
Types A and C gain share. The components sum to zero, so the increases exactly balance B's decrease.
\end{solution}
\end{example}

For $n$ types, the closed state space is $\Delta^{n-1}$, of dimension $n-1$. The constraint removes one independent coordinate. A differentiable curve in its normalization plane has tangent velocity satisfying
\begin{equation}\label{eq:tangent}
\sum_i p_i'=0.
\end{equation}
For the replicator field this follows algebraically:
\[
\sum_i p_i(f_i-\bar f)=\bar f-\bar f\sum_i p_i=0.
\]

\section{Invariant faces and positive solutions}\index{invariant face}
If $p_k=0$, then the $k$th replicator component is zero. With smooth fitness functions, uniqueness therefore makes every face of the simplex invariant. A population that begins with a type absent retains that absence under the model.

An explicit identity describes positivity. Along a solution with $p_i(0)>0$,
\begin{equation}\label{eq:positivity}
p_i(t)=p_i(0)\exp\left(\int_0^t[f_i(\p(\tau))-\bar f(\p(\tau))]\dd\tau\right).
\end{equation}
If the fitness functions are smooth on a neighborhood of the compact simplex, the integrand is bounded on every finite time interval. The exponential is then positive. A type present initially remains present at finite times, although its proportion can tend to zero as time increases without bound.

\begin{example}[An edge carries a two-strategy equation]
Consider the matrix
\[
M_s=\begin{pmatrix}0&1&-s\\-s&0&1\\1&-s&0\end{pmatrix},\qquad s>0.
\]
Reduce its dynamics to the AB edge, where $p_C=0$.
\begin{solution}
Set $p_A=x$ and $p_B=1-x$. Then $f_A=1-x$ and $f_B=-sx$. The edge equation is
\[
x'=x(1-x)\bigl[1+(s-1)x\bigr].
\]
The bracket is positive on $[0,1]$ because its endpoint values are $1$ and $s$. A therefore wins along the AB edge. The absent C type stays absent, even though it would interact with A if it were introduced.
\end{solution}
\end{example}

Interior and edge trajectories can consequently have different limiting behavior. A small positive introduction of a missing type changes the state from one invariant face to a larger state space. This distinction is essential when interpreting the triangle's boundary.

\section{Cyclic competition}\index{cyclic competition}
For a payoff matrix $M$, random mixing gives
\[
\f=M\p,\qquad \bar f=\p^TM\p.
\]
The cyclic family $M_s$ above assigns payoff $+1$ to a win, $-s$ to a loss, and zero to an encounter with the same type. A outperforms B, B outperforms C, and C outperforms A. Its fitnesses are
\begin{equation}\label{eq:cyclicfitness}
f_A=p_B-sp_C,\quad f_B=p_C-sp_A,\quad f_C=p_A-sp_B.
\end{equation}
At the uniform state $\p^*=(1/3,1/3,1/3)$, all three fitnesses equal $(1-s)/3$. Hence $\p^*$ is an equilibrium for every $s>0$.

\begin{example}[A first direction of cyclic motion]
For $s=1$ and $\p=(1/2,1/3,1/6)$, compute the fitnesses and the velocity.
\begin{solution}
The fitness vector is $(1/6,-1/3,1/6)$. Its mean is zero, so
\[
\p'=\left(\frac1{12},-\frac19,\frac1{36}\right).
\]
The sum is zero. A and C increase at this instant, while B decreases. The identities of the increasing types change as the state moves, producing circulation.
\end{solution}
\end{example}

When $s=1$, the matrix is antisymmetric: $M_1^T=-M_1$. For every vector $\p$, the scalar $\p^TM_1\p$ equals its transpose, which is also $-\p^TM_1\p$. Thus the mean fitness is identically zero. Individual fitnesses can differ from zero and drive motion.

\section{A product that detects inward and outward motion}
The pictures of cyclic motion suggest three different tendencies: inward spiraling, closed circulation, and outward spiraling. To distinguish them beyond the linear approximation near the center, we seek a scalar quantity whose change has a definite sign throughout the interior. The product of the three probabilities serves this purpose.
For an interior state, define
\[
P=p_Ap_Bp_C,\qquad Q=p_Ap_B+p_Bp_C+p_Cp_A.
\]
Differentiating the logarithm of the product gives
\begin{align*}
\frac{d}{dt}\ln P
&=\sum_i\frac{p_i'}{p_i}
=\sum_i f_i-3\bar f.
\end{align*}
The cyclic fitnesses satisfy $\sum_i f_i=1-s$, while
\[
\bar f=(1-s)Q.
\]
Therefore
\begin{equation}\label{eq:productderivative}
\boxed{\frac{d}{dt}\ln P=(1-s)(1-3Q).}
\end{equation}
The second factor is nonnegative. Indeed,
\[
1-3Q=\frac12\bigl[(p_A-p_B)^2+(p_B-p_C)^2+(p_C-p_A)^2\bigr],
\]
and it vanishes precisely at the uniform state. This identity determines the direction of product change throughout the interior.

If $s<1$, the product increases at every nonstationary interior point. Its maximum is $1/27$, attained at the center. If $s>1$, it decreases. If $s=1$, it is conserved.

\begin{proposition}[Balanced cyclic orbits]\index{periodic orbit}\index{conserved quantity}
For $s=1$, every nonstationary interior trajectory lies on a periodic orbit surrounding $\p^*$.
\end{proposition}
\begin{proof}
The function $\ln P=\ln p_A+\ln p_B+\ln p_C$ is strictly concave in the simplex interior, tends to $-\infty$ at the boundary, and has a unique critical point at the center. For $0<c<1/27$, the level $P=c$ is a smooth closed curve surrounding the center. Conservation makes the vector field tangent to this curve. An interior equilibrium requires equal fitnesses, which in this balanced model forces all three probabilities to agree. Thus the field is nonzero on $P=c$. A continuous nonzero velocity on a compact smooth circle carries the state around the circle in finite time and repeats.
\end{proof}

The monotonicity also explains the other two regimes globally. For $s<1$, a trajectory remains in the compact interior set $P\geq P(0)>0$. Any limiting invariant set must have constant limiting product, hence must lie where \eqref{eq:productderivative} vanishes. That is the center. For $s>1$, a noncentral trajectory has decreasing product. If its limit were positive, the same compactness argument would force convergence to the center, whose product is larger than the initial product. This is impossible; therefore $P\to0$, and the trajectory approaches the boundary as a set.

The last conclusion allows prolonged visits near different vertices. It does not specify convergence to one pure state. The directions along the edges form a cycle, so the currently predominant type can continue changing.

\begin{figure}[htbp]\centering\begin{tikzpicture}
\begin{scope}[xshift=0.0cm]\begin{scope}[scale=4.0]
\fill[soft] (0,0)--(1,0)--(.5,.8660254)--cycle;
\draw[ink!75] (0,0)--(1,0)--(.5,.8660254)--cycle;
\node[font=\small] at (.5,1.08){$s=0.7$};
\draw[teal,line width=.85pt] plot coordinates {(0.3500000,0.1732051) (0.3462275,0.1782764) (0.3429071,0.1835900) (0.3400367,0.1891428) (0.3376130,0.1949304) (0.3356314,0.2009474) (0.3340858,0.2071872) (0.3329690,0.2136417) (0.3322725,0.2203018) (0.3319870,0.2271566) (0.3321018,0.2341941) (0.3326053,0.2414007) (0.3334853,0.2487616) (0.3347285,0.2562603) (0.3363210,0.2638792) (0.3382480,0.2715994) (0.3404946,0.2794009) (0.3430450,0.2872623) (0.3458833,0.2951617) (0.3489934,0.3030760) (0.3523588,0.3109816) (0.3559632,0.3188542) (0.3597902,0.3266694) (0.3638236,0.3344024) (0.3680475,0.3420283) (0.3724463,0.3495225) (0.3770050,0.3568605) (0.3817087,0.3640184) (0.3865436,0.3709728) (0.3914961,0.3777010) (0.3965536,0.3841812) (0.4017039,0.3903925) (0.4069357,0.3963151) (0.4122385,0.4019303) (0.4176025,0.4072208) (0.4230186,0.4121705) (0.4284783,0.4167646) (0.4339740,0.4209898) (0.4394985,0.4248344) (0.4450455,0.4282878) (0.4506089,0.4313412) (0.4561832,0.4339872) (0.4617633,0.4362201) (0.4673445,0.4380356) (0.4729222,0.4394309) (0.4784919,0.4404048) (0.4840496,0.4409576) (0.4895909,0.4410911) (0.4951115,0.4408088) (0.5006070,0.4401153) (0.5060729,0.4390170) (0.5115043,0.4375216) (0.5168961,0.4356381) (0.5222429,0.4333769) (0.5275387,0.4307497) (0.5327775,0.4277695) (0.5379524,0.4244504) (0.5430565,0.4208075) (0.5480821,0.4168571) (0.5530213,0.4126164) (0.5578657,0.4081033) (0.5626065,0.4033365) (0.5672346,0.3983353) (0.5717404,0.3931196) (0.5761142,0.3877096) (0.5803460,0.3821258) (0.5844256,0.3763889) (0.5883427,0.3705196) (0.5920871,0.3645386) (0.5956483,0.3584662) (0.5990161,0.3523227) (0.6021804,0.3461277) (0.6051314,0.3399006) (0.6078594,0.3336599) (0.6103552,0.3274239) (0.6126099,0.3212096) (0.6146152,0.3150337) (0.6163632,0.3089119) (0.6178467,0.3028589) (0.6190590,0.2968887) (0.6199943,0.2910143) (0.6206473,0.2852477) (0.6210137,0.2796001) (0.6210899,0.2740818) (0.6208731,0.2687022) (0.6203615,0.2634696) (0.6195543,0.2583917) (0.6184515,0.2534754) (0.6170541,0.2487266) (0.6153643,0.2441506) (0.6133850,0.2397521) (0.6111204,0.2355348) (0.6085757,0.2315023) (0.6057569,0.2276572) (0.6026713,0.2240019) (0.5993271,0.2205380) (0.5957336,0.2172670) (0.5919009,0.2141897) (0.5878404,0.2113070) (0.5835640,0.2086190) (0.5790847,0.2061259) (0.5744164,0.2038275) (0.5695735,0.2017235) (0.5645713,0.1998132) (0.5594258,0.1980962) (0.5541533,0.1965715) (0.5487708,0.1952384) (0.5432955,0.1940958) (0.5377450,0.1931428) (0.5321372,0.1923783) (0.5264900,0.1918011) (0.5208212,0.1914101) (0.5151488,0.1912042) (0.5094905,0.1911821) (0.5038637,0.1913426) (0.4982855,0.1916844) (0.4927728,0.1922062) (0.4873417,0.1929065) (0.4820078,0.1937839) (0.4767863,0.1948368) (0.4716915,0.1960637) (0.4667371,0.1974626) (0.4619359,0.1990318) (0.4572999,0.2007691) (0.4528404,0.2026722) (0.4485676,0.2047387) (0.4444910,0.2069659) (0.4406191,0.2093507) (0.4369594,0.2118900) (0.4335188,0.2145802) (0.4303028,0.2174173) (0.4273164,0.2203971) (0.4245635,0.2235150) (0.4220471,0.2267660) (0.4197695,0.2301445) (0.4177321,0.2336447) (0.4159353,0.2372604) (0.4143789,0.2409846) (0.4130619,0.2448102) (0.4119827,0.2487296) (0.4111387,0.2527345) (0.4105270,0.2568164) (0.4101439,0.2609664) (0.4099852,0.2651750) (0.4100460,0.2694323) (0.4103211,0.2737283) (0.4108050,0.2780525) (0.4114913,0.2823939) (0.4123738,0.2867417) (0.4134457,0.2910844) (0.4146999,0.2954106) (0.4161292,0.2997089) (0.4177261,0.3039674) (0.4194831,0.3081746) (0.4213925,0.3123188) (0.4234465,0.3163884) (0.4256372,0.3203721) (0.4279570,0.3242586) (0.4303980,0.3280370) (0.4329524,0.3316966) (0.4356125,0.3352270) (0.4383709,0.3386184) (0.4412200,0.3418613) (0.4441525,0.3449467) (0.4471611,0.3478662) (0.4502389,0.3506118) (0.4533787,0.3531763) (0.4565739,0.3555531) (0.4598179,0.3577361) (0.4631041,0.3597200) (0.4664262,0.3615002) (0.4697779,0.3630729) (0.4731532,0.3644348) (0.4765461,0.3655836) (0.4799507,0.3665175) (0.4833613,0.3672356) (0.4867720,0.3677377) (0.4901774,0.3680244) (0.4935718,0.3680969) (0.4969496,0.3679572) (0.5003054,0.3676080) (0.5036338,0.3670525) (0.5069291,0.3662949) (0.5101861,0.3653398) (0.5133992,0.3641924) (0.5165631,0.3628585) (0.5196723,0.3613445) (0.5227213,0.3596572) (0.5257048,0.3578040) (0.5286173,0.3557927) (0.5314535,0.3536314) (0.5342079,0.3513286) (0.5368753,0.3488931) (0.5394504,0.3463340) (0.5419280,0.3436605) (0.5443030,0.3408821) (0.5465704,0.3380084) (0.5487251,0.3350490) (0.5507626,0.3320137) (0.5526782,0.3289122) (0.5544675,0.3257542) (0.5561262,0.3225492) (0.5576505,0.3193068) (0.5590364,0.3160364) (0.5602807,0.3127470) (0.5613800,0.3094477) (0.5623315,0.3061472) (0.5631328,0.3028541) (0.5637815,0.2995765) (0.5642758,0.2963223) (0.5646144,0.2930992) (0.5647962,0.2899144) (0.5648206,0.2867750) (0.5646873,0.2836876) (0.5643966,0.2806584) (0.5639492,0.2776934) (0.5633463,0.2747982) (0.5625893,0.2719780) (0.5616804,0.2692378) (0.5606221,0.2665820) (0.5594174,0.2640151) (0.5580696,0.2615408) (0.5565826,0.2591627) (0.5549608,0.2568843) (0.5532089,0.2547084) (0.5513320,0.2526378) (0.5493356,0.2506749) (0.5472258,0.2488220) (0.5450087,0.2470808) (0.5426911,0.2454531) (0.5402797,0.2439403) (0.5377820,0.2425437) (0.5352053,0.2412643) (0.5325573,0.2401028) (0.5298460,0.2390599) (0.5270795,0.2381360) (0.5242660,0.2373313) (0.5214137,0.2366459) (0.5185312,0.2360798) (0.5156268,0.2356326) (0.5127090,0.2353039) (0.5097862,0.2350931) (0.5068667,0.2349996) (0.5039587,0.2350224) (0.5010704,0.2351605) (0.4982097,0.2354128) (0.4953844,0.2357778) (0.4926021,0.2362542) (0.4898701,0.2368402) (0.4871953,0.2375342) (0.4845847,0.2383341) (0.4820445,0.2392379) (0.4795810,0.2402432) (0.4772000,0.2413478) (0.4749068,0.2425489) (0.4727066,0.2438439) (0.4706039,0.2452297) (0.4686032,0.2467034) (0.4667082,0.2482615) (0.4649225,0.2499007) (0.4632492,0.2516173) (0.4616910,0.2534076) (0.4602501,0.2552674) (0.4589285,0.2571928) (0.4577277,0.2591793) (0.4566489,0.2612225) (0.4556927,0.2633177) (0.4548596,0.2654601) (0.4541496,0.2676447) (0.4535624,0.2698665) (0.4530974,0.2721203) (0.4527536,0.2744006) (0.4525297,0.2767022) (0.4524241,0.2790193) (0.4524352,0.2813464) (0.4525607,0.2836779) (0.4527983,0.2860079) (0.4531456,0.2883307) (0.4535997,0.2906406) (0.4541577,0.2929318) (0.4548164,0.2951985) (0.4555726,0.2974351) (0.4564229,0.2996358) (0.4573636,0.3017953) (0.4583910,0.3039079) (0.4595015,0.3059683) (0.4606910,0.3079715) (0.4619558,0.3099123) (0.4632916,0.3117859) (0.4646945,0.3135877) (0.4661604,0.3153133) (0.4676851,0.3169583) (0.4692645,0.3185190) (0.4708943,0.3199915) (0.4725703,0.3213724) (0.4742885,0.3226587) (0.4760445,0.3238475) (0.4778343,0.3249363) (0.4796535,0.3259227) (0.4814982,0.3268050) (0.4833642,0.3275816) (0.4852474,0.3282512) (0.4871437,0.3288129) (0.4890491,0.3292662) (0.4909596,0.3296108) (0.4928713,0.3298468) (0.4947801,0.3299746) (0.4966824,0.3299949) (0.4985741,0.3299089) (0.5004516,0.3297179) (0.5023110,0.3294234) (0.5041488,0.3290276) (0.5059612,0.3285327) (0.5077448,0.3279410) (0.5094959,0.3272555) (0.5112112,0.3264791) (0.5128873,0.3256151) (0.5145208,0.3246668) (0.5161086,0.3236380) (0.5176476,0.3225324) (0.5191345,0.3213540) (0.5205667,0.3201070) (0.5219411,0.3187956) (0.5232550,0.3174242) (0.5245059,0.3159972) (0.5256911,0.3145193) (0.5268085,0.3129949) (0.5278556,0.3114288) (0.5288305,0.3098257) (0.5297312,0.3081901) (0.5305559,0.3065269) (0.5313030,0.3048407) (0.5319711,0.3031360) (0.5325589,0.3014176) (0.5330652,0.2996899) (0.5334893,0.2979574) (0.5338304,0.2962245) (0.5340880,0.2944954) (0.5342618,0.2927744) (0.5343517,0.2910656) (0.5343579,0.2893728) (0.5342806,0.2877000) (0.5341205,0.2860508) (0.5338782,0.2844288) (0.5335547,0.2828374) (0.5331512,0.2812799) (0.5326692,0.2797595) (0.5321101,0.2782791) (0.5314759,0.2768416) (0.5307685,0.2754496) (0.5299901,0.2741056) (0.5291431,0.2728120) (0.5282301,0.2715710) (0.5272538,0.2703847) (0.5262172,0.2692550) (0.5251234,0.2681835) (0.5239755,0.2671719) (0.5227770,0.2662217) (0.5215313,0.2653341) (0.5202421,0.2645102) (0.5189130,0.2637511) (0.5175480,0.2630577) (0.5161507,0.2624305) (0.5147254,0.2618703) (0.5132758,0.2613774) (0.5118062,0.2609522) (0.5103205,0.2605948) (0.5088228,0.2603053) (0.5073173,0.2600835) (0.5058080,0.2599293) (0.5042990,0.2598423) (0.5027942,0.2598221) (0.5012976,0.2598681) (0.4998131,0.2599795) (0.4983444,0.2601556) (0.4968954,0.2603954) (0.4954696,0.2606978) (0.4940705,0.2610617) (0.4927015,0.2614859) (0.4913658,0.2619688) (0.4900666,0.2625091) (0.4888069,0.2631052) (0.4875893,0.2637552) (0.4864167,0.2644574) (0.4852914,0.2652100) (0.4842159,0.2660108) (0.4831922,0.2668579) (0.4822223,0.2677489) (0.4813080,0.2686817) (0.4804509,0.2696538) (0.4796525,0.2706629) (0.4789139,0.2717063) (0.4782363,0.2727816) (0.4776205,0.2738860) (0.4770672,0.2750168) (0.4765769,0.2761713) (0.4761501,0.2773466) (0.4757868,0.2785398) (0.4754871,0.2797480) (0.4752509,0.2809682) (0.4750777,0.2821976) (0.4749672,0.2834330) (0.4749187,0.2846715) (0.4749315,0.2859101) (0.4750047,0.2871457) (0.4751372,0.2883754) (0.4753279,0.2895961) (0.4755755,0.2908050) (0.4758786,0.2919992) (0.4762358,0.2931756) (0.4766455,0.2943316) (0.4771059,0.2954644) (0.4776153,0.2965712) (0.4781718,0.2976496) (0.4787735,0.2986969) (0.4794185,0.2997107) (0.4801046,0.3006886) (0.4808297,0.3016285) (0.4815916,0.3025283) (0.4823882,0.3033858) (0.4832172,0.3041993) (0.4840764,0.3049670) (0.4849633,0.3056873) (0.4858756,0.3063587) (0.4868110,0.3069798) (0.4877671,0.3075496) (0.4887416,0.3080669) (0.4897319,0.3085309) (0.4907358,0.3089409) (0.4917508,0.3092962) (0.4927746,0.3095965) (0.4938047,0.3098414) (0.4948389,0.3100310) (0.4958747,0.3101651) (0.4969098,0.3102441) (0.4979419,0.3102683) (0.4989688,0.3102381) (0.4999882,0.3101542) (0.5009979,0.3100174) (0.5019958,0.3098286) (0.5029796,0.3095889) (0.5039473,0.3092994) (0.5048970,0.3089615) (0.5058265,0.3085764) (0.5067340,0.3081459) (0.5076176,0.3076715) (0.5084754,0.3071549) (0.5093058,0.3065980) (0.5101071,0.3060026) (0.5108775,0.3053708) (0.5116157,0.3047046) (0.5123202,0.3040061) (0.5129895,0.3032775) (0.5136224,0.3025210) (0.5142177,0.3017389) (0.5147742,0.3009334) (0.5152909,0.3001069) (0.5157669,0.2992617) (0.5162013,0.2984001) (0.5165935,0.2975246) (0.5169427,0.2966374) (0.5172485,0.2957410) (0.5175103,0.2948376) (0.5177279,0.2939296) (0.5179010,0.2930192) (0.5180295,0.2921088) (0.5181135,0.2912005) (0.5181530,0.2902965) (0.5181481,0.2893990) (0.5180994,0.2885100) (0.5180070,0.2876315) (0.5178717,0.2867656) (0.5176940,0.2859142) (0.5174746,0.2850790) (0.5172144,0.2842620) (0.5169144,0.2834647) (0.5165755,0.2826888) (0.5161990,0.2819359) (0.5157859,0.2812075) (0.5153377,0.2805050) (0.5148557,0.2798296) (0.5143414,0.2791827) (0.5137963,0.2785654) (0.5132221,0.2779788) (0.5126205,0.2774239) (0.5119931,0.2769016) (0.5113418,0.2764127) (0.5106684,0.2759580) (0.5099749,0.2755381) (0.5092631,0.2751536) (0.5085351,0.2748050) (0.5077928,0.2744926) (0.5070382,0.2742169) (0.5062735,0.2739781) (0.5055006,0.2737763) (0.5047215,0.2736115) (0.5039384,0.2734839) (0.5031532,0.2733933) (0.5023680,0.2733395) (0.5015847,0.2733223) (0.5008054,0.2733413) (0.5000320,0.2733963) (0.4992664,0.2734867) (0.4985105,0.2736120) (0.4977660,0.2737716) (0.4970348,0.2739648) (0.4963186,0.2741909) (0.4956190,0.2744490) (0.4949376,0.2747383) (0.4942759,0.2750578) (0.4936355,0.2754066) (0.4930175,0.2757836) (0.4924235,0.2761876) (0.4918545,0.2766176) (0.4913117,0.2770724) (0.4907962,0.2775506) (0.4903090,0.2780509) (0.4898509,0.2785720) (0.4894227,0.2791126) (0.4890252,0.2796712) (0.4886589,0.2802463) (0.4883244,0.2808365) (0.4880222,0.2814402) (0.4877525,0.2820560) (0.4875158,0.2826822) (0.4873121,0.2833173) (0.4871415,0.2839597) (0.4870041,0.2846078) (0.4868998,0.2852600) (0.4868284,0.2859147) (0.4867898,0.2865703) (0.4867836,0.2872252) (0.4868094,0.2878779) (0.4868668,0.2885268) (0.4869552,0.2891702) (0.4870741,0.2898068) (0.4872229,0.2904349) (0.4874007,0.2910532) (0.4876069,0.2916602) (0.4878405,0.2922544) (0.4881008,0.2928345) (0.4883867,0.2933993) (0.4886974,0.2939474) (0.4890317,0.2944776) (0.4893886,0.2949887) (0.4897671,0.2954797) (0.4901659,0.2959494) (0.4905839,0.2963970) (0.4910199,0.2968214) (0.4914726,0.2972218) (0.4919409,0.2975975) (0.4924235,0.2979476) (0.4929190,0.2982716) (0.4934262,0.2985689) (0.4939437,0.2988390) (0.4944704,0.2990815) (0.4950047,0.2992959) (0.4955455,0.2994821) (0.4960913,0.2996398) (0.4966410,0.2997689) (0.4971930,0.2998693) (0.4977462,0.2999412) (0.4982993,0.2999845) (0.4988509,0.2999995) (0.4993998,0.2999865) (0.4999447,0.2999456) (0.5004845,0.2998774) (0.5010179,0.2997823) (0.5015437,0.2996608) (0.5020608,0.2995134) (0.5025681,0.2993409) (0.5030644,0.2991438) (0.5035487,0.2989231) (0.5040199,0.2986794) (0.5044772,0.2984137) (0.5049195,0.2981269) (0.5053460,0.2978198) (0.5057557,0.2974936) (0.5061479,0.2971491) (0.5065218,0.2967875) (0.5068766,0.2964099) (0.5072117,0.2960174) (0.5075265,0.2956110) (0.5078203,0.2951921) (0.5080927,0.2947617) (0.5083432,0.2943211) (0.5085714,0.2938715) (0.5087769,0.2934140) (0.5089594,0.2929499) (0.5091187,0.2924804) (0.5092546,0.2920066) (0.5093669,0.2915299) (0.5094556,0.2910514) (0.5095207,0.2905723) (0.5095621,0.2900938) (0.5095801,0.2896169) (0.5095746,0.2891430) (0.5095461,0.2886730) (0.5094946,0.2882080) (0.5094206,0.2877491) (0.5093243,0.2872974) (0.5092063,0.2868539) (0.5090670,0.2864195) (0.5089069,0.2859951) (0.5087266,0.2855817) (0.5085266,0.2851801) (0.5083078,0.2847912) (0.5080707,0.2844157) (0.5078161,0.2840544) (0.5075448,0.2837080) (0.5072576,0.2833770) (0.5069554,0.2830623) (0.5066390,0.2827642) (0.5063094,0.2824834) (0.5059675,0.2822203) (0.5056142,0.2819753) (0.5052505,0.2817489) (0.5048775,0.2815413) (0.5044961,0.2813529) (0.5041074,0.2811839) (0.5037124,0.2810345) (0.5033122,0.2809048) (0.5029077,0.2807950) (0.5025001,0.2807051) (0.5020904,0.2806351) (0.5016796,0.2805849) (0.5012688,0.2805546) (0.5008589,0.2805439) (0.5004511,0.2805528) (0.5000462,0.2805809) (0.4996454,0.2806280) (0.4992494,0.2806939) (0.4988593,0.2807782) (0.4984760,0.2808806) (0.4981003,0.2810005) (0.4977331,0.2811377) (0.4973752,0.2812915) (0.4970274,0.2814615) (0.4966905,0.2816471) (0.4963652,0.2818478) (0.4960522,0.2820629) (0.4957520,0.2822918) (0.4954654,0.2825338) (0.4951928,0.2827883) (0.4949348,0.2830544) (0.4946919,0.2833316) (0.4944645,0.2836189) (0.4942530,0.2839157) (0.4940576,0.2842212) (0.4938788,0.2845345) (0.4937168,0.2848548) (0.4935718,0.2851814) (0.4934440,0.2855133) (0.4933334,0.2858497) (0.4932402,0.2861898) (0.4931643,0.2865327) (0.4931059,0.2868776) (0.4930648,0.2872236) (0.4930409,0.2875699) (0.4930341,0.2879156) (0.4930442,0.2882599) (0.4930711,0.2886020) (0.4931144,0.2889411) (0.4931739,0.2892764) (0.4932493,0.2896070) (0.4933401,0.2899323) (0.4934461,0.2902515) (0.4935667,0.2905638) (0.4937016,0.2908686) (0.4938503,0.2911651) (0.4940122,0.2914528) (0.4941868,0.2917310) (0.4943736,0.2919991) (0.4945719,0.2922565) (0.4947813,0.2925028) (0.4950010,0.2927373) (0.4952305,0.2929597) (0.4954691,0.2931695) (0.4957161,0.2933663) (0.4959708,0.2935497) (0.4962326,0.2937195) (0.4965007,0.2938752) (0.4967746,0.2940168) (0.4970534,0.2941439) (0.4973364,0.2942564) (0.4976229,0.2943541) (0.4979123,0.2944370) (0.4982038,0.2945050) (0.4984966,0.2945580) (0.4987902,0.2945962) (0.4990837,0.2946195) (0.4993765,0.2946280) (0.4996679,0.2946219) (0.4999572,0.2946012) (0.5002438,0.2945664) (0.5005269,0.2945174) (0.5008060,0.2944547) (0.5010805,0.2943785) (0.5013497,0.2942891) (0.5016130,0.2941869) (0.5018700,0.2940723)};
\fill[ink] (0.3500000,0.1732051) circle(.011);
\fill[ink] (.5,.2886751) circle(.008);\end{scope}\end{scope}
\begin{scope}[xshift=4.8cm]\begin{scope}[scale=4.0]
\fill[soft] (0,0)--(1,0)--(.5,.8660254)--cycle;
\draw[ink!75] (0,0)--(1,0)--(.5,.8660254)--cycle;
\node[font=\small] at (.5,1.08){$s=1.0$};
\draw[ink,line width=.85pt] plot coordinates {(0.3500000,0.1732051) (0.3445481,0.1785791) (0.3396990,0.1843096) (0.3354598,0.1903970) (0.3318346,0.1968402) (0.3288246,0.2036365) (0.3264280,0.2107813) (0.3246402,0.2182682) (0.3234541,0.2260882) (0.3228600,0.2342300) (0.3228458,0.2426797) (0.3233974,0.2514206) (0.3244984,0.2604333) (0.3261305,0.2696953) (0.3282739,0.2791811) (0.3309067,0.2888624) (0.3340061,0.2987081) (0.3375476,0.3086841) (0.3415060,0.3187541) (0.3458551,0.3288791) (0.3505682,0.3390181) (0.3556183,0.3491286) (0.3609781,0.3591663) (0.3666206,0.3690862) (0.3725192,0.3788424) (0.3786479,0.3883889) (0.3849816,0.3976799) (0.3914963,0.4066702) (0.3981692,0.4153155) (0.4049790,0.4235729) (0.4119059,0.4314014) (0.4189318,0.4387618) (0.4260404,0.4456174) (0.4332172,0.4519340) (0.4404493,0.4576800) (0.4477255,0.4628270) (0.4550364,0.4673497) (0.4623741,0.4712257) (0.4697317,0.4744360) (0.4771040,0.4769651) (0.4844865,0.4788007) (0.4918753,0.4799339) (0.4992673,0.4803592) (0.5066597,0.4800746) (0.5140494,0.4790814) (0.5214334,0.4773844) (0.5288080,0.4749919) (0.5361690,0.4719154) (0.5435111,0.4681700) (0.5508279,0.4637738) (0.5581117,0.4587486) (0.5653533,0.4531192) (0.5725420,0.4469133) (0.5796653,0.4401619) (0.5867088,0.4328987) (0.5936568,0.4251600) (0.6004913,0.4169843) (0.6071929,0.4084126) (0.6137407,0.3994872) (0.6201122,0.3902523) (0.6262835,0.3807528) (0.6322298,0.3710345) (0.6379252,0.3611434) (0.6433431,0.3511254) (0.6484567,0.3410258) (0.6532387,0.3308889) (0.6576621,0.3207579) (0.6617001,0.3106738) (0.6653263,0.3006762) (0.6685152,0.2908019) (0.6712424,0.2810854) (0.6734846,0.2715585) (0.6752197,0.2622501) (0.6764275,0.2531861) (0.6770893,0.2443897) (0.6771885,0.2358808) (0.6767104,0.2276768) (0.6756424,0.2197921) (0.6739745,0.2122384) (0.6716990,0.2050251) (0.6688108,0.1981593) (0.6653076,0.1916459) (0.6611899,0.1854878) (0.6564612,0.1796866) (0.6511281,0.1742420) (0.6452004,0.1691528) (0.6386913,0.1644166) (0.6316174,0.1600304) (0.6239986,0.1559902) (0.6158587,0.1522919) (0.6072246,0.1489309) (0.5981269,0.1459027) (0.5885996,0.1432026) (0.5786797,0.1408262) (0.5684076,0.1387692) (0.5578264,0.1370278) (0.5469819,0.1355986) (0.5359220,0.1344787) (0.5246965,0.1336656) (0.5133567,0.1331578) (0.5019550,0.1329539) (0.4905442,0.1330537) (0.4791773,0.1334574) (0.4679068,0.1341657) (0.4567845,0.1351803) (0.4458607,0.1365033) (0.4351840,0.1381374) (0.4248011,0.1400859) (0.4147559,0.1423524) (0.4050897,0.1449412) (0.3958408,0.1478565) (0.3870443,0.1511030) (0.3787319,0.1546852) (0.3709319,0.1586078) (0.3636689,0.1628751) (0.3569644,0.1674910) (0.3508359,0.1724590) (0.3452978,0.1777816) (0.3403612,0.1834607) (0.3360336,0.1894966) (0.3323197,0.1958885) (0.3292210,0.2026341) (0.3267361,0.2097291) (0.3248609,0.2171671) (0.3235885,0.2249396) (0.3229096,0.2330358) (0.3228126,0.2414420) (0.3232835,0.2501421) (0.3243063,0.2591169) (0.3258630,0.2683443) (0.3279339,0.2777995) (0.3304975,0.2874545) (0.3335310,0.2972785) (0.3370103,0.3072379) (0.3409101,0.3172967) (0.3452045,0.3274162) (0.3498667,0.3375558) (0.3548698,0.3476731) (0.3601866,0.3577241) (0.3657899,0.3676638) (0.3716532,0.3774465) (0.3777503,0.3870261) (0.3840560,0.3963568) (0.3905460,0.4053932) (0.3971974,0.4140909) (0.4039887,0.4224069) (0.4108998,0.4302998) (0.4179125,0.4377301) (0.4250101,0.4446607) (0.4321779,0.4510570) (0.4394027,0.4568872)};
\fill[ink] (0.3500000,0.1732051) circle(.011);
\fill[ink] (.5,.2886751) circle(.008);\end{scope}\end{scope}
\begin{scope}[xshift=9.6cm]\begin{scope}[scale=4.0]
\fill[soft] (0,0)--(1,0)--(.5,.8660254)--cycle;
\draw[ink!75] (0,0)--(1,0)--(.5,.8660254)--cycle;
\node[font=\small] at (.5,1.08){$s=1.3$};
\draw[copper,line width=.85pt] plot coordinates {(0.3500000,0.1732051) (0.3428733,0.1788898) (0.3365140,0.1850624) (0.3309452,0.1917286) (0.3261837,0.1988925) (0.3222408,0.2065561) (0.3191223,0.2147190) (0.3168285,0.2233778) (0.3153549,0.2325255) (0.3146920,0.2421515) (0.3148259,0.2522405) (0.3157386,0.2627724) (0.3174078,0.2737222) (0.3198077,0.2850588) (0.3229089,0.2967458) (0.3266787,0.3087405) (0.3310816,0.3209945) (0.3360794,0.3334532) (0.3416319,0.3460566) (0.3476967,0.3587393) (0.3542306,0.3714314) (0.3611891,0.3840586) (0.3685278,0.3965438) (0.3762025,0.4088073) (0.3841699,0.4207684) (0.3923885,0.4323458) (0.4008187,0.4434594) (0.4094238,0.4540308) (0.4181702,0.4639844) (0.4270277,0.4732485) (0.4359703,0.4817558) (0.4449758,0.4894443) (0.4540262,0.4962578) (0.4631077,0.5021460) (0.4722102,0.5070653) (0.4813274,0.5109782) (0.4904561,0.5138543) (0.4995961,0.5156694) (0.5087491,0.5164060) (0.5179186,0.5160530) (0.5271091,0.5146059) (0.5363251,0.5120662) (0.5455707,0.5084420) (0.5548490,0.5037478) (0.5641610,0.4980043) (0.5735053,0.4912388) (0.5828773,0.4834851) (0.5922685,0.4747838) (0.6016665,0.4651822) (0.6110541,0.4547342) (0.6204094,0.4435004) (0.6297053,0.4315478) (0.6389098,0.4189493) (0.6479860,0.4057835) (0.6568920,0.3921337) (0.6655821,0.3780871) (0.6740064,0.3637338) (0.6821120,0.3491656) (0.6898436,0.3344745) (0.6971442,0.3197519) (0.7039561,0.3050865) (0.7102212,0.2905634) (0.7158825,0.2762628) (0.7208845,0.2622586) (0.7251735,0.2486179) (0.7286990,0.2353999) (0.7314132,0.2226558) (0.7332719,0.2104282) (0.7342350,0.1987511) (0.7342660,0.1876506) (0.7333327,0.1771447) (0.7314068,0.1672443) (0.7284647,0.1579536) (0.7244868,0.1492710) (0.7194583,0.1411901) (0.7133692,0.1337000) (0.7062142,0.1267866) (0.6979939,0.1204334) (0.6887145,0.1146218) (0.6783886,0.1093322) (0.6670360,0.1045444) (0.6546839,0.1002382) (0.6413678,0.0963937) (0.6271321,0.0929921) (0.6120304,0.0900156) (0.5961260,0.0874477) (0.5794918,0.0852737) (0.5622107,0.0834807) (0.5443749,0.0820577) (0.5260854,0.0809959) (0.5074509,0.0802883) (0.4885868,0.0799305) (0.4696135,0.0799199) (0.4506544,0.0802566) (0.4318344,0.0809424) (0.4132773,0.0819818) (0.3951040,0.0833813) (0.3774306,0.0851495) (0.3603662,0.0872974) (0.3440116,0.0898379) (0.3284577,0.0927860) (0.3137851,0.0961585) (0.3000631,0.0999741) (0.2873496,0.1042534) (0.2756913,0.1090180) (0.2651242,0.1142913) (0.2556743,0.1200973) (0.2473580,0.1264611) (0.2401839,0.1334079) (0.2341530,0.1409629) (0.2292598,0.1491507) (0.2254936,0.1579947) (0.2228390,0.1675165) (0.2212764,0.1777351) (0.2207828,0.1886662) (0.2213320,0.2003210) (0.2228947,0.2127055) (0.2254387,0.2258196) (0.2289289,0.2396557) (0.2333269,0.2541981) (0.2385908,0.2694218) (0.2446754,0.2852918) (0.2515314,0.3017625) (0.2591058,0.3187770) (0.2673416,0.3362672) (0.2761779,0.3541536) (0.2855504,0.3723461) (0.2953918,0.3907449) (0.3056327,0.4092413) (0.3162021,0.4277199) (0.3270291,0.4460600) (0.3380434,0.4641382) (0.3491775,0.4818304) (0.3603673,0.4990143) (0.3715535,0.5155717) (0.3826833,0.5313903) (0.3937107,0.5463660) (0.4045977,0.5604039) (0.4153150,0.5734193) (0.4258417,0.5853386) (0.4361659,0.5960991) (0.4462840,0.6056489) (0.4562006,0.6139463) (0.4659276,0.6209589) (0.4754836,0.6266627) (0.4848930,0.6310409) (0.4941847,0.6340827) (0.5033917,0.6357826) (0.5125496,0.6361392) (0.5216958,0.6351543) (0.5308685,0.6328328) (0.5401057,0.6291824) (0.5494443,0.6242129) (0.5589192,0.6179370) (0.5685622,0.6103708) (0.5784014,0.6015338) (0.5884600,0.5914504) (0.5987557,0.5801502) (0.6092998,0.5676699) (0.6200964,0.5540533) (0.6311419,0.5393534) (0.6424240,0.5236327) (0.6539217,0.5069639) (0.6656052,0.4894308) (0.6774352,0.4711280) (0.6893638,0.4521608) (0.7013347,0.4326442) (0.7132843,0.4127017) (0.7251422,0.3924638) (0.7368332,0.3720653) (0.7482785,0.3516429) (0.7593974,0.3313325) (0.7701088,0.3112662) (0.7803333,0.2915695) (0.7899946,0.2723585) (0.7990205,0.2537377) (0.8073443,0.2357981) (0.8149056,0.2186160) (0.8216503,0.2022521) (0.8275310,0.1867518) (0.8325069,0.1721450) (0.8365426,0.1584475) (0.8396082,0.1456619) (0.8416782,0.1337791) (0.8427303,0.1227803) (0.8427450,0.1126381) (0.8417045,0.1033190) (0.8395920,0.0947843) (0.8363910,0.0869918) (0.8320851,0.0798975) (0.8266577,0.0734562) (0.8200919,0.0676228) (0.8123710,0.0623529) (0.8034784,0.0576037) (0.7933990,0.0533342) (0.7821193,0.0495055) (0.7696293,0.0460814) (0.7559232,0.0430283) (0.7410011,0.0403150) (0.7248706,0.0379134) (0.7075486,0.0357977) (0.6890627,0.0339450) (0.6694530,0.0323347) (0.6487737,0.0309486) (0.6270940,0.0297710) (0.6044993,0.0287883) (0.5810911,0.0279890) (0.5569873,0.0273635) (0.5323210,0.0269043) (0.5072388,0.0266056) (0.4818987,0.0264634) (0.4564674,0.0264754) (0.4311161,0.0266409) (0.4060174,0.0269610) (0.3813405,0.0274381) (0.3572479,0.0280765) (0.3338911,0.0288818) (0.3114076,0.0298613) (0.2899180,0.0310241) (0.2695246,0.0323806) (0.2503100,0.0339431) (0.2323373,0.0357255) (0.2156501,0.0377436) (0.2002746,0.0400148) (0.1862204,0.0425585) (0.1734835,0.0453961) (0.1620481,0.0485508) (0.1518889,0.0520480) (0.1429740,0.0559149) (0.1352666,0.0601809) (0.1287270,0.0648774) (0.1233145,0.0700379) (0.1189889,0.0756978) (0.1157114,0.0818939) (0.1134454,0.0886651) (0.1121577,0.0960514) (0.1118179,0.1040935) (0.1123993,0.1128331) (0.1138785,0.1223112) (0.1162351,0.1325687) (0.1194513,0.1436443) (0.1235110,0.1555744) (0.1283997,0.1683918) (0.1341030,0.1821239) (0.1406058,0.1967921) (0.1478913,0.2124097) (0.1559399,0.2289802) (0.1647281,0.2464965) (0.1742274,0.2649382) (0.1844031,0.2842711) (0.1952137,0.3044454) (0.2066101,0.3253950) (0.2185354,0.3470371) (0.2309246,0.3692724) (0.2437051,0.3919857) (0.2567975,0.4150474) (0.2701166,0.4383156) (0.2835731,0.4616395) (0.2970757,0.4848618) (0.3105329,0.5078235) (0.3238557,0.5303675) (0.3369595,0.5523425) (0.3497668,0.5736073) (0.3622085,0.5940334) (0.3742261,0.6135084) (0.3857723,0.6319370) (0.3968117,0.6492430) (0.4073213,0.6653683) (0.4172896,0.6802734) (0.4267165,0.6939353) (0.4356119,0.7063460) (0.4439949,0.7175106) (0.4518922,0.7274442) (0.4593369,0.7361705) (0.4663675,0.7437184) (0.4730264,0.7501208) (0.4793589,0.7554117) (0.4854124,0.7596252) (0.4912357,0.7627935) (0.4968783,0.7649462) (0.5023897,0.7661089) (0.5078196,0.7663028) (0.5132168,0.7655442) (0.5186299,0.7638442) (0.5241065,0.7612087) (0.5296934,0.7576382) (0.5354363,0.7531284) (0.5413799,0.7476707) (0.5475675,0.7412520) (0.5540409,0.7338563) (0.5608402,0.7254648) (0.5680033,0.7160575) (0.5755658,0.7056137) (0.5835599,0.6941137) (0.5920144,0.6815404) (0.6009534,0.6678804) (0.6103962,0.6531262) (0.6203554,0.6372777) (0.6308371,0.6203443) (0.6418392,0.6023469) (0.6533507,0.5833193) (0.6653511,0.5633099) (0.6778093,0.5423831) (0.6906840,0.5206202) (0.7039227,0.4981194) (0.7174628,0.4749955) (0.7312319,0.4513787) (0.7451490,0.4274128) (0.7591264,0.4032523) (0.7730713,0.3790590) (0.7868884,0.3549979) (0.8004826,0.3312333) (0.8137609,0.3079236) (0.8266359,0.2852179) (0.8390269,0.2632512) (0.8508627,0.2421421) (0.8620824,0.2219896) (0.8726364,0.2028723) (0.8824866,0.1848474) (0.8916064,0.1679513) (0.8999797,0.1522007) (0.9076000,0.1375940) (0.9144691,0.1241141) (0.9205956,0.1117303) (0.9259935,0.1004011) (0.9306803,0.0900766) (0.9346760,0.0807010) (0.9380016,0.0722143) (0.9406778,0.0645546) (0.9427244,0.0576596) (0.9441591,0.0514676) (0.9449967,0.0459187) (0.9452491,0.0409554) (0.9449245,0.0365234) (0.9440269,0.0325717) (0.9425565,0.0290529) (0.9405087,0.0259231) (0.9378750,0.0231420) (0.9346422,0.0206731) (0.9307926,0.0184828) (0.9263044,0.0165411) (0.9211516,0.0148207) (0.9153039,0.0132972) (0.9087276,0.0119486) (0.9013853,0.0107553) (0.8932368,0.0096999) (0.8842391,0.0087668) (0.8743474,0.0079421) (0.8635160,0.0072137) (0.8516985,0.0065707) (0.8388498,0.0060034) (0.8249267,0.0055033) (0.8098898,0.0050631) (0.7937048,0.0046762) (0.7763444,0.0043367) (0.7577908,0.0040397) (0.7380370,0.0037807) (0.7170895,0.0035560) (0.6949703,0.0033621) (0.6717188,0.0031962) (0.6473936,0.0030560) (0.6220735,0.0029392) (0.5958584,0.0028443) (0.5688691,0.0027697) (0.5412466,0.0027144) (0.5131499,0.0026775) (0.4847533,0.0026584) (0.4562427,0.0026566) (0.4278113,0.0026721) (0.3996539,0.0027050) (0.3719622,0.0027554) (0.3449189,0.0028240) (0.3186931,0.0029113) (0.2934357,0.0030183) (0.2692762,0.0031461) (0.2463199,0.0032961) (0.2246474,0.0034699) (0.2043139,0.0036693) (0.1853503,0.0038964) (0.1677652,0.0041535) (0.1515471,0.0044435) (0.1366673,0.0047692) (0.1230829,0.0051340) (0.1107401,0.0055417) (0.0995769,0.0059963) (0.0895259,0.0065026) (0.0805164,0.0070656) (0.0724773,0.0076908) (0.0653376,0.0083846) (0.0590289,0.0091537) (0.0534855,0.0100057) (0.0486460,0.0109488) (0.0444531,0.0119922) (0.0408544,0.0131459) (0.0378023,0.0144209) (0.0352543,0.0158293) (0.0331728,0.0173841) (0.0315250,0.0190999) (0.0302828,0.0209925) (0.0294224,0.0230789) (0.0289245,0.0253781) (0.0287739,0.0279104) (0.0289590,0.0306979) (0.0294721,0.0337647) (0.0303091,0.0371367) (0.0314688,0.0408419) (0.0329535,0.0449105) (0.0347684,0.0493746) (0.0369214,0.0542687) (0.0394232,0.0596294) (0.0422867,0.0654956) (0.0455273,0.0719079) (0.0491624,0.0789091) (0.0532111,0.0865434) (0.0576941,0.0948564) (0.0626331,0.1038944) (0.0680507,0.1137040) (0.0739696,0.1243313) (0.0804123,0.1358210) (0.0874003,0.1482155) (0.0949533,0.1615533) (0.1030884,0.1758682) (0.1118192,0.1911871) (0.1211547,0.2075286) (0.1310984,0.2249013) (0.1416469,0.2433017) (0.1527891,0.2627124) (0.1645049,0.2831004) (0.1767648,0.3044159) (0.1895287,0.3265908) (0.2027461,0.3495385) (0.2163556,0.3731536) (0.2302860,0.3973133) (0.2444566,0.4218786) (0.2587793,0.4466969) (0.2731599,0.4716057) (0.2875011,0.4964361) (0.3017044,0.5210176) (0.3156733,0.5451831) (0.3293161,0.5687731) (0.3425480,0.5916407) (0.3552937,0.6136547) (0.3674890,0.6347034) (0.3790821,0.6546961) (0.3900338,0.6735641) (0.4003180,0.6912611) (0.4099210,0.7077621) (0.4188404,0.7230615) (0.4270841,0.7371716) (0.4346689,0.7501194) (0.4416185,0.7619444) (0.4479627,0.7726956) (0.4537351,0.7824289) (0.4589724,0.7912051) (0.4637128,0.7990874) (0.4679952,0.8061399) (0.4718584,0.8124258) (0.4753403,0.8180068) (0.4784774,0.8229416) (0.4813046,0.8272858) (0.4838552,0.8310911) (0.4861600,0.8344051) (0.4882482,0.8372713) (0.4901467,0.8397291) (0.4918806,0.8418137) (0.4934730,0.8435561) (0.4949456,0.8449837) (0.4963183,0.8461201) (0.4976098,0.8469854) (0.4988372,0.8475964) (0.5000171,0.8479666) (0.5011648,0.8481066) (0.5022950,0.8480243) (0.5034219,0.8477247) (0.5045591,0.8472101) (0.5057202,0.8464806) (0.5069185,0.8455336) (0.5081672,0.8443641) (0.5094800,0.8429650) (0.5108704,0.8413265) (0.5123527,0.8394367) (0.5139415,0.8372813) (0.5156521,0.8348436) (0.5175003,0.8321045) (0.5195031,0.8290425) (0.5216781,0.8256336) (0.5240439,0.8218512) (0.5266203,0.8176665) (0.5294281,0.8130480) (0.5324891,0.8079616) (0.5358264,0.8023710) (0.5394640,0.7962376) (0.5434272,0.7895204) (0.5477418,0.7821770) (0.5524348,0.7741629) (0.5575332,0.7654328) (0.5630647,0.7559406) (0.5690566,0.7456406) (0.5755354,0.7344876) (0.5825268,0.7224388) (0.5900542,0.7094540) (0.5981387,0.6954979) (0.6067978,0.6805410) (0.6160443,0.6645614) (0.6258857,0.6475467) (0.6363230,0.6294959) (0.6473493,0.6104212) (0.6589491,0.5903498) (0.6710975,0.5693252) (0.6837590,0.5474091) (0.6968875,0.5246816) (0.7104260,0.5012414) (0.7243071,0.4772057) (0.7384535,0.4527082) (0.7527797,0.4278974) (0.7671933,0.4029333) (0.7815975,0.3779833) (0.7958937,0.3532184) (0.8099841,0.3288080) (0.8237745,0.3049153) (0.8371769,0.2816925) (0.8501119,0.2592768) (0.8625108,0.2377875) (0.8743168,0.2173227) (0.8854856,0.1979588) (0.8959862,0.1797494) (0.9058000,0.1627265) (0.9149203,0.1469014) (0.9233510,0.1322671) (0.9311053,0.1188003) (0.9382042,0.1064643) (0.9446747,0.0952119) (0.9505486,0.0849876) (0.9558607,0.0757306) (0.9606478,0.0673765) (0.9649477,0.0598593) (0.9687978,0.0531133) (0.9722352,0.0470738) (0.9752951,0.0416785) (0.9780114,0.0368680) (0.9804157,0.0325864) (0.9825376,0.0287814) (0.9844042,0.0254045) (0.9860404,0.0224114) (0.9874688,0.0197612) (0.9887097,0.0174169) (0.9897813,0.0153450) (0.9907000,0.0135152) (0.9914799,0.0119002) (0.9921338,0.0104758) (0.9926725,0.0092200) (0.9931053,0.0081134) (0.9934403,0.0071387) (0.9936841,0.0062804) (0.9938420,0.0055248) (0.9939184,0.0048599) (0.9939163,0.0042749) (0.9938380,0.0037602) (0.9936845,0.0033076) (0.9934560,0.0029096) (0.9931518,0.0025595) (0.9927703,0.0022518) (0.9923089,0.0019813) (0.9917643,0.0017435) (0.9911319,0.0015344) (0.9904067,0.0013507) (0.9895822,0.0011892) (0.9886512,0.0010472) (0.9876054,0.0009224) (0.9864354,0.0008127) (0.9851303,0.0007163) (0.9836785,0.0006315) (0.9820666,0.0005570) (0.9802799,0.0004915) (0.9783023,0.0004339) (0.9761158,0.0003833) (0.9737010,0.0003388) (0.9710364,0.0002996) (0.9680986,0.0002652) (0.9648622,0.0002349) (0.9612993,0.0002082) (0.9573801,0.0001848) (0.9530721,0.0001641) (0.9483406,0.0001460) (0.9431480,0.0001300) (0.9374544,0.0001159) (0.9312174,0.0001035) (0.9243919,0.0000926) (0.9169305,0.0000830) (0.9087836,0.0000745) (0.8999000,0.0000671) (0.8902266,0.0000605) (0.8797097,0.0000547) (0.8682953,0.0000496) (0.8559299,0.0000451) (0.8425617,0.0000412) (0.8281418,0.0000377) (0.8126256,0.0000347) (0.7959744,0.0000320) (0.7781574,0.0000297) (0.7591536,0.0000276)};
\fill[ink] (0.3500000,0.1732051) circle(.011);
\fill[ink] (.5,.2886751) circle(.008);\end{scope}\end{scope}
\end{tikzpicture}\caption{Cyclic replicator trajectories from the same initial law $(.55,.25,.20)$. Left: attraction toward the center. Middle: a closed product level. Right: outward motion toward the boundary. Curves are numerically integrated from the stated vector fields.}\end{figure}

\section{Planar linearization at the center}\index{Jacobian matrix}
Use independent coordinates $x=p_A$ and $y=p_B$, so $p_C=1-x-y$. For a planar system $(x',y')=(F,G)$, the Jacobian at an equilibrium is
\[
J=\begin{pmatrix}\partial_xF&\partial_yF\\\partial_xG&\partial_yG\end{pmatrix}.
\]
For the cyclic family, substituting the fitnesses and evaluating at $(1/3,1/3)$ yields
\begin{equation}\label{eq:cyclicjacobian}
J_s=\frac13\begin{pmatrix}s&s+1\\-(s+1)&-1\end{pmatrix}.
\end{equation}
For instance, at the center $f_A-\bar f=0$, so differentiating $x(f_A-\bar f)$ leaves $x\,\partial(f_A-\bar f)$. Carrying this out for the four entries gives the displayed matrix.

Its trace and determinant are
\[
\operatorname{tr}J_s=\frac{s-1}{3},\qquad
\det J_s=\frac{s^2+s+1}{9}.
\]
The characteristic equation gives
\begin{equation}\label{eq:cycliceigenvalues}
\lambda_\pm=\frac{s-1}{6}\ \pm\ i\frac{\sqrt3(s+1)}6.
\end{equation}
For $s>0$, the imaginary parts are nonzero. The real parts are negative for $s<1$ and positive for $s>1$, giving a locally attracting spiral and repelling spiral, respectively. At $s=1$, linearization produces a center. The conserved product establishes the periodic behavior of the nonlinear system in that case.

\begin{example}[Damping and oscillation]
For $s=0.7$, estimate the decay rate and angular frequency near the center.
\begin{solution}
The eigenvalues have real part $-0.05$ and imaginary parts approximately $\pm0.49075$. In the linearized system a perturbation has an exponentially decaying envelope $e^{-0.05t}$ while it rotates with angular frequency about $0.49075$. The corresponding small-oscillation period is approximately $2\pi/0.49075=12.80$.
\end{solution}
\end{example}

\section{Finding equilibria and testing invasion}
At a replicator equilibrium, every present type has fitness equal to the population mean. Types with zero proportion impose no stationarity condition on their own fitness values, because their equations already vanish through the factor $p_i$.

For an interior matrix equilibrium, all types are present. The condition can therefore be written
\[
M\p=c\ones,\qquad\ones^T\p=1,
\]
where $c$ is the common fitness. This is a linear system for the probabilities together with $c$. A solution represents an interior equilibrium when all its probability coordinates are positive. Singular systems can have no solution or a family of solutions, so positivity and uniqueness must be checked after solving the equations.

For the cyclic family, subtracting pairs of the equal-fitness equations gives the uniform solution. One can also see uniqueness by restricting $M_s$ to the tangent plane: the corresponding linear equal-fitness system has nonzero determinant for $s>0$. The three vertices remain equilibria independently of that interior calculation.

Stability on an invariant edge concerns perturbations within that edge. A missing type can introduce another direction of perturbation. If a face equilibrium is $q$ and type $k$ is absent, its equation near the face has leading term
\[
p_k'\approx p_k\bigl(f_k(q)-\bar f(q)\bigr).
\]
The difference is the type's invasion rate at the resident equilibrium. A positive value makes a small introduction grow to first order, even if the resident state attracts along its own face.

\begin{example}[An edge attractor with an invading third type]
Consider
\[
M=\begin{pmatrix}0&2&0\\2&0&0\\3&3&0\end{pmatrix}.
\]
Analyze the AB equilibrium $q=(1/2,1/2,0)$ along the edge and under introduction of C.
\begin{solution}
On the edge, $f_A=2(1-x)$ and $f_B=2x$, giving $x'=x(1-x)(2-4x)$. The equilibrium $x=1/2$ attracts, with derivative $-1$.

At the same state in the full model, the fitnesses are $(1,1,3)$ and mean fitness is 1. A rare C type has invasion rate $3-1=2$, so its initial growth is approximately $p_C'=2p_C$. The equilibrium attracts along the edge and repels in the direction introducing C. It is unstable in the full simplex.
\end{solution}
\end{example}

At a pure resident state $e_j$, the invasion rate of type $i\ne j$ is particularly simple: $m_{ij}-m_{jj}$. These are the diagonal entries of the tangent linearization in the rare-type coordinates. In balanced cyclic competition, a pure A population resists rare B with rate $-1$, while rare C invades with rate $+1$. Each vertex similarly has one attracting and one repelling edge direction. These local directions organize the circulation around the boundary.

Thus the triangle carries several related stability questions: behavior within each invariant edge, invasion of the missing type, and motion near an interior equilibrium. Their coordinate calculations are elementary, but they describe distinct allowable perturbations of the population state.

\section{Distributions, support, and state-space dimension}
The \emph{support}\index{support} of a probability vector is the set of indices with positive coordinates. Smooth replicator dynamics preserve support at finite times. A face of the simplex is therefore both a geometric object and the state space of a population restricted to its initially present types.

The same normalized vector also describes the result of sampling one individual. The triangle is a space of categorical distributions, and the replicator field supplies one rule for moving through it. A distribution can be studied independently of that rule, which permits the probability and information constructions developed next.

\FloatBarrier
\section{Exercises}
\begin{exercises}
\item Plot $(0.2,0.5,0.3)$ using \eqref{eq:barycentric}. Identify the absent type on each edge of the triangle.
\item Show that the three-category simplex has two independent coordinates and that its tangent velocities have components summing to zero.
\item For $\p=(1/2,1/4,1/4)$ and $\f=(2,0,1)$, compute $\bar f$ and $\p'$.
\item Derive the $n$-type replicator equation from $N_i'=f_iN_i$. Explain the subtraction of mean fitness.
\item Verify \eqref{eq:positivity} by differentiating its right-hand side. State the hypotheses that keep its exponent finite on a finite time interval.
\item Reduce $M_s$ to the BC and CA edges. Determine the directions around the boundary cycle.
\item At $s=1$, prove algebraically that the mean fitness is zero for every state. Is every state an equilibrium?
\item Verify $\bar f=(1-s)Q$ and derive \eqref{eq:productderivative}.
\item Prove the identity expressing $1-3Q$ as a sum of squared coordinate differences. Use it to determine its zero set.
\item Prove $p_Ap_Bp_C\leq1/27$ for normalized nonnegative coordinates. Give both a calculus proof and an arithmetic-geometric-mean proof.
\item Compute the trace, determinant, and eigenvalues of \eqref{eq:cyclicjacobian} when $s=1$. Explain what extra information the conserved product supplies.
\item Starting from the cyclic fitnesses, calculate all four entries of the reduced Jacobian at the center.
\item For $s=1.3$, calculate the real and imaginary parts of the center eigenvalues. Interpret both quantities.
\item \emph{Exploration.} Compare trajectories from $(0.55,0.25,0.20)$ at $s=0.7,1,1.3$. Record the product at several times and relate its behavior to the direction of spiraling.
\item \emph{Exploration.} Compare initial states $(0.4,0.6,0)$ and $(0.4,0.59,0.01)$ for $s=1$. Explain why their long-term paths differ.
\item If all $n$ fitnesses are equal at a state, prove the state is stationary. Prove the converse for an interior state. Explain how the converse changes on a face.
\item For an interior replicator trajectory, prove $\frac{d}{dt}\ln(p_i/p_j)=f_i-f_j$. Interpret the result for types with equal fitness along a trajectory.
\item \emph{Exploration.} Implement a small-step numerical integration for the balanced cyclic model. Compare the error in $\sum_i p_i$ with the error in $p_Ap_Bp_C$. Explain why preserving one identity numerically need not preserve the other.
\end{exercises}

\chapter{Probability Distributions as States}
\label{ch:probability}

A normalized population state becomes a probability distribution when one individual is sampled. The same mathematical structure describes any finite list of possible observations. This chapter develops probability, expectation, and sampling directly on the simplex, providing the statistical interpretation needed for information geometry.

There are two ways to use a probability vector. It may describe the current mixture in a changing population, or it may specify the chances of possible outcomes in a fixed observation model. Both interpretations assign the same coordinates to the same simplex. They differ in the questions we ask: how the state moves in the first setting, and what we expect to observe in the second.

\section{Outcomes, events, and probability laws}\index{probability law}
Let $\Omega=\{\omega_1,\ldots,\omega_n\}$ be a finite set of possible outcomes. A probability law assigns a number $p_i\geq0$ to each outcome, with $\sum_i p_i=1$. An \emph{event} is a subset $E\subseteq\Omega$, and its probability is the sum of the probabilities of its outcomes:
\[
\Pr(E)=\sum_{\omega_i\in E}p_i.
\]
This includes $\Pr(\Omega)=1$ and $\Pr(\varnothing)=0$. For disjoint events $E$ and $F$, their union has probability $\Pr(E)+\Pr(F)$, since the two sums have no outcomes in common.

A law $P=(p_1,\ldots,p_n)$ is a point of the simplex $\Delta^{n-1}$. It describes the probabilities before an observation is made. An observation is one realized element of $\Omega$. A sequence of observations can be summarized by another simplex point, the empirical distribution, developed below.

\begin{example}[A finite observation model]
A device returns one of three symbols A, B, and C, with probabilities $(1/2,1/3,1/6)$. Find the probability of a return other than A.
\begin{solution}
The event is $\{B,C\}$, whose probability is $1/3+1/6=1/2$. Equivalently it is the complement of $\{A\}$, so its probability is $1-1/2$.
\end{solution}
\end{example}

The outcome labels carry no numerical information by themselves. To compute a mean, one must specify the numbers assigned to them. Assigning measurement values, monetary returns, or indicator values produces different random variables on the same probability law.

\section{Random variables and expectation}\index{expectation}
A \emph{random variable}\index{random variable} on a finite outcome space is a function $Z:\Omega\to\R$. Write $z_i=Z(\omega_i)$. Its expectation is
\begin{equation}\label{eq:expectation}
\E_P[Z]=\sum_i p_i z_i.
\end{equation}
The subscript emphasizes which distribution supplies the weights. If the law changes while the numerical assignment stays fixed, the expectation changes as a function on the simplex.

\begin{example}[A weighted mean]
Let $(z_A,z_B,z_C)=(0,1,3)$ and $P=(1/2,1/3,1/6)$. Compute the expectation.
\begin{solution}
Equation \eqref{eq:expectation} gives
\[
\E_P[Z]=0\left(\frac12\right)+1\left(\frac13\right)+3\left(\frac16\right)=\frac56.
\]
An individual reading can be $0$, $1$, or $3$. The mean $5/6$ describes the average under repeated independent observation, so it need not be a possible individual reading.
\end{solution}
\end{example}

\begin{figure}[htbp]
\centering
\begin{tikzpicture}
\begin{axis}[width=.83\textwidth,height=5cm,ybar,bar width=19pt,xmin=-.55,xmax=3.55,ymin=0,ymax=.6,xtick={0,1,3},ytick={0,.1666667,.3333333,.5},yticklabels={0,$1/6$,$1/3$,$1/2$},xlabel={numerical reading $z$},ylabel={probability},grid=major]
\addplot[teal,fill=teal!55] coordinates {(0,.5)(1,.33333333)(3,.16666667)};
\end{axis}
\end{tikzpicture}
\caption{The law of a numerical random variable. The horizontal positions are its actual values, including the larger gap between 1 and 3; the bar heights are probabilities.}
\end{figure}

Expectation is linear. For constants $a,b$ and random variables $Z,W$ on the same space,
\[
\E[aZ+bW]=a\E[Z]+b\E[W].
\]
This follows by distributing the finite sum. No independence assumption is needed. In particular, a constant random variable $c$ has expectation $c$, and adding $c$ to every reading adds $c$ to its mean.

An indicator variable $\mathbf{1}_E$ equals 1 when event $E$ occurs and 0 otherwise. Its expectation is $\Pr(E)$. Thus event probabilities themselves are expectations of simple numerical assignments.

\section{Variance and covariance}
Writing $\mu=\E[Z]$, define the \emph{variance}\index{variance} by
\[
\Var_P(Z)=\E_P[(Z-\mu)^2].
\]
Expanding the square and using linearity gives
\begin{equation}\label{eq:variance}
\Var_P(Z)=\E_P[Z^2]-\bigl(\E_P[Z]\bigr)^2.
\end{equation}
Variance is nonnegative because it averages squares. It is zero exactly when all outcomes having positive probability receive the same numerical value. If the readings have physical units, variance has the square of those units; the standard deviation $\sqrt{\Var(Z)}$ has the original units.

\begin{example}[Computing a variance]
For the readings and probabilities of the previous example, calculate the variance.
\begin{solution}
The second moment is
\[
\E[Z^2]=0+\frac13+\frac{9}{6}=\frac{11}{6}.
\]
Thus $\Var(Z)=11/6-25/36=41/36$. The standard deviation is $\sqrt{41}/6$, approximately $1.067$.
\end{solution}
\end{example}

For two random variables, their \emph{covariance}\index{covariance} is
\begin{equation}\label{eq:covariance}
\Cov_P(Z,W)=\E_P[(Z-\E_PZ)(W-\E_PW)]
=\E_P[ZW]-\E_P[Z]\E_P[W].
\end{equation}
Positive covariance means that readings above their respective means tend to occur together in the weighted average. Covariance with itself is variance. Direct expansion gives
\[
\Var(aZ+b)=a^2\Var(Z),
\]
so shifting all readings leaves variance unchanged while rescaling them changes variance quadratically.

\section{The Bernoulli family}\index{Bernoulli family}
A Bernoulli variable $X$ has outcomes 0 and 1, with $\Pr(X=1)=p$. The single parameter $p\in[0,1]$ specifies its law. Since $X^2=X$ on these two possible values,
\begin{equation}\label{eq:bernoullimoments}
\E[X]=p,\qquad \Var(X)=p-p^2=p(1-p).
\end{equation}
Completing the square gives $p(1-p)=1/4-(p-1/2)^2$. Thus the variance is maximized at an equal split and vanishes at either certain outcome.

An indicator variable is Bernoulli with parameter $\Pr(E)$. This observation connects Bernoulli calculations with the sampling of any event, even in a much larger outcome space.

\begin{example}[What the probability parameter describes]
Suppose a device returns a positive signal with probability $0.2$. Let $X$ indicate a positive signal. Find its expectation and variance, then repeat after the signal is coded as $Y=5X$.
\begin{solution}
For $X$, the mean is $0.2$ and variance is $0.2(0.8)=0.16$. For the recoded reading $Y$, the mean is $1$ and variance is $25(0.16)=4$. The probability law of positive versus negative signals has the same parameter; the numerical readings use a different scale.
\end{solution}
\end{example}

\section{Functions and curves on the simplex}
With fixed numerical values $z_i$, the mean $\mu(\p)=\sum_i p_i z_i$ is affine in the probabilities. For three outcomes with readings $(0,1,3)$ and coordinates $(x,y,1-x-y)$,
\[
\mu(x,y)=3-3x-2y.
\]
Its level sets are line segments in the simplex. Every state on a given segment has the same expected reading, even though its three probabilities can differ.

\begin{figure}[htbp]
\centering
\begin{tikzpicture}[scale=4.3]
\fill[soft] (0,0)--(1,0)--(.5,.866025)--cycle;
\draw[ink] (0,0)--(1,0)--(.5,.866025)--cycle;
\draw[teal,thick] (.5,0)--(.083333,.144338);
\draw[teal,thick] (1,0)--(.166667,.288675);
\draw[copper,thick] (.75,.433013)--(.333333,.577350);
\node[below left] at (0,0){$A:0$};\node[below right] at (1,0){$B:1$};\node[above] at (.5,.866025){$C:3$};
\node[font=\small,above] at (.41,.205){$\mu=1$};
\node[font=\small,above] at (.55,.51){$\mu=2$};
\end{tikzpicture}
\caption{Expectation as an affine function on the probability triangle. Each line segment consists of laws with the same mean reading.}
\end{figure}

For a smooth curve $\p(t)$, differentiating normalization gives $\sum_i p_i'=0$. Differentiating expectation gives
\begin{equation}\label{eq:meanvelocity}
\frac{d}{dt}\E_{\p(t)}[Z]=\sum_i z_i p_i'(t).
\end{equation}
The numerical values are fixed in this formula; only their probabilities vary. Under replicator dynamics it becomes
\begin{equation}\label{eq:covarianceidentity}
\frac{d}{dt}\E_\p[Z]=\sum_i z_i p_i(f_i-\bar f)=\Cov_\p(Z,f).
\end{equation}
Here the fitness values at the current state form another random variable on the type labels. A reading increases in mean when it is positively associated with fitness.

\begin{example}[A mean changing under selection]
At $\p=(0.2,0.3,0.5)$, take readings $z=(0,1,3)$ and fitnesses $f=(3,1,2)$. Find the instantaneous mean change.
\begin{solution}
The replicator velocity is $(0.22,-0.27,0.05)$. Therefore $\mu'=0(0.22)+1(-0.27)+3(0.05)=-0.12$. Alternatively, $\E[Zf]=3.3$, $\E[Z]=1.8$, and $\E[f]=1.9$, so covariance is $3.3-1.8(1.9)=-0.12$.
\end{solution}
\end{example}

\section{Mixtures of probability laws}\index{mixture}
Given laws $P$ and $R$ on the same outcome set, their mixture is
\begin{equation}\label{eq:mixture}
M(t)=(1-t)P+tR,\qquad 0\leq t\leq1.
\end{equation}
Each coordinate is nonnegative, and their sum is one. This proves that the simplex is convex. Operationally, one can generate the mixture by selecting the first law with probability $1-t$ and the second with probability $t$, then drawing an observation from the selected law.

Its mean is $(1-t)\mu_P+t\mu_R$. To find its variance, use the corresponding mixture of second moments and subtract the squared mixture mean. The result is
\begin{equation}\label{eq:mixturevariance}
\Var_{M(t)}(Z)=(1-t)\Var_P(Z)+t\Var_R(Z)
+t(1-t)(\mu_P-\mu_R)^2.
\end{equation}
The first two terms average the variances within the two component laws. The final term accounts for the separation of their means. A mixture of two concentrated laws can have substantial variance when their centers are far apart.

\begin{example}[A straight distribution path]
Let $P=(0.6,0.3,0.1)$, $R=(0.2,0.3,0.5)$, and $z=(0,1,3)$. Find the mean and variance along their mixture.
\begin{solution}
The path is $(0.6-0.4t,0.3,0.1+0.4t)$, with tangent velocity $(-0.4,0,0.4)$. Its mean is $0.6+1.2t$, and its second moment is $1.2+3.6t$. Consequently
\[
\Var_{M(t)}(Z)=0.84+2.16t-1.44t^2.
\]
The mean is affine along this path, while the variance is a quadratic function. At the endpoints the variances are $0.84$ and $1.56$.
\end{solution}
\end{example}

\begin{figure}[htbp]
\centering
\begin{tikzpicture}
\begin{axis}[width=.85\textwidth,height=5cm,xmin=0,xmax=1,ymin=0,ymax=2,xlabel={mixture parameter $t$},ylabel={moment value},grid=major,legend pos=south east]
\addplot[teal,domain=0:1]{.6+1.2*x};\addlegendentry{mean}
\addplot[copper,domain=0:1]{.84+2.16*x-1.44*x^2};\addlegendentry{variance}
\end{axis}
\end{tikzpicture}
\caption{Two scalar functions along the same mixture path. Their different dependence on $t$ follows from their definitions as functions of the probabilities.}
\end{figure}

\section{Samples and empirical distributions}
An abstract probability vector assigns chances to possible outcomes. A sample gives us a second probability vector obtained from observed relative frequencies. These two points of the simplex generally differ, and the difference is random. The first calculation measures the size of this sampling variation in each coordinate.
Draw $N$ observations independently from a fixed law $\p$. Independence means that the joint probability of a specified sequence is the product of the individual outcome probabilities. Let $N_i$ count the occurrences of outcome $i$. The \emph{empirical distribution}\index{empirical distribution} is
\[
\widehat p_i=\frac{N_i}{N}.
\]
For example, counts $(27,21,12)$ in 60 observations give $(0.45,0.35,0.20)$. Sampling from the same law again generally gives a different point.

Write $I_{ri}$ for the indicator that observation $r$ is outcome $i$. Then $\widehat p_i=N^{-1}\sum_{r=1}^N I_{ri}$. Each indicator has mean $p_i$ and variance $p_i(1-p_i)$. Independence between observations makes their cross-covariances zero, giving
\begin{equation}\label{eq:empiricalmoments}
\E[\widehat p_i]=p_i,\qquad
\Var(\widehat p_i)=\frac{p_i(1-p_i)}{N}.
\end{equation}
Thus the empirical proportion is centered at the model probability, and its standard deviation decreases like $N^{-1/2}$. Multiplying the sample size by four halves that standard deviation.

We can turn the variance calculation into a guarantee about large sampling errors. A general inequality bounds the chance of deviating from a mean by at least $\varepsilon$: any random variable $Y$ with mean $m$ satisfies
\[
\Pr(|Y-m|\geq\varepsilon)\leq\frac{\Var(Y)}{\varepsilon^2}.
\]
Indeed, $(Y-m)^2\geq\varepsilon^2$ on the indicated event, so averaging the square gives at least $\varepsilon^2$ times the event's probability. Applying this inequality to $\widehat p_i$ proves convergence in probability to $p_i$ as $N\to\infty$, a finite-outcome version of the \emph{law of large numbers}\index{law of large numbers}.

The normalized sampling errors satisfy $\sum_i(\widehat p_i-p_i)=0$. Sampling therefore produces random displacements within the same affine plane that contains the dynamical tangent vectors. Information geometry will supply a location-dependent way to measure those displacements.

\section{The shape of sampling fluctuations}\index{sampling covariance}
The empirical probabilities fluctuate together because their sum is always one. To describe this dependence, let $i\ne j$ be two categories and consider their indicators in one observation. They cannot both equal 1, so
\[
\E[I_iI_j]=0,\qquad \Cov(I_i,I_j)=-p_ip_j.
\]
Indicators belonging to different independent observations have covariance zero. Averaging over the $N$ observations therefore gives
\[
\Cov(\widehat p_i,\widehat p_j)=-\frac{p_ip_j}{N}\qquad(i\ne j).
\]
Together with the coordinate variances, this yields the full covariance matrix
\begin{equation}\label{eq:empiricalcovariance}
\operatorname{Cov}(\widehat\p)=\frac1N\bigl(\diag(\p)-\p\p^T\bigr).
\end{equation}
Multiplying this matrix by $\ones$ gives zero, expressing the absence of fluctuation in the total probability. The random cloud of empirical states spreads only within the normalization plane.

For three outcomes, retaining only the first two independent coordinates gives
\[
\operatorname{Cov}\begin{pmatrix}\widehat p_A\\\widehat p_B\end{pmatrix}
=\frac1N\begin{pmatrix}p_A(1-p_A)&-p_Ap_B\\-p_Ap_B&p_B(1-p_B)\end{pmatrix}.
\]
Its off-diagonal entries are negative: an unusually large observed share of A uses some of the fixed total share, tending to accompany smaller shares elsewhere. The effect follows from normalization rather than from a physical interaction among the categories.

\begin{example}[Uncertainty in two empirical coordinates]
For $\p=(0.5,0.3,0.2)$ and $N=100$, find the covariance matrix for the first two empirical probabilities.
\begin{solution}
The matrix is
\[
\begin{pmatrix}0.0025&-0.0015\\-0.0015&0.0021\end{pmatrix}.
\]
The first coordinate's standard deviation is $0.05$, and the second's is $\sqrt{0.0021}\approx0.04583$. Their covariance is negative. Repeating independent samples produces a cloud centered at $(0.5,0.3)$ whose spread depends on direction.
\end{solution}
\end{example}

Finite samples can also miss a possible category entirely. For a category of probability $p_i$, the probability of no occurrence in $N$ independent observations is $(1-p_i)^N$: every observation must avoid that category.

\begin{example}[A possible category absent from a sample]
A category has model probability $0.02$. Find the chance that it is absent from a sample of 25 observations.
\begin{solution}
The probability is $0.98^{25}\approx0.6035$. Thus a zero empirical coordinate is quite plausible for this sample size even though the model coordinate is positive. An empirical distribution reports what occurred in the sample; the probability model describes the process that produced it.
\end{solution}
\end{example}

The probability of missing a fixed positive-probability category tends to zero as the sample grows. At each finite sample size, however, the empirical distribution can lie on a face of the simplex while the generating law lies in its interior. Support conventions in information calculations must distinguish these two laws carefully.

\FloatBarrier
\section{Exercises}
\begin{exercises}
\item For $P=(1/4,1/2,1/4)$ and readings $(0,1,3)$, compute the mean, second moment, and variance.
\item Prove linearity of expectation directly from the finite sum. Explain why independence is unnecessary.
\item If $Z$ takes values $-2,1,4$ with probabilities $1/3,1/2,1/6$, compute its mean and variance.
\item Prove $\Var(aZ+b)=a^2\Var(Z)$.
\item Derive the Bernoulli mean and variance from their definitions. Find the maximum variance and interpret both endpoints.
\item Find two different laws on readings $(0,1,2)$ with the same mean. Compare their variances.
\item For the readings $(0,1,3)$, find the directional derivative of the mean in tangent direction $v=(-0.1,0.05,0.05)$. Interpret the probability transfers.
\item Derive \eqref{eq:mixturevariance} by expanding the second moments and means.
\item For the mixture in the worked example, find the parameter at which variance is largest and calculate that variance.
\item Prove the replicator covariance identity \eqref{eq:covarianceidentity}. What happens when the reading $z_i$ is a fixed constant fitness $f_i$?
\item A sample of 80 observations gives counts $(38,27,15)$. Find the empirical distribution. Compare its mean reading for $(0,1,3)$ with that of $(1/2,1/3,1/6)$.
\item Derive the expectation and variance of an empirical proportion using independent indicator variables.
\item For a Bernoulli model with $p=0.4$, calculate the standard deviations of the empirical proportion for $N=100$ and $N=400$.
\item Show that for distinct categories $i,j$, $\Cov(\widehat p_i,\widehat p_j)=-p_ip_j/N$. Explain why these empirical coordinates are dependent.
\item \emph{Exploration.} Simulate samples of sizes 10, 100, and 1000 from $(1/2,1/3,1/6)$. Compare empirical simplex points and empirical mean readings across several independent runs.
\item For four readings $(0,1,2,3)$, verify that the distributions $(1/4,1/4,1/4,1/4)$ and $(5/24,9/24,3/24,7/24)$ have the same mean and variance. What does this show about moment summaries?
\item Use the variance bound above to give a sufficient sample size for $\Pr(|\widehat p-p|\geq0.05)\leq0.1$, uniformly over all Bernoulli parameters $p$.
\item \emph{Exploration.} Trace a straight mixture between two categorical laws. Plot mean and variance against the mixture parameter and verify their predicted linear and quadratic forms.
\end{exercises}

\chapter{Entropy, Log-Odds, and Exponential Updating}
\label{ch:entropy}

Logarithms translate products of probabilities into sums. This property supports a measure of information carried by observations, and it also simplifies the relative growth of competing types. Shannon entropy and logarithmic probability ratios bring these two uses of the logarithm into the same setting.

Consider two related tasks: measuring the information provided by an observation and comparing how rapidly two population types change in proportion. The first involves the logarithm of an event probability; the second involves the logarithm of a ratio of probabilities. The logarithm makes both calculations additive, even when the underlying probabilities or abundances evolve multiplicatively.

\section{Surprisal and information units}\index{nat}\index{bit}
For an event of probability $p>0$, its \emph{surprisal}\index{surprisal} is
\begin{equation}\label{eq:surprisal}
I(p)=-\ln p.
\end{equation}
A certain event has zero surprisal. Less probable events have larger surprisal, and $I(p)\to\infty$ as $p\to0^+$. The value is determined by the probability assigned before the observation.

If events E and F are independent, $\Pr(E\cap F)=\Pr(E)\Pr(F)$. Therefore
\[
I(E\cap F)=I(E)+I(F).
\]
Independent observations accumulate information additively. With continuity and monotonicity, this additivity under multiplication characterizes a positive multiple of the negative logarithm. The multiplicative constant selects the information unit.

Natural logarithms give \emph{nats}. Base-2 logarithms give \emph{bits}. Since $\log_2p=\ln p/\ln2$, a quantity in nats is converted to bits by dividing by $\ln2$.

\begin{example}[Information from joint events]
An event has probability $1/2$ and an independent event has probability $1/5$. Find their joint surprisal.
\begin{solution}
The joint probability is $1/10$. Its surprisal is $\ln10\approx2.3026$ nats, or $\log_2 10\approx3.3219$ bits. It is also $\ln2+\ln5$, the sum of the individual surprisals.
\end{solution}
\end{example}

For dependent observations the joint probability instead factors into a first probability and a conditional probability. The same logarithmic algebra then adds the information of the first event to the information of the second given the first. Independence is the condition that allows the unconditional probabilities to be used in both terms.

\section{Shannon entropy as expected surprisal}
For a law $P=(p_i)$, average the surprisal of each possible observation using that same law:
\begin{equation}\label{eq:entropy}
\boxed{H(P)=-\sum_i p_i\ln p_i.}
\end{equation}
This is \emph{Shannon entropy}\index{entropy!Shannon}. The convention $0\ln0=0$ follows from $\lim_{p\to0^+}p\ln p=0$, making $H$ continuous on the closed simplex.

Entropy quantifies the expected information of the next observation under the specified law. A deterministic law has entropy zero. For a uniform law on $n$ outcomes, every observation has surprisal $\ln n$, so the entropy is $\ln n$ nats.

\begin{example}[Three-outcome entropy]
Calculate the entropy of $P=(1/2,1/3,1/6)$.
\begin{solution}
The expected surprisal is
\[
H(P)=\frac12\ln2+\frac13\ln3+\frac16\ln6
\approx1.0114\text{ nats}.
\]
Dividing by $\ln2$ gives approximately $1.4591$ bits. The entropy of the uniform three-outcome law is $\ln3\approx1.0986$ nats.
\end{solution}
\end{example}

Entropy depends on the category probabilities. Recoding category labels as numerical readings leaves this value unchanged. Means and variances depend on the readings as well as the probabilities, so they answer a different statistical question.

For two independent finite random variables with laws $(p_i)$ and $(r_j)$, the joint law is $(p_ir_j)$. Expanding its entropy gives
\begin{align*}
H(X,Y)&=-\sum_{i,j}p_ir_j\ln(p_ir_j)\\
&=-\sum_i p_i\ln p_i\sum_jr_j-\sum_jr_j\ln r_j\sum_ip_i\\
&=H(X)+H(Y).
\end{align*}
This is the expected-information counterpart of surprisal additivity.

\section{Binary entropy and the uniform maximum}\index{maximum entropy}\index{concavity}
For a two-outcome law,
\[
H(p)=-p\ln p-(1-p)\ln(1-p).
\]
Differentiating in the interior gives
\begin{equation}\label{eq:binaryentropyderivatives}
H'(p)=\ln\frac{1-p}{p},\qquad
H''(p)=-\frac1{p(1-p)}.
\end{equation}
The second derivative is negative, so the graph is strictly concave. Its unique interior critical point is $p=1/2$, where the entropy is $\ln2$. Symmetry under $p\mapsto1-p$ corresponds to interchanging the labels of the two outcomes.

\begin{figure}[htbp]
\centering
\begin{tikzpicture}
\begin{axis}[width=.88\textwidth,height=5cm,xmin=0,xmax=1,ymin=0,ymax=.77,xlabel={probability $p$},ylabel={$H(p)$ in nats},xtick={0,.25,.5,.75,1},grid=major]
\addplot[teal,domain=.0001:.9999,samples=150]{-x*ln(x)-(1-x)*ln(1-x)};
\addplot[only marks,ink] coordinates {(0,0)(1,0)(.5,.69314718)};
\end{axis}
\end{tikzpicture}
\caption{Binary entropy. Its maximum is $\ln2$ at equal probabilities, and its continuous extension vanishes at the certain outcomes.}
\end{figure}

The corresponding multivariable calculation is equally explicit. On positive probabilities, the ambient Hessian of $H$ is $-\diag(1/p_i)$. For every nonzero tangent vector $v$,
\[
v^T(\operatorname{Hess}H)v=-\sum_i\frac{v_i^2}{p_i}<0.
\]
Thus entropy is strictly concave on the positive simplex. A constrained critical point satisfies $-(\ln p_i+1)=\lambda$ for all $i$, forcing $p_i=1/n$.

\begin{proposition}[Maximum entropy on a finite outcome space]
For $P\in\Delta^{n-1}$,
\[
0\leq H(P)\leq\ln n.
\]
The upper bound is attained only at the uniform law. The lower bound is attained exactly at pure laws.
\end{proposition}
\begin{proof}
Each summand $-p_i\ln p_i$ is nonnegative and is positive for $0<p_i<1$, giving the lower statement. Continuity on the compact simplex guarantees a maximum. At a maximizer, let $k$ be the number of positive coordinates. Lagrange multipliers within that support force each to equal $1/k$, giving entropy $\ln k\leq\ln n$. The uniform $n$-outcome law reaches $\ln n$, with equality requiring $k=n$ and equal coordinates.
\end{proof}

Entropy can also be maximized subject to a prescribed expectation. Such constraints retain information about the numerical readings. They lead to exponential weights, as developed after the dynamical derivation below.

\section{Odds and log-odds coordinates}
For an interior two-type state $x$, the odds of A relative to B are $x/(1-x)$. Define the \emph{log-odds}\index{log-odds} coordinate
\[
L=\ln\frac{x}{1-x}.
\]
It maps $(0,1)$ bijectively to $\R$, with inverse
\[
x=\frac{e^L}{1+e^L}=\frac1{1+e^{-L}}.
\]
The midpoint $x=1/2$ corresponds to $L=0$; probabilities approaching 0 or 1 correspond to $L\to-\infty$ or $L\to+\infty$.

Along a replicator solution,
\begin{equation}\label{eq:logoddsevolution}
L'=\left(\frac1x+\frac1{1-x}\right)x'
=\frac{x'}{x(1-x)}=f_A-f_B.
\end{equation}
The derivative of the log-odds is exactly the fitness advantage. A constant advantage gives $L(t)=L(0)+kt$, and applying the inverse transformation recovers the logistic solution.

\begin{example}[A nonlinear curve in a new coordinate]
For $x'=2x(1-x)$ with $x(0)=1/4$, determine $L(t)$ and the time when $x=3/4$.
\begin{solution}
Initially $L(0)=-\ln3$, so $L(t)=-\ln3+2t$. The target has log-odds $\ln3$. Thus $2t=2\ln3$ and the target time is $t=\ln3$.
\end{solution}
\end{example}

\begin{figure}[htbp]
\centering
\begin{tikzpicture}
\begin{axis}[width=.46\textwidth,height=4.6cm,xmin=0,xmax=3,ymin=0,ymax=1,xlabel={$t$},ylabel={$x(t)$},grid=major]
\addplot[teal,domain=0:3,samples=70]{1/(1+3*exp(-2*x))};
\end{axis}
\begin{axis}[at={(.51\textwidth,0)},anchor=south west,width=.46\textwidth,height=4.6cm,xmin=0,xmax=3,ymin=-1.2,ymax=5,xlabel={$t$},ylabel={$L(t)$},grid=major]
\addplot[copper,domain=0:3]{-ln(3)+2*x};
\end{axis}
\end{tikzpicture}
\caption{The same constant-advantage solution in proportion and log-odds coordinates. The linear log-odds curve gives an exact solution method.}
\end{figure}

For general state-dependent fitnesses, the identity $L'=\Delta(x(t))$ still holds. Its right-hand side then varies along the solution, so the coordinate change produces a useful equation without necessarily giving a straight time graph.

\section{Exponential normalization for constant fitness}\index{exponential normalization}
Suppose $n$ types have constant fitnesses $c_i$. Their absolute growth weights are
\[
w_i(t)=p_i(0)e^{c_it}.
\]
Normalize by $Z(t)=\sum_jw_j(t)$ to obtain
\begin{equation}\label{eq:expnormalization}
\boxed{p_i(t)=\frac{p_i(0)e^{c_it}}{\sum_jp_j(0)e^{c_jt}}.}
\end{equation}
Differentiation verifies the equation:
\[
p_i'=p_i\left(c_i-\frac{Z'}Z\right),\qquad
\frac{Z'}Z=\sum_jp_jc_j=\bar c.
\]
Thus these normalized weights solve the constant-fitness replicator system exactly.

\begin{example}[An exact three-type update]
Let $\p(0)=(1/5,3/10,1/2)$ and $c=(1,0,-1)$. Find the state at $t=\ln2$.
\begin{solution}
The weights are $(2/5,3/10,1/4)$, with sum $19/20$. After normalization,
\[
\p(\ln2)=\left(\frac8{19},\frac6{19},\frac5{19}\right).
\]
The first type gains share; the third loses share. The middle type has unchanged unnormalized weight, but its normalized proportion also changes.
\end{solution}
\end{example}

Taking ratios eliminates the normalizing sum:
\[
\ln\frac{p_i(t)}{p_j(t)}=\ln\frac{p_i(0)}{p_j(0)}+(c_i-c_j)t.
\]
Adding the same constant to all fitnesses multiplies all weights by the same exponential factor and leaves the normalized state unchanged.

The \emph{softmax}\index{softmax} of a score vector $s$ is the normalized vector $e^{s_i}/\sum_je^{s_j}$. For positive initial probabilities, \eqref{eq:expnormalization} is the softmax of $s_i=\ln p_i(0)+c_it$. To compute it reliably for large scores, subtract $m=\max_i s_i$ before exponentiating. The common factor $e^{-m}$ cancels, while every exponent $s_i-m$ is nonpositive.

If the maximal fitness occurs only at type $k$ and $p_k(0)>0$, dividing the formula by $e^{c_kt}$ shows $p_k(t)\to1$. If several initially present types tie for maximal fitness, their final proportions retain their initial ratios after the other types vanish asymptotically.

\section{The log-normalizer and constrained entropy}\index{log-normalizer}
The denominator of an exponential probability update looks at first like a bookkeeping device: it keeps all the probabilities summing to one. Its logarithm has a more substantial role. Its first derivative records mean fitness, while its second derivative records the variability of fitness. This connects the changing distribution to the convexity of a scalar function.
For constant fitnesses define $\psi(t)=\ln Z(t)$. Differentiating the normalizer gives
\begin{equation}\label{eq:logpartition}
\psi'(t)=\bar c(t),\qquad
\psi''(t)=\Var_{\p(t)}(c)\geq0.
\end{equation}
To verify the second identity, differentiate $\sum_i p_ic_i$ and substitute $p_i'=p_i(c_i-\bar c)$. The result is $\sum_i p_ic_i^2-\bar c^2$. Thus the log-normalizer is convex, and mean fitness increases under constant-fitness selection.

Exponential normalization also appears when maximizing entropy subject to a prescribed mean $\sum_i p_i z_i=m$. At an interior constrained critical point, Lagrange multipliers give
\[
-(\ln p_i+1)+\lambda z_i+\eta=0.
\]
Solving and normalizing produces
\begin{equation}\label{eq:maxent}
p_i(\lambda)=\frac{e^{\lambda z_i}}{\sum_je^{\lambda z_j}}.
\end{equation}
The mean constraint determines $\lambda$. Strict concavity of entropy makes any such feasible interior critical point the unique maximum. If the readings are not all equal, differentiation shows $dm/d\lambda=\Var_{p(\lambda)}(Z)>0$, so the mean varies strictly with $\lambda$.

\begin{example}[Maximum entropy with a mean constraint]
For readings $(0,1,2)$, find the maximum-entropy law with mean $10/7$.
\begin{solution}
Take $\lambda=\ln2$. Equation \eqref{eq:maxent} gives weights $(1,2,4)$ and law $(1/7,2/7,4/7)$. Its mean is $2/7+8/7=10/7$, so it satisfies the constraint and is the unique entropy maximum.
\end{solution}
\end{example}

\section{Reading an exponentially normalized trajectory}
The exact solution lets each coordinate's behavior be understood through its relation to the changing mean fitness. A fixed fitness value does not require a monotone proportion. The population mean itself moves as the weights of the other types change.

\begin{example}[A middle type that first gains and then loses share]
Take constant fitnesses $(1,0,-1)$ and initial law $(0.2,0.3,0.5)$. Determine when the second type's proportion is largest.
\begin{solution}
The normalizer and second coordinate are
\[
Z(t)=0.2e^t+0.3+0.5e^{-t},\qquad p_B(t)=\frac{0.3}{Z(t)}.
\]
The denominator decreases while $Z'<0$ and increases while $Z'>0$. Its only stationary point satisfies $0.2e^t=0.5e^{-t}$, giving
\[
t_* =\frac12\ln\frac52.
\]
At that time the A and C weights are both $\sqrt{0.1}$, and
\[
p_B(t_*)=\frac{0.3}{0.3+2\sqrt{0.1}}\approx0.3217.
\]
Initially the mean fitness is $-0.3$, below B's fitness 0, so B gains share. At $t_*$ the mean reaches zero. Thereafter the mean is positive and B loses share. Its final proportion is zero as the uniquely fittest A type dominates.
\end{solution}
\end{example}

The convexity of $\ln Z$ supplies another explanation of the unique turning point. The mean fitness $\psi'$ increases, and in this example it crosses B's fixed fitness once. More generally, a constant-fitness type gains share exactly while its fitness is above the current mean. Logarithmic ratios between any two types still have constant derivatives throughout this change in a single coordinate's direction.

The transformation to log-odds also clarifies binary entropy near equal probabilities. Since $H'(p)=-L(p)$, an infinitesimal transfer toward the less probable outcome increases entropy. At $p=1/2$ the first derivative vanishes; the leading effect of a small imbalance is quadratic. Taylor expansion gives
\[
H\left(\frac12+h\right)=\ln2-2h^2+O(h^4).
\]
The cubic term vanishes by symmetry. For instance, a displacement $h=0.05$ gives the approximation $H(0.55)\approx\ln2-0.005$. Entropy has a horizontal tangent at its maximum, so the information change is smaller than first order in this small probability displacement.

For $n$ outcomes, write a perturbation of uniformity as $p_i=1/n+h_i$ with $\sum_i h_i=0$. The constant gradient components cancel against normalization, and the entropy Hessian gives
\[
H(U+h)=\ln n-\frac n2\sum_i h_i^2+O(\|h\|^3).
\]
This measures the entropy deficit near uniformity. Away from uniformity, the reciprocal-probability coefficients vary with the reference state, a dependence developed explicitly through relative entropy in the next chapter.

\section{Entropy along a dynamical trajectory}
Entropy is a function on the simplex, so a population trajectory produces a time-dependent entropy simply by evaluating that function along the path. The direction of motion alone does not determine whether entropy increases: its rate also depends on how the current probability weights are distributed. Differentiation makes this dependence explicit.
For an interior probability curve, differentiation of entropy gives
\[
\dot H=-\sum_i(\ln p_i+1)p_i'=-\sum_i\ln p_i\,p_i',
\]
since $\sum_i p_i'=0$. Under replicator dynamics,
\begin{equation}\label{eq:entropyrate}
\dot H=-\Cov_\p(\ln p,f).
\end{equation}
Its sign depends on the current law and fitness field. In a two-type model the formula is
\[
\dot H=\ln\frac{1-x}{x}\,x(1-x)\Delta(x).
\]

\begin{example}[Entropy during stable coexistence]
Follow $x'=x(1-x)(3-4x)$ from $x_0=0.1$.
\begin{solution}
The state increases toward $3/4$. Until it reaches $1/2$, both $x'$ and $H'(x)$ are positive, so entropy increases. At $1/2$, entropy reaches $\ln2$. For the rest of the trajectory, $x'>0$ and $H'(x)<0$, so entropy decreases toward $H(3/4)$. The dynamics approach a stable equilibrium throughout this change of entropy behavior.
\end{solution}
\end{example}

Entropy measures the uncertainty of observations at the current state. To measure mismatch with a specified equilibrium or predictive model, the next chapter introduces a function of two distributions.

\FloatBarrier
\section{Exercises}
\begin{exercises}
\item Compute surprisals for probabilities $1/2,1/4,1/10$ in nats and bits.
\item Prove entropy additivity for two independent finite laws by expanding the joint sum.
\item Calculate $H(1/2,1/4,1/4)$ in nats and bits. Compare with the uniform three-outcome entropy.
\item Derive both formulas in \eqref{eq:binaryentropyderivatives}. Explain the endpoint values by continuity.
\item Show that a uniform law on $k$ of $n$ available outcomes has entropy $\ln k$. Identify all minimum-entropy laws.
\item Prove strict concavity of entropy along every nonconstant straight segment of positive distributions.
\item For $x'=1.2x(1-x)$ and $x_0=1/4$, find $L(t)$ and the first time $x=1/2$.
\item If $x'=x(1-x)(3-4x)$, express the log-odds equation entirely in terms of $L$ using $x=(1+e^{-L})^{-1}$.
\item For $c=(1,0,-1)$ and $\p(0)=(0.2,0.3,0.5)$, calculate the state at $t=1$ to four decimal places.
\item Prove the softmax is unchanged by adding the same constant to every score. Use this to compute the softmax of $(1000,1001,999)$.
\item Derive both identities in \eqref{eq:logpartition} and explain why mean constant fitness is nondecreasing.
\item Suppose $c=(2,2,0)$ and $\p(0)=(0.2,0.3,0.5)$. Find the limiting distribution.
\item For readings $(0,1,2)$, derive the exponential form of a positive entropy maximizer subject to a mean constraint.
\item Prove that $m'(\lambda)=\Var(Z)$ for the law in \eqref{eq:maxent}. What happens when all readings agree?
\item For constant advantage $k>0$ and an initial state below $1/2$, determine where entropy increases and where it decreases. Find its long-term limit.
\item Derive \eqref{eq:entropyrate} for general replicator fitnesses. State precisely which normalization identity removes the extra term.
\item \emph{Exploration.} Plot a constant-fitness three-type trajectory and its entropy. Choose an initial law whose entropy first rises and later falls. Use the limiting law to explain the eventual behavior.
\item \emph{Exploration.} Compare proportion and log-odds time graphs for a constant advantage and a frequency-dependent advantage. Identify which features are consequences of the coordinate identity and which require constant fitness.
\end{exercises}

\chapter{Relative Entropy and Lyapunov Stability}
\label{ch:kl}

Predicting an observation requires a probability law. If the predictive law differs from the generating law, logarithmic scoring measures the resulting average penalty. The excess penalty is relative entropy, a comparison that also supplies useful functions for studying dynamical stability.

The direction of this comparison matters. Observations have one distribution, while predictions use another; reversing their roles changes the expected penalty. After developing this statistical meaning, we will fix one distribution as a reference and let the other move. Relative entropy then becomes a function along a dynamical trajectory, and its rate of change reveals stability.

\section{Prediction and expected logarithmic loss}\index{logarithmic loss}
Suppose observations are generated by $Q=(q_i)$ and predicted using $P=(p_i)$. If outcome $i$ occurs, the logarithmic loss assigned to the prediction is $-\ln p_i$. Averaging over the actual outcome law gives the \emph{cross-entropy}\index{cross-entropy}
\[
\operatorname{CE}(Q,P)=-\sum_iq_i\ln p_i.
\]
Using $Q$ itself for prediction gives expected loss $H(Q)$. The difference is
\begin{equation}\label{eq:kl}
\boxed{D(Q\|P)=\operatorname{CE}(Q,P)-H(Q)
=\sum_iq_i\ln\frac{q_i}{p_i}.}
\end{equation}
It is called \emph{relative entropy} or \emph{Kullback--Leibler divergence}\index{relative entropy|see{KL divergence}}\index{KL divergence}. Its first argument supplies the averaging distribution. This convention determines which expected prediction loss is being calculated.

\begin{example}[The cost of using an equal-probability model]
Let $Q=(3/4,1/4)$ and $P=(1/2,1/2)$. Compute the entropy, cross-entropy, and divergence.
\begin{solution}
Every outcome receives loss $\ln2$ under $P$, so $\operatorname{CE}(Q,P)=\ln2\approx0.693147$. The entropy is
\[
H(Q)=-\frac34\ln\frac34-\frac14\ln\frac14\approx0.562335.
\]
Their difference is
\[
D(Q\|P)=\frac34\ln\frac32+\frac14\ln\frac12\approx0.130812.
\]
The equal-probability model incurs about $0.1308$ extra nats of expected log loss per observation.
\end{solution}
\end{example}

For $N$ independent observations, the model probability of a sequence is the product of its predicted probabilities. Negative logarithms turn this product into a sum of individual losses. The expected excess total loss is therefore $N D(Q\|P)$. This gives divergence a direct interpretation at different sample sizes.

\section{Nonnegativity and the role of support}
The elementary inequality
\[
\ln u\leq u-1\qquad(u>0)
\]
has equality exactly at $u=1$. It follows by differentiating $u-1-\ln u$, whose unique minimum is zero at 1.

\begin{theorem}[Gibbs' inequality]\index{Gibbs' inequality}
For finite probability distributions, $D(Q\|P)\geq0$, with equality exactly when $P=Q$.
\end{theorem}
\begin{proof}
First assume every coordinate is positive. Apply the logarithm inequality to $u=p_i/q_i$, multiply by $q_i$, and sum:
\[
D(Q\|P)\geq\sum_iq_i\left(1-\frac{p_i}{q_i}\right)
=\sum_i(q_i-p_i)=0.
\]
Equality forces $p_i/q_i=1$ for every $i$.

At the boundary, terms with $q_i=0$ are assigned value zero. If $q_i>0$ but $p_i=0$, the divergence is infinite and the inequality holds. Otherwise apply the same calculation over the support $S$ of $Q$ to obtain $D(Q\|P)\geq1-\sum_{i\in S}p_i\geq0$. Equality requires $p_i=q_i$ on $S$ and no mass outside $S$, which again means $P=Q$.
\end{proof}

A predictive model assigning zero probability to a possible observation has infinite log loss when that observation occurs. The divergence becomes infinite whenever this happens with positive generating probability. Its finiteness condition is $\operatorname{supp}Q\subseteq\operatorname{supp}P$.

\begin{example}[Boundary direction matters]
Compare $Q=(1,0)$ and $P=(1/2,1/2)$ in both directions.
\begin{solution}
$D(Q\|P)=\ln2$ is finite: the single outcome generated by $Q$ receives positive probability under $P$. In the reverse direction, $Q$ assigns zero predictive probability to an outcome generated with probability $1/2$ under $P$, giving $D(P\|Q)=\infty$.
\end{solution}
\end{example}

\section{Divergence as a scalar function on the simplex}\index{convexity}
Fix a positive law $Q$. Then $P\mapsto D(Q\|P)$ is a smooth function on the positive simplex, with unique minimum zero at $P=Q$. Its ambient first and second derivatives are
\[
\frac{\partial D}{\partial p_i}=-\frac{q_i}{p_i},\qquad
\frac{\partial^2D}{\partial p_i\partial p_j}
=\begin{cases}q_i/p_i^2&i=j,\\0&i\ne j.\end{cases}
\]
The Hessian is positive definite, so the function is strictly convex. It diverges toward every boundary face when $Q$ has full support.

\begin{example}[Reversing two interior laws]
Reverse the laws in the first worked example.
\begin{solution}
\[
D(P\|Q)=\frac12\ln\frac{2}{3}+\frac12\ln2
=\frac12\ln\frac43\approx0.143841.
\]
The two values use different averaging laws. Both are nonnegative and vanish at equality, while their difference records the direction of prediction.
\end{solution}
\end{example}

KL divergence is generally asymmetric. It also fails the triangle inequality. For example, take $A=(1,0)$, $B=(1/2,1/2)$, and $C=(1/4,3/4)$. Then
\[
D(A\|C)=\ln4>\ln2+\frac12\ln\frac43
=D(A\|B)+D(B\|C).
\]
Thus divergence provides a prediction comparison rather than a metric distance. A metric will emerge from its local second-order behavior.

\begin{figure}[htbp]
\centering
\begin{tikzpicture}
\begin{axis}[width=.88\textwidth,height=5.3cm,xmin=.03,xmax=.97,ymin=0,ymax=2.3,xlabel={current probability $x$},ylabel={divergence in nats},grid=major,legend pos=north east]
\addplot[teal,domain=.03:.97,samples=150]{.75*ln(.75/x)+.25*ln(.25/(1-x))};\addlegendentry{$D(3/4\|x)$}
\addplot[copper,domain=.03:.97,samples=150]{x*ln(x/.75)+(1-x)*ln((1-x)/.25)};\addlegendentry{$D(x\|3/4)$}
\addplot[only marks,ink] coordinates {(.75,0)};
\end{axis}
\end{tikzpicture}
\caption{Two directed comparisons with the same reference probability. They share the minimum at equality and have different global shapes.}
\end{figure}

Along the mixture $M(t)=(1-t)P+tQ$, set $V(t)=D(Q\|M(t))$. Its second derivative is nonnegative, while $V'(1)=-\sum_i(q_i-p_i)=0$. Hence $V'(t)\leq0$ on $[0,1]$, strictly below zero for $t<1$ when $P\ne Q$. Moving along this mixture toward the reference reduces expected excess log loss.

\section{Fitting a model to observed frequencies}
Relative entropy also appears when a model is fitted to a finite data set. Suppose $N$ independent observations have category counts $N_i$, and set $\widehat q_i=N_i/N$. The probability assigned to the observed ordered sequence by a proposed law $P$ is
\[
L(P)=\prod_i p_i^{N_i}.
\]
As a function of the proposed law with the observed sequence fixed, this expression is called its \emph{likelihood}\index{likelihood}. The average negative log likelihood is
\[
-\frac1N\ln L(P)=-\sum_i\widehat q_i\ln p_i
=H(\widehat Q)+D(\widehat Q\|P).
\]
The empirical entropy does not depend on the proposed model. Maximizing likelihood therefore amounts to minimizing divergence from the empirical law to the model law. The direction of the divergence is determined by the observed frequencies that weight the logarithmic terms.

Among all laws on the closed simplex, Gibbs' inequality shows that the unique maximizing law is $P=\widehat Q$. When some observed counts are zero, this maximum can lie on the boundary. If the proposed laws are restricted to a smaller family, the fitted law minimizes the same directed divergence within that family. The data need not themselves lie in the chosen family.

\begin{example}[An empirical law and a uniform model]
A sample of 40 observations gives counts $(20,12,8)$. Compare the likelihood of the empirical law with that of the uniform three-category law.
\begin{solution}
The empirical law is $\widehat Q=(0.5,0.3,0.2)$. Let $U=(1/3,1/3,1/3)$. The divergence is
\[
D(\widehat Q\|U)=0.5\ln1.5+0.3\ln0.9+0.2\ln0.6.
\]
The average negative log likelihood under the uniform law is $\ln3$, whereas under the empirical law it is $H(\widehat Q)$. Their total log-likelihood difference is $40D(\widehat Q\|U)$, and the likelihood ratio is $\exp[40D(\widehat Q\|U)]$. This compares the two models' probabilities for the same observations.
\end{solution}
\end{example}

If only the counts are recorded, the probability of those counts includes the number of sequences that produce them. That multiplicative factor depends on the counts and sample size, while being the same for every proposed law. It cancels from likelihood ratios and leaves the fitted probabilities unchanged.

For a Bernoulli sample with $k$ successes in $N$ observations, the log likelihood is
\[
\ell(p)=k\ln p+(N-k)\ln(1-p).
\]
When $0<k<N$, differentiation gives
\[
\ell'(p)=\frac{N(\widehat q-p)}{p(1-p)},\qquad \widehat q=\frac{k}{N}.
\]
It is positive below the empirical proportion and negative above it, so the maximum occurs at $p=\widehat q$. At that maximum,
\[
-\ell''(\widehat q)=\frac{N}{\widehat q(1-\widehat q)}.
\]
The reciprocal-variance coefficient from the KL Taylor expansion has appeared again, now as local curvature of a data-fitting objective. The next chapter relates it to score fluctuations and Fisher information.

\section{Lyapunov functions}
A scalar function can establish stability by controlling the possible motion of a trajectory. Suppose an equilibrium $q$ lies in an open state domain and a continuously differentiable function $V$ satisfies
\[
V(q)=0,\qquad V(x)>0\quad(x\ne q)
\]
near $q$. Its derivative along $x'=F(x)$ is $\dot V=\nabla V\cdot F$. A nonpositive derivative keeps a solution within the sublevel set $V\leq V(x_0)$.

\begin{definition}[Lyapunov function]\index{Lyapunov function}
A \emph{Lyapunov function} is a function with a strict local minimum at the equilibrium and nonpositive derivative along nearby trajectories. If its derivative is strictly negative away from equilibrium, it is a \emph{strict Lyapunov function}.
\end{definition}

The usual Lyapunov stability theorem follows from nested sublevel sets. A sufficiently small sublevel component around a strict local minimum lies inside any prescribed small neighborhood and is forward invariant when $\dot V\leq0$. This proves stability. When $\dot V<0$ away from the equilibrium, compact annuli around the equilibrium have a negative upper bound for $\dot V$. A trajectory cannot remain in one of these annuli forever, since $V$ would eventually become negative. Applying this observation to successively smaller neighborhoods gives convergence.

The argument applies in any finite dimension when the relevant sublevel components are compactly contained in the domain. A globally strict minimum together with compact sublevels and strict decay gives convergence from every initial state in that domain.

\section{A relative-entropy proof of coexistence}
Consider the family
\[
x'=x(1-x)(a-bx),\qquad b>0,\quad q=\frac ab\in(0,1).
\]
The equilibrium $q$ is the attracting interior state. Define
\begin{equation}\label{eq:binarylyap}
V(x)=D(q\|x)=q\ln\frac qx+(1-q)\ln\frac{1-q}{1-x}.
\end{equation}
Gibbs' inequality gives the required strict minimum. Differentiate with respect to the evolving probability:
\[
V'(x)=-\frac qx+\frac{1-q}{1-x}
=\frac{x-q}{x(1-x)}.
\]
The chain rule now yields
\begin{equation}\label{eq:kldecay}
\boxed{\dot V=(x-q)(a-bx)=-b(x-q)^2.}
\end{equation}
Thus $V$ decreases strictly along every nonstationary interior trajectory.

Because $q$ is positive and less than 1, $V(x)\to\infty$ at both endpoints. Every finite sublevel is compactly contained in $(0,1)$. The strict Lyapunov argument therefore gives convergence to $q$ for every interior initial condition. It also gives a quantitative function with which to monitor the approach.

\begin{example}[The recurring coexistence model]
For $x'=x(1-x)(3-4x)$, write its KL Lyapunov function and its decay rate.
\begin{solution}
Here $q=3/4$ and $b=4$, so
\[
V(x)=\frac34\ln\frac{3}{4x}+\frac14\ln\frac{1}{4(1-x)},
\qquad \dot V=-4\left(x-\frac34\right)^2.
\]
The reference is the equilibrium law. This choice yields decay throughout the interior, including portions of the trajectory where Shannon entropy changes its direction of variation.
\end{solution}
\end{example}

\begin{figure}[htbp]\centering\begin{tikzpicture}\begin{axis}[width=.9\textwidth,height=5.4cm,xmin=0,xmax=12,ymin=0,ymax=1.4,xlabel={time $t$},ylabel={function value},grid=major,legend pos=north east]
\addplot[ink] coordinates {(0.0000000,0.1000000) (0.0750000,0.1186821) (0.1500000,0.1395951) (0.2250000,0.1626096) (0.3000000,0.1874889) (0.3750000,0.2138985) (0.4500000,0.2414290) (0.5250000,0.2696284) (0.6000000,0.2980378) (0.6750000,0.3262250) (0.7500000,0.3538100) (0.8250000,0.3804818) (0.9000000,0.4060045) (0.9750000,0.4302148) (1.0500000,0.4530152) (1.1250000,0.4743629) (1.2000000,0.4942583) (1.2750000,0.5127344) (1.3500000,0.5298465) (1.4250000,0.5456649) (1.5000000,0.5602682) (1.5750000,0.5737389) (1.6500000,0.5861598) (1.7250000,0.5976116) (1.8000000,0.6081719) (1.8750000,0.6179136) (1.9500000,0.6269049) (2.0250000,0.6352086) (2.1000000,0.6428830) (2.1750000,0.6499810) (2.2500000,0.6565513) (2.3250000,0.6626382) (2.4000000,0.6682819) (2.4750000,0.6735191) (2.5500000,0.6783830) (2.6250000,0.6829039) (2.7000000,0.6871094) (2.7750000,0.6910244) (2.8500000,0.6946717) (2.9250000,0.6980721) (3.0000000,0.7012446) (3.0750000,0.7042062) (3.1500000,0.7069729) (3.2250000,0.7095590) (3.3000000,0.7119777) (3.3750000,0.7142411) (3.4500000,0.7163603) (3.5250000,0.7183454) (3.6000000,0.7202060) (3.6750000,0.7219505) (3.7500000,0.7235869) (3.8250000,0.7251225) (3.9000000,0.7265642) (3.9750000,0.7279181) (4.0500000,0.7291901) (4.1250000,0.7303854) (4.2000000,0.7315092) (4.2750000,0.7325659) (4.3500000,0.7335599) (4.4250000,0.7344951) (4.5000000,0.7353753) (4.5750000,0.7362038) (4.6500000,0.7369839) (4.7250000,0.7377186) (4.8000000,0.7384106) (4.8750000,0.7390626) (4.9500000,0.7396769) (5.0250000,0.7402560) (5.1000000,0.7408018) (5.1750000,0.7413163) (5.2500000,0.7418015) (5.3250000,0.7422590) (5.4000000,0.7426906) (5.4750000,0.7430976) (5.5500000,0.7434817) (5.6250000,0.7438440) (5.7000000,0.7441858) (5.7750000,0.7445085) (5.8500000,0.7448130) (5.9250000,0.7451003) (6.0000000,0.7453716) (6.0750000,0.7456277) (6.1500000,0.7458695) (6.2250000,0.7460977) (6.3000000,0.7463133) (6.3750000,0.7465168) (6.4500000,0.7467089) (6.5250000,0.7468904) (6.6000000,0.7470619) (6.6750000,0.7472237) (6.7500000,0.7473767) (6.8250000,0.7475211) (6.9000000,0.7476575) (6.9750000,0.7477864) (7.0500000,0.7479081) (7.1250000,0.7480232) (7.2000000,0.7481318) (7.2750000,0.7482345) (7.3500000,0.7483315) (7.4250000,0.7484231) (7.5000000,0.7485097) (7.5750000,0.7485915) (7.6500000,0.7486689) (7.7250000,0.7487419) (7.8000000,0.7488109) (7.8750000,0.7488762) (7.9500000,0.7489378) (8.0250000,0.7489961) (8.1000000,0.7490511) (8.1750000,0.7491031) (8.2500000,0.7491523) (8.3250000,0.7491988) (8.4000000,0.7492427) (8.4750000,0.7492842) (8.5500000,0.7493234) (8.6250000,0.7493605) (8.7000000,0.7493955) (8.7750000,0.7494286) (8.8500000,0.7494599) (8.9250000,0.7494895) (9.0000000,0.7495175) (9.0750000,0.7495439) (9.1500000,0.7495689) (9.2250000,0.7495925) (9.3000000,0.7496148) (9.3750000,0.7496359) (9.4500000,0.7496558) (9.5250000,0.7496747) (9.6000000,0.7496925) (9.6750000,0.7497093) (9.7500000,0.7497252) (9.8250000,0.7497403) (9.9000000,0.7497545) (9.9750000,0.7497679) (10.0500000,0.7497806) (10.1250000,0.7497926) (10.2000000,0.7498040) (10.2750000,0.7498147) (10.3500000,0.7498248) (10.4250000,0.7498344) (10.5000000,0.7498435) (10.5750000,0.7498520) (10.6500000,0.7498601) (10.7250000,0.7498678) (10.8000000,0.7498750) (10.8750000,0.7498819) (10.9500000,0.7498883) (11.0250000,0.7498944) (11.1000000,0.7499002) (11.1750000,0.7499057) (11.2500000,0.7499108) (11.3250000,0.7499157) (11.4000000,0.7499203) (11.4750000,0.7499247) (11.5500000,0.7499288) (11.6250000,0.7499327) (11.7000000,0.7499364) (11.7750000,0.7499399) (11.8500000,0.7499432) (11.9250000,0.7499463) (12.0000000,0.7499492)};\addlegendentry{$x(t)$}
\addplot[copper] coordinates {(0.0000000,0.3250830) (0.0750000,0.3642908) (0.1500000,0.4042277) (0.2250000,0.4439720) (0.3000000,0.4825613) (0.3750000,0.5190762) (0.4500000,0.5527215) (0.5250000,0.5828889) (0.6000000,0.6091926) (0.6750000,0.6314729) (0.7500000,0.6497733) (0.8250000,0.6642995) (0.9000000,0.6753713) (0.9750000,0.6833754) (1.0500000,0.6887255) (1.1250000,0.6918321) (1.2000000,0.6930812) (1.2750000,0.6928228) (1.3500000,0.6913645) (1.4250000,0.6889708) (1.5000000,0.6858650) (1.5750000,0.6822326) (1.6500000,0.6782258) (1.7250000,0.6739682) (1.8000000,0.6695588) (1.8750000,0.6650763) (1.9500000,0.6605824) (2.0250000,0.6561253) (2.1000000,0.6517414) (2.1750000,0.6474584) (2.2500000,0.6432965) (2.3250000,0.6392700) (2.4000000,0.6353887) (2.4750000,0.6316584) (2.5500000,0.6280823) (2.6250000,0.6246612) (2.7000000,0.6213940) (2.7750000,0.6182786) (2.8500000,0.6153116) (2.9250000,0.6124889) (3.0000000,0.6098061) (3.0750000,0.6072582) (3.1500000,0.6048399) (3.2250000,0.6025461) (3.3000000,0.6003714) (3.3750000,0.5983105) (3.4500000,0.5963581) (3.5250000,0.5945092) (3.6000000,0.5927587) (3.6750000,0.5911017) (3.7500000,0.5895337) (3.8250000,0.5880500) (3.9000000,0.5866463) (3.9750000,0.5853186) (4.0500000,0.5840628) (4.1250000,0.5828751) (4.2000000,0.5817520) (4.2750000,0.5806900) (4.3500000,0.5796858) (4.4250000,0.5787364) (4.5000000,0.5778389) (4.5750000,0.5769903) (4.6500000,0.5761881) (4.7250000,0.5754297) (4.8000000,0.5747128) (4.8750000,0.5740351) (4.9500000,0.5733946) (5.0250000,0.5727890) (5.1000000,0.5722166) (5.1750000,0.5716756) (5.2500000,0.5711642) (5.3250000,0.5706807) (5.4000000,0.5702238) (5.4750000,0.5697919) (5.5500000,0.5693836) (5.6250000,0.5689977) (5.7000000,0.5686330) (5.7750000,0.5682882) (5.8500000,0.5679623) (5.9250000,0.5676542) (6.0000000,0.5673631) (6.0750000,0.5670878) (6.1500000,0.5668277) (6.2250000,0.5665818) (6.3000000,0.5663493) (6.3750000,0.5661296) (6.4500000,0.5659219) (6.5250000,0.5657256) (6.6000000,0.5655401) (6.6750000,0.5653647) (6.7500000,0.5651989) (6.8250000,0.5650422) (6.9000000,0.5648940) (6.9750000,0.5647540) (7.0500000,0.5646216) (7.1250000,0.5644965) (7.2000000,0.5643782) (7.2750000,0.5642665) (7.3500000,0.5641608) (7.4250000,0.5640609) (7.5000000,0.5639665) (7.5750000,0.5638772) (7.6500000,0.5637928) (7.7250000,0.5637131) (7.8000000,0.5636377) (7.8750000,0.5635664) (7.9500000,0.5634991) (8.0250000,0.5634354) (8.1000000,0.5633752) (8.1750000,0.5633183) (8.2500000,0.5632645) (8.3250000,0.5632137) (8.4000000,0.5631656) (8.4750000,0.5631202) (8.5500000,0.5630772) (8.6250000,0.5630366) (8.7000000,0.5629983) (8.7750000,0.5629620) (8.8500000,0.5629277) (8.9250000,0.5628953) (9.0000000,0.5628646) (9.0750000,0.5628357) (9.1500000,0.5628083) (9.2250000,0.5627824) (9.3000000,0.5627579) (9.3750000,0.5627348) (9.4500000,0.5627130) (9.5250000,0.5626923) (9.6000000,0.5626727) (9.6750000,0.5626543) (9.7500000,0.5626368) (9.8250000,0.5626203) (9.9000000,0.5626047) (9.9750000,0.5625900) (10.0500000,0.5625760) (10.1250000,0.5625629) (10.2000000,0.5625504) (10.2750000,0.5625386) (10.3500000,0.5625275) (10.4250000,0.5625170) (10.5000000,0.5625070) (10.5750000,0.5624976) (10.6500000,0.5624887) (10.7250000,0.5624803) (10.8000000,0.5624724) (10.8750000,0.5624649) (10.9500000,0.5624578) (11.0250000,0.5624511) (11.1000000,0.5624447) (11.1750000,0.5624387) (11.2500000,0.5624331) (11.3250000,0.5624277) (11.4000000,0.5624227) (11.4750000,0.5624179) (11.5500000,0.5624133) (11.6250000,0.5624091) (11.7000000,0.5624050) (11.7750000,0.5624012) (11.8500000,0.5623976) (11.9250000,0.5623942) (12.0000000,0.5623909)};\addlegendentry{$H(x(t))$}
\addplot[teal] coordinates {(0.0000000,1.1909438) (0.0750000,1.0677294) (0.1500000,0.9520099) (0.2250000,0.8443332) (0.3000000,0.7450980) (0.3750000,0.6545226) (0.4500000,0.5726296) (0.5250000,0.4992483) (0.6000000,0.4340349) (0.6750000,0.3765057) (0.7500000,0.3260767) (0.8250000,0.2821058) (0.9000000,0.2439290) (0.9750000,0.2108918) (1.0500000,0.1823706) (1.1250000,0.1577879) (1.2000000,0.1366199) (1.2750000,0.1183993) (1.3500000,0.1027148) (1.4250000,0.0892076) (1.5000000,0.0775672) (1.5750000,0.0675263) (1.6500000,0.0588559) (1.7250000,0.0513599) (1.8000000,0.0448711) (1.8750000,0.0392469) (1.9500000,0.0343657) (2.0250000,0.0301236) (2.1000000,0.0264322) (2.1750000,0.0232159) (2.2500000,0.0204101) (2.3250000,0.0179593) (2.4000000,0.0158162) (2.4750000,0.0139400) (2.5500000,0.0122958) (2.6250000,0.0108533) (2.7000000,0.0095866) (2.7750000,0.0084732) (2.8500000,0.0074937) (2.9250000,0.0066313) (3.0000000,0.0058713) (3.0750000,0.0052011) (3.1500000,0.0046096) (3.2250000,0.0040873) (3.3000000,0.0036257) (3.3750000,0.0032176) (3.4500000,0.0028565) (3.5250000,0.0025369) (3.6000000,0.0022538) (3.6750000,0.0020029) (3.7500000,0.0017805) (3.8250000,0.0015833) (3.9000000,0.0014083) (3.9750000,0.0012529) (4.0500000,0.0011150) (4.1250000,0.0009925) (4.2000000,0.0008836) (4.2750000,0.0007869) (4.3500000,0.0007008) (4.4250000,0.0006243) (4.5000000,0.0005563) (4.5750000,0.0004957) (4.6500000,0.0004418) (4.7250000,0.0003938) (4.8000000,0.0003511) (4.8750000,0.0003130) (4.9500000,0.0002791) (5.0250000,0.0002489) (5.1000000,0.0002220) (5.1750000,0.0001981) (5.2500000,0.0001767) (5.3250000,0.0001577) (5.4000000,0.0001407) (5.4750000,0.0001255) (5.5500000,0.0001120) (5.6250000,0.0001000) (5.7000000,0.0000892) (5.7750000,0.0000796) (5.8500000,0.0000711) (5.9250000,0.0000635) (6.0000000,0.0000567) (6.0750000,0.0000506) (6.1500000,0.0000452) (6.2250000,0.0000403) (6.3000000,0.0000360) (6.3750000,0.0000322) (6.4500000,0.0000287) (6.5250000,0.0000256) (6.6000000,0.0000229) (6.6750000,0.0000205) (6.7500000,0.0000183) (6.8250000,0.0000163) (6.9000000,0.0000146) (6.9750000,0.0000130) (7.0500000,0.0000116) (7.1250000,0.0000104) (7.2000000,0.0000093) (7.2750000,0.0000083) (7.3500000,0.0000074) (7.4250000,0.0000066) (7.5000000,0.0000059) (7.5750000,0.0000053) (7.6500000,0.0000047) (7.7250000,0.0000042) (7.8000000,0.0000038) (7.8750000,0.0000034) (7.9500000,0.0000030) (8.0250000,0.0000027) (8.1000000,0.0000024) (8.1750000,0.0000021) (8.2500000,0.0000019) (8.3250000,0.0000017) (8.4000000,0.0000015) (8.4750000,0.0000014) (8.5500000,0.0000012) (8.6250000,0.0000011) (8.7000000,0.0000010) (8.7750000,0.0000009) (8.8500000,0.0000008) (8.9250000,0.0000007) (9.0000000,0.0000006) (9.0750000,0.0000006) (9.1500000,0.0000005) (9.2250000,0.0000004) (9.3000000,0.0000004) (9.3750000,0.0000004) (9.4500000,0.0000003) (9.5250000,0.0000003) (9.6000000,0.0000003) (9.6750000,0.0000002) (9.7500000,0.0000002) (9.8250000,0.0000002) (9.9000000,0.0000002) (9.9750000,0.0000001) (10.0500000,0.0000001) (10.1250000,0.0000001) (10.2000000,0.0000001) (10.2750000,0.0000001) (10.3500000,0.0000001) (10.4250000,0.0000001) (10.5000000,0.0000001) (10.5750000,0.0000001) (10.6500000,0.0000001) (10.7250000,0.0000000) (10.8000000,0.0000000) (10.8750000,0.0000000) (10.9500000,0.0000000) (11.0250000,0.0000000) (11.1000000,0.0000000) (11.1750000,0.0000000) (11.2500000,0.0000000) (11.3250000,0.0000000) (11.4000000,0.0000000) (11.4750000,0.0000000) (11.5500000,0.0000000) (11.6250000,0.0000000) (11.7000000,0.0000000) (11.7750000,0.0000000) (11.8500000,0.0000000) (11.9250000,0.0000000) (12.0000000,0.0000000)};\addlegendentry{$D(3/4\|x(t))$}
\end{axis}\end{tikzpicture}\caption{State, entropy, and relative entropy along the coexistence trajectory from $x_0=0.1$. The KL function decreases throughout the approach to equilibrium.}\end{figure}

\section{The multi-strategy derivative}
For a fixed positive reference $\q$ and an interior replicator solution,
\begin{align}
\frac{d}{dt}D(\q\|\p(t))
&=-\sum_iq_i\frac{p_i'}{p_i}\notag\\
&=-\sum_iq_i(f_i-\bar f)
=\bar f-\sum_iq_if_i.\label{eq:multiklderivative}
\end{align}
The sign depends on the fitness field and reference. The derivative is nonpositive precisely when the reference-weighted fitness is at least the current mean fitness.

\begin{example}[Relative entropy in the cyclic family]
Take $\q=(1/3,1/3,1/3)$ and the cyclic matrix $M_s$ from Chapter 3.
\begin{solution}
The divergence can be written using the product $P=p_Ap_Bp_C$:
\[
D(\q\|\p)=-\frac13\ln(27P).
\]
Consequently
\[
\frac{d}{dt}D(\q\|\p)=-\frac{1-s}{3}(1-3Q),
\]
where $Q=p_Ap_B+p_Bp_C+p_Cp_A$. For $s<1$, this is a strict Lyapunov function for the center. For $s=1$, it is conserved along the closed orbits. For $s>1$, it increases along every noncentral interior trajectory.
\end{solution}
\end{example}

This calculation reinterprets the earlier product criterion in information-theoretic terms. The same divergence supports both decay and conservation when paired with different dynamics.

\section{The local quadratic approximation}\index{KL divergence!quadratic expansion}
Relative entropy has so far measured the difference between distributions at finite separation. Near a fixed reference law, the comparison is governed by the first nonzero term of its Taylor expansion. That term will supply an inner product on the tangent plane of the simplex. It is the precise point at which a statistical comparison begins to produce geometry.
Fix a positive reference vector $\q$ and write a nearby law as $\q+h$, where $\sum_i h_i=0$. Expand each logarithm:
\begin{align*}
D(\q\|\q+h)
&=-\sum_iq_i\ln\left(1+\frac{h_i}{q_i}\right)\\
&=-\sum_i h_i+\frac12\sum_i\frac{h_i^2}{q_i}+O(\|h\|^3).
\end{align*}
The linear term vanishes by normalization, giving
\begin{equation}\label{eq:klquadratic}
\boxed{D(\q\|\q+h)=\frac12\sum_i\frac{h_i^2}{q_i}+O(\|h\|^3).}
\end{equation}
The estimate is local around a fixed positive law. In particular, the allowed displacements must be small compared with the probabilities that approach zero near the boundary.

For a two-outcome displacement $(h,-h)$, the expression reduces to
\begin{equation}\label{eq:binaryklquadratic}
D(q\|q+h)=\frac{h^2}{2q(1-q)}+O(|h|^3).
\end{equation}
Reversing the arguments gives the same leading quadratic term. Directional differences first enter at higher order. The symmetric positive quadratic expression is the object from which Fisher geometry is constructed.

\begin{example}[Equal coordinate changes at different states]
Use the quadratic approximation for $h=0.02$ at $q=0.2$ and $q=0.5$.
\begin{solution}
The approximations, in nats, are
\begin{align*}
q=0.2:&\quad \frac{0.02^2}{2(0.2)(0.8)}=0.00125,\\
q=0.5:&\quad \frac{0.02^2}{2(0.5)(0.5)}=0.0008.
\end{align*}
The same small probability displacement has a larger leading information cost at the first reference state.
\end{solution}
\end{example}

There is also an exact connection with entropy. Let $F(P)=-H(P)=\sum_i p_i\ln p_i$. Then
\begin{equation}\label{eq:bregman}
D(Q\|P)=F(Q)-F(P)-\sum_i(\ln p_i+1)(q_i-p_i).
\end{equation}
The divergence is the difference between the value of the convex function $F$ and its tangent-plane approximation at $P$. Its Hessian in probability coordinates is $\diag(1/p_i)$, the same quadratic form seen in \eqref{eq:klquadratic}.

\FloatBarrier
\section{Exercises}
\begin{exercises}
\item For $Q=(3/4,1/4)$ and $P=(1/2,1/2)$, compute both directed divergences to six decimal places.
\item Prove $\ln u\leq u-1$ using calculus and characterize equality.
\item Determine which directed divergences are finite between $(1/2,1/2,0)$ and $(1/3,1/3,1/3)$. Compute the finite value.
\item Prove the strict convexity of $P\mapsto D(Q\|P)$ for positive $Q$.
\item For $q=2/5$, derive $V'(x)$ when $V(x)=D(q\|x)$. Apply it to $x'=x(1-x)(2-5x)$.
\item Give a complete Lyapunov argument for convergence in the preceding problem, including compactness of sublevels.
\item Show that along the mixture from $P$ to $Q$, the function $D(Q\|M(t))$ decreases strictly unless $P=Q$.
\item For $Q=(0.5,0.3,0.2)$ and $P=(0.15,0.20,0.65)$, compute cross-entropy and both divergences numerically.
\item Derive \eqref{eq:multiklderivative}. State the condition that makes its right-hand side nonpositive.
\item For the balanced cyclic model, express $D((1/3,1/3,1/3)\|\p)$ in terms of $p_Ap_Bp_C$ and explain its conservation.
\item Expand the binary divergence at $x=q+h$ through quadratic order. Identify the reason its linear term vanishes.
\item Derive the cubic term in $D(\q\|\q+h)$. Explain why the expansion is not uniform over the entire positive simplex.
\item For $q=0.2$ and $h=0.02$, compare the exact divergence with its quadratic approximation. Repeat after halving $h$.
\item Verify the exact entropy identity \eqref{eq:bregman} by expanding all sums.
\item Show that $D(q\|x)\geq2(x-q)^2$ for interior binary probabilities. Use $V'(x)=(x-q)/[x(1-x)]$ and integrate from $q$ to $x$.
\item On a compact interval $[\delta,1-\delta]$ containing both $x$ and $q$, prove $D(q\|x)\leq(x-q)^2/[2\delta(1-\delta)]$. Combine this with \eqref{eq:kldecay} to obtain an exponential bound for $V$ along a trapped trajectory.
\item \emph{Exploration.} For the coexistence model from $x_0=0.1$, graph $x(t)$, $H(x(t))$, and $D(3/4\|x(t))$. Explain each curve using its derivative.
\item \emph{Exploration.} Fix a positive categorical reference and inspect both directed KL landscapes. Describe their boundary behavior by identifying which coordinates create infinite loss.
\end{exercises}

\chapter{Fisher Information and Local Geometry}
\label{ch:fisher}

The quadratic term of relative entropy assigns a positive squared length to an infinitesimal probability change. Allowing this inner product to vary with the underlying distribution produces the Fisher information metric. Calculus then supplies lengths of curves, angles between tangent vectors, and rates of motion through statistical state space.

The key transition is from comparing two finite distributions to measuring a small displacement from one distribution. The leading quadratic term of relative entropy gives that displacement a length, much as a quadratic form in Euclidean geometry measures a vector. Because the coefficients depend on the starting distribution, the resulting geometry changes from point to point. This chapter develops that local measurement before the next chapter identifies its global shape.

\section{A metric from the KL expansion}
At a positive law $\p$, tangent vectors belong to
\[
T_\p\Delta^{n-1}_{\!\circ}=\{v\in\R^n:\sum_i v_i=0\}.
\]
For a tangent displacement $h$, Chapter 6 gives
\[
D(\p\|\p+h)=\frac12\sum_i\frac{h_i^2}{p_i}+O(\|h\|^3).
\]
This motivates the inner product
\begin{equation}\label{eq:fishermetric}
\boxed{g_\p(u,v)=\sum_i\frac{u_iv_i}{p_i}.}
\end{equation}
It is bilinear and symmetric, and $g_\p(v,v)>0$ for every nonzero tangent vector because each denominator is positive. Its coefficients depend smoothly on $\p$ in the interior.

\begin{definition}[Fisher information metric]\index{Fisher information metric}\index{Riemannian metric}
The inner products \eqref{eq:fishermetric}, one at each positive law, form the \emph{Fisher information metric}. A smoothly varying inner product on tangent spaces is called a \emph{Riemannian metric}. The associated tangent norm is $\|v\|_F=\sqrt{g_\p(v,v)}$.
\end{definition}

An inner product must be evaluated at a specified base point. A velocity vector can have the same numerical components at two distributions and different Fisher norms. A rare category has a large reciprocal weight, so changes in that probability make a larger contribution to the local quadratic log-loss comparison.

The normalization used here gives $D(\p\|\p+h)\approx\tfrac12\|h\|_F^2$. The factor $1/2$ belongs to Taylor's formula; the Fisher metric itself is the quadratic form $\sum_i h_i^2/p_i$.

\section{Scores and parameterized families}\index{Fisher information matrix}
An arbitrary tangent vector represents an infinitesimal change in all category probabilities. In a statistical family, only changes produced by varying the model's parameters are available. The score measures the change in log probability of each outcome when a parameter changes; its expected square gives the Fisher length of that parameter velocity.
Let $P_\theta$ be a smoothly parameterized positive categorical family, with probabilities $p_i(\theta)$. For one observed outcome $i$, its \emph{score}\index{score} is
\[
S_\theta(i)=\frac{\partial}{\partial\theta}\ln p_i(\theta)
=\frac{p_i'(\theta)}{p_i(\theta)}.
\]
The expected score is zero:
\[
\E_\theta[S_\theta]=\sum_i p_i\frac{p_i'}{p_i}=\sum_i p_i'=0.
\]
Its expected square is the \emph{Fisher information}
\begin{equation}\label{eq:scoreinfo}
I(\theta)=\E_\theta[S_\theta^2]=\sum_i\frac{(p_i'(\theta))^2}{p_i(\theta)}.
\end{equation}
This is exactly the squared Fisher norm of the family velocity $\partial_\theta P$. A parameter change $d\theta$ has squared length $I(\theta)d\theta^2$.

For a parameter vector $\theta=(\theta^1,\ldots,\theta^m)$, the score components $S_a(i)=\partial_a\ln p_i$ give
\begin{equation}\label{eq:informationmatrix}
G_{ab}(\theta)=\E_\theta[S_aS_b]
=\sum_i\frac{\partial_ap_i\,\partial_bp_i}{p_i}.
\end{equation}
Substituting $dp_i=\sum_a\partial_ap_i\,d\theta^a$ into the categorical metric produces
\[
ds^2=d\theta^TG(\theta)d\theta.
\]
Thus the score definition is the categorical metric expressed along the chosen family. The matrix is positive definite when the parameter derivatives give independent tangent directions. Redundant parameters can instead give a singular matrix.

\begin{example}[Bernoulli information from scores]
For $\Pr_p(X=1)=p$, compute the score and information.
\begin{solution}
The score is $1/p$ for outcome 1 and $-1/(1-p)$ for outcome 0. Its mean is zero, and
\[
I(p)=p\left(\frac1p\right)^2+(1-p)\left(\frac1{1-p}\right)^2
=\frac1{p(1-p)}.
\]
The same result follows from the tangent displacement $(dp,-dp)$ in \eqref{eq:fishermetric}.
\end{solution}
\end{example}

Differentiating $\sum_i\partial_ap_i=0$ once more also gives
\[
G_{ab}=-\E_\theta[\partial_a\partial_b\ln p_i].
\]
The identity relates average score products to average log-likelihood curvature. For finite positive laws all differentiations pass through finite sums. Continuous families require regularity assumptions allowing the analogous differentiations under an integral.

\section{Coordinate changes and statistical meaning}\index{coordinate transformation}
Suppose $\theta=\theta(\eta)$ is a smooth change of parameters, with Jacobian $J=\partial\theta/\partial\eta$. The chain rule gives $d\theta=Jd\eta$, so
\begin{equation}\label{eq:metrictransform}
G_\eta=J^TG_\theta J.
\end{equation}
The matrix coefficients change, while the squared length of the same distribution change remains the same.

\begin{example}[Information in log-odds coordinates]
Use $L=\ln[p/(1-p)]$ in the Bernoulli family.
\begin{solution}
Since $dp/dL=p(1-p)$, the transformation law gives
\[
I_L=I_p\left(\frac{dp}{dL}\right)^2=p(1-p).
\]
Thus $ds^2=p(1-p)dL^2$. The log-odds coordinate simplifies constant-fitness dynamics, but its metric coefficient still varies. A different coordinate below makes the length coefficient constant.
\end{solution}
\end{example}

For nearby parameters, the KL comparison becomes
\[
D(P_\theta\|P_{\theta+d\theta})
=\frac12d\theta^TG(\theta)d\theta+O(\|d\theta\|^3).
\]
This formula connects the geometric length with the expected excess log loss of neighboring statistical models. Its value is independent of the coordinates used to describe the same displacement to the displayed order.

\section{Information from independent observations}\index{independent observations}
For $N$ independent observations, the log likelihood is a sum, so its score is $S_N=\sum_{r=1}^N S^{(r)}$. The individual scores have mean zero, and independence eliminates the cross terms in the expected square. Consequently
\begin{equation}\label{eq:sampleinformation}
I_N(\theta)=N I(\theta).
\end{equation}
The corresponding local squared lengths for the full sample model scale by $N$, and lengths scale by $\sqrt N$. A small change that is difficult to distinguish with one observation becomes more detectable with many observations.

An estimator $T$ is a numerical function of the observed sample, intended to estimate $\theta$. It is \emph{unbiased} on an interval when $\E_\theta[T]=\theta$ throughout that interval. For a finite positive sample space, differentiate this identity:
\[
1=\frac{d}{d\theta}\E_\theta[T]=\E_\theta[T S_N]
=\E_\theta[(T-\theta)S_N],
\]
where the last equality uses $\E_\theta[S_N]=0$. Cauchy--Schwarz now gives
\[
1\leq\E_\theta[(T-\theta)^2]\E_\theta[S_N^2]
=\Var_\theta(T)\,N I(\theta).
\]
Thus the \emph{Cram\'er--Rao bound}\index{Cramer--Rao bound@Cram\'er--Rao bound} is
\begin{equation}\label{eq:crbound}
\Var_\theta(T)\geq\frac1{NI(\theta)}.
\end{equation}
The proof uses unbiasedness in a neighborhood, positivity, and differentiability. Those hypotheses specify the estimation problem to which the bound applies.

\begin{example}[An estimator attaining the bound]
Take $N$ independent Bernoulli observations and estimate $p$ by their empirical proportion.
\begin{solution}
Chapter 4 gives $\E[\widehat p]=p$ and $\Var(\widehat p)=p(1-p)/N$. Since $I(p)=1/[p(1-p)]$, this variance equals $1/(NI(p))$. The empirical proportion attains the bound for every interior parameter.
\end{solution}
\end{example}

\section{Curve length and the Bernoulli distance}\index{curve length}\index{Fisher distance!Bernoulli}
The metric assigns a length to one tangent vector at one distribution. A curve length accumulates these infinitesimal measurements as the distribution moves. It provides a way to compare paths whose coordinate speeds differ and, by taking the shortest available length, to define distance between endpoint laws.
For a smooth probability curve $\p(t)$, its Fisher length is
\begin{equation}\label{eq:fisherlength}
L_F(\p)=\int_a^b\sqrt{\sum_i\frac{(p_i'(t))^2}{p_i(t)}}\dd t.
\end{equation}
The integrand is speed, computed with the metric at the current law. A monotone change of the time parameter leaves this integral unchanged, because the velocity scaling is canceled by the time differential.

The \emph{intrinsic distance}\index{intrinsic distance} between two laws in a family is the infimum of lengths of curves in that family connecting them. In a Bernoulli family,
\[
ds^2=\frac{dp^2}{p(1-p)},\qquad
d_F(a,b)=\left|\int_a^b\frac{dp}{\sqrt{p(1-p)}}\right|.
\]
Set $p=\sin^2(\theta/2)$ with $0<\theta<\pi$. Then
\[
dp=\tfrac12\sin\theta\,d\theta,
\qquad \sqrt{p(1-p)}=\tfrac12\sin\theta.
\]
The line element becomes $ds^2=d\theta^2$, yielding
\begin{equation}\label{eq:bernoullidistance}
\boxed{d_F(a,b)=2\left|\arcsin\sqrt b-\arcsin\sqrt a\right|.}
\end{equation}
Any monotone Bernoulli curve between $a$ and $b$ has this length. A path that reverses direction adds extra travel. Constant-speed travel corresponds to an affine function of the arc coordinate $\theta$, and its probabilities are $p(t)=\sin^2(\theta(t)/2)$.

\begin{example}[An exact Bernoulli distance]
Find the Fisher distance from $p=1/4$ to $p=3/4$.
\begin{solution}
The corresponding $\theta$ values are $\pi/3$ and $2\pi/3$. Their separation is $\pi/3$. The coordinate probability difference is $1/2$; Fisher distance uses the integrated statistical line element.
\end{solution}
\end{example}

\begin{figure}[htbp]
\centering
\begin{tikzpicture}
\begin{axis}[width=.88\textwidth,height=5.1cm,xmin=0,xmax=1,ymin=0,ymax=3.2,xlabel={probability $p$},ylabel={arc coordinate $\theta$},ytick={0,1.04719755,1.57079633,2.0943951,3.14159265},yticklabels={0,$\pi/3$,$\pi/2$,$2\pi/3$,$\pi$},grid=major]
\addplot[teal,domain=.0001:.9999,samples=160]{2*asin(sqrt(x))*pi/180};
\addplot[only marks,ink] coordinates {(.25,1.04719755)(.5,1.57079633)(.75,2.0943951)};
\end{axis}
\end{tikzpicture}
\caption{The Bernoulli arc coordinate. Its difference is exact Fisher distance for monotone travel, and its total range has length $\pi$.}
\end{figure}

The coefficient $1/[p(1-p)]$ tends to infinity near either boundary, but its square root has an integrable singularity. The limiting distance from 0 to 1 is $\pi$. Boundary states can therefore be added at finite distance even though the metric was defined on positive laws.

\section{The categorical metric in independent coordinates}
For $\p=(x,y,z)$ with $z=1-x-y$, a tangent displacement has $dz=-dx-dy$. Hence
\begin{align*}
ds^2&=\frac{dx^2}{x}+\frac{dy^2}{y}+\frac{(dx+dy)^2}{z}\\
&=\begin{pmatrix}dx&dy\end{pmatrix}
\underbrace{\begin{pmatrix}1/x+1/z&1/z\\1/z&1/y+1/z\end{pmatrix}}_{G(x,y)}
\begin{pmatrix}dx\\dy\end{pmatrix}.
\end{align*}
Its determinant is $1/(xyz)>0$, so the matrix is positive definite throughout the interior. Its inverse is
\begin{equation}\label{eq:inversefisher}
G^{-1}(x,y)=\begin{pmatrix}x(1-x)&-xy\\-xy&y(1-y)\end{pmatrix}.
\end{equation}
The off-diagonal terms record the effect of the dependent third coordinate. Changing $x$ or $y$ necessarily changes $z$ as well.

For nonzero tangent vectors, the Fisher angle is defined by
\[
\cos\alpha_F=\frac{g_\p(u,v)}{\sqrt{g_\p(u,u)g_\p(v,v)}}.
\]
The usual inner-product identities and Cauchy--Schwarz inequality apply in each tangent space.

\begin{example}[Angles at a categorical law]
At $\q=(0.5,0.3,0.2)$ take
\[
u=(0.1,-0.05,-0.05),\qquad v=(0,0.1,-0.1).
\]
Compare their Fisher and Euclidean angles.
\begin{solution}
The Euclidean dot product is zero. In the Fisher metric,
\[
g_\q(u,v)=\frac1{120},\qquad
g_\q(u,u)=\frac{49}{1200},\qquad
g_\q(v,v)=\frac1{12}.
\]
Thus $\cos\alpha_F=1/7$ and $\alpha_F\approx81.79^\circ$, while the Euclidean angle is $90^\circ$.
\end{solution}
\end{example}

\begin{figure}[htbp]\centering\begin{tikzpicture}[scale=6]
\fill[soft] (0,0)--(1,0)--(.5,.8660254)--cycle;
\draw[ink!75] (0,0)--(1,0)--(.5,.8660254)--cycle;
\node[below left,font=\small] at (0,0){$A$};\node[below right,font=\small] at (1,0){$B$};\node[above,font=\small] at (.5,.8660254){$C$};
\filldraw[fill=teal!13,draw=teal,line width=1pt] plot coordinates {(0.3939340,0.2274379) (0.3979884,0.2208988) (0.4024454,0.2146272) (0.4072874,0.2086478) (0.4124953,0.2029842) (0.4180486,0.1976589) (0.4239252,0.1926927) (0.4301021,0.1881053) (0.4365549,0.1839148) (0.4432581,0.1801378) (0.4501851,0.1767891) (0.4573088,0.1738820) (0.4646010,0.1714279) (0.4720328,0.1694366) (0.4795751,0.1679158) (0.4871979,0.1668716) (0.4948713,0.1663081) (0.5025649,0.1662275) (0.5102484,0.1666302) (0.5178914,0.1675145) (0.5254639,0.1688770) (0.5329358,0.1707123) (0.5402778,0.1730131) (0.5474608,0.1757704) (0.5544565,0.1789733) (0.5612372,0.1826091) (0.5677763,0.1866635) (0.5740480,0.1911205) (0.5800274,0.1959626) (0.5856909,0.2011705) (0.5910163,0.2067237) (0.5959825,0.2126004) (0.6005698,0.2187773) (0.6047603,0.2252300) (0.6085373,0.2319332) (0.6118860,0.2388603) (0.6147931,0.2459840) (0.6172472,0.2532761) (0.6192386,0.2607080) (0.6207594,0.2682502) (0.6218036,0.2758731) (0.6223671,0.2835464) (0.6224476,0.2912400) (0.6220450,0.2989235) (0.6211606,0.3065666) (0.6197981,0.3141390) (0.6179628,0.3216110) (0.6156620,0.3289529) (0.6129047,0.3361359) (0.6097018,0.3431316) (0.6060660,0.3499124) (0.6020116,0.3564515) (0.5975546,0.3627231) (0.5927126,0.3687025) (0.5875047,0.3743660) (0.5819514,0.3796914) (0.5760748,0.3846576) (0.5698979,0.3892450) (0.5634451,0.3934354) (0.5567419,0.3972125) (0.5498149,0.4005611) (0.5426912,0.4034683) (0.5353990,0.4059223) (0.5279672,0.4079137) (0.5204249,0.4094345) (0.5128021,0.4104787) (0.5051287,0.4110422) (0.4974351,0.4111228) (0.4897516,0.4107201) (0.4821086,0.4098358) (0.4745361,0.4084733) (0.4670642,0.4066380) (0.4597222,0.4043371) (0.4525392,0.4015799) (0.4455435,0.3983770) (0.4387628,0.3947412) (0.4322237,0.3906867) (0.4259520,0.3862297) (0.4199726,0.3813877) (0.4143091,0.3761798) (0.4089837,0.3706266) (0.4040175,0.3647499) (0.3994302,0.3585730) (0.3952397,0.3521202) (0.3914627,0.3454171) (0.3881140,0.3384900) (0.3852069,0.3313663) (0.3827528,0.3240742) (0.3807614,0.3166423) (0.3792406,0.3091001) (0.3781964,0.3014772) (0.3776329,0.2938038) (0.3775524,0.2861102) (0.3779550,0.2784267) (0.3788394,0.2707837) (0.3802019,0.2632113) (0.3820372,0.2557393) (0.3843380,0.2483974) (0.3870953,0.2412144) (0.3902982,0.2342187) (0.3939340,0.2274379)}--cycle;
\fill[teal] (0.5000000,0.2886751) circle(.008);
\node[font=\small] at (0.6000000,0.4800000){$(1/3,1/3,1/3)$};
\filldraw[fill=copper!13,draw=copper,line width=1pt] plot coordinates {(0.0600000,0.0346410) (0.0622837,0.0310958) (0.0649135,0.0277695) (0.0678791,0.0246755) (0.0711688,0.0218259) (0.0747697,0.0192319) (0.0786674,0.0169038) (0.0828466,0.0148508) (0.0872909,0.0130810) (0.0919827,0.0116013) (0.0969034,0.0104175) (0.1020337,0.0095345) (0.1073532,0.0089556) (0.1128411,0.0086832) (0.1184756,0.0087182) (0.1242346,0.0090607) (0.1300952,0.0097091) (0.1360344,0.0106611) (0.1420287,0.0119127) (0.1480544,0.0134591) (0.1540879,0.0152942) (0.1601052,0.0174106) (0.1660826,0.0198002) (0.1719965,0.0224534) (0.1778237,0.0253597) (0.1835410,0.0285078) (0.1891260,0.0318851) (0.1945565,0.0354784) (0.1998112,0.0392734) (0.2048694,0.0432552) (0.2097109,0.0474081) (0.2143169,0.0517157) (0.2186690,0.0561610) (0.2227501,0.0607264) (0.2265440,0.0653939) (0.2300359,0.0701451) (0.2332120,0.0749613) (0.2360596,0.0798234) (0.2385676,0.0847123) (0.2407260,0.0896086) (0.2425264,0.0944931) (0.2439617,0.0993464) (0.2450261,0.1041494) (0.2457155,0.1088832) (0.2460272,0.1135290) (0.2459598,0.1180686) (0.2455138,0.1224840) (0.2446908,0.1267578) (0.2434941,0.1308731) (0.2419285,0.1348137) (0.2400000,0.1385641) (0.2377163,0.1421093) (0.2350865,0.1454355) (0.2321209,0.1485296) (0.2288312,0.1513792) (0.2252303,0.1539732) (0.2213326,0.1563012) (0.2171534,0.1583543) (0.2127091,0.1601241) (0.2080173,0.1616038) (0.2030966,0.1627875) (0.1979663,0.1636706) (0.1926468,0.1642495) (0.1871589,0.1645219) (0.1815244,0.1644868) (0.1757654,0.1641444) (0.1699048,0.1634959) (0.1639656,0.1625440) (0.1579713,0.1612924) (0.1519456,0.1597460) (0.1459121,0.1579109) (0.1398948,0.1557944) (0.1339174,0.1534049) (0.1280035,0.1507517) (0.1221763,0.1478453) (0.1164590,0.1446973) (0.1108740,0.1413200) (0.1054435,0.1377267) (0.1001888,0.1339317) (0.0951306,0.1299498) (0.0902891,0.1257969) (0.0856831,0.1214894) (0.0813310,0.1170441) (0.0772499,0.1124787) (0.0734560,0.1078112) (0.0699641,0.1030600) (0.0667880,0.0982438) (0.0639404,0.0933817) (0.0614324,0.0884928) (0.0592740,0.0835965) (0.0574736,0.0787120) (0.0560383,0.0738587) (0.0549739,0.0690557) (0.0542845,0.0643219) (0.0539728,0.0596760) (0.0540402,0.0551365) (0.0544862,0.0507211) (0.0553092,0.0464473) (0.0565059,0.0423320) (0.0580715,0.0383913) (0.0600000,0.0346410)}--cycle;
\fill[copper] (0.1500000,0.0866025) circle(.008);
\node[font=\small] at (0.2000000,-0.0900000){$(.8,.1,.1)$};
\end{tikzpicture}\caption{Tangent ellipses at two base laws. Each displacement from its ellipse center has Fisher norm $0.3$ in the metric frozen at that center. These are tangent-norm contours, rather than finite-distance circles.}\end{figure}

\section{Tangent velocities as centered random variables}
There is another way to read the reciprocal probabilities in the Fisher metric. At a positive law $\p$, assign to a tangent vector $v$ the category values
\[
a_i=\frac{v_i}{p_i}.
\]
Their probability-weighted mean is $\sum_i p_i a_i=\sum_i v_i=0$. Thus each tangent vector corresponds to a mean-zero random variable on the outcome labels. Conversely, any values $a_i$ of mean zero give a tangent vector $v_i=p_i a_i$.

For two tangent vectors, this correspondence gives
\[
g_\p(u,v)=\sum_i p_i\left(\frac{u_i}{p_i}\right)
\left(\frac{v_i}{p_i}\right).
\]
Their Fisher inner product is the expected product, or covariance, of their centered random-variable representatives. Their squared norms are the variances of those representatives. The metric weights become the ordinary probability weights after velocities are expressed as relative rates of change.

For a curve $\p(t)$ through $\p$, the values $p_i'/p_i$ are the derivatives of its log probabilities. They are precisely its score values when $t$ is treated as the parameter. The score definition of information and the tangent-vector definition therefore agree at the level of an explicit isometry between two finite-dimensional inner-product spaces.

\begin{example}[A tangent velocity and its score values]
At $\p=(1/2,1/4,1/4)$ take $v=(0,0.1,-0.1)$. Calculate its centered random-variable representative and norm.
\begin{solution}
The values $v_i/p_i$ are $(0,0.4,-0.4)$, of weighted mean zero. Their variance is $0.25(0.4)^2+0.25(-0.4)^2=0.08$. Direct calculation gives $g_\p(v,v)=0.01/(1/4)+0.01/(1/4)=0.08$, so the norm is $\sqrt{0.08}$.
\end{solution}
\end{example}

Sampling covariance gives a complementary description. In the two independent coordinates of a three-category law, Chapter 4 gave
\[
N\operatorname{Cov}\begin{pmatrix}\widehat p_A\\\widehat p_B\end{pmatrix}
=\begin{pmatrix}x(1-x)&-xy\\-xy&y(1-y)\end{pmatrix}.
\]
This is exactly the inverse information matrix \eqref{eq:inversefisher}. Directions with larger empirical variance have smaller information cost in the corresponding inverse quadratic form. The normalization constraint is already incorporated by using two independent coordinates; the full three-coordinate covariance matrix has a zero direction and cannot be inverted as an ordinary three-by-three matrix.

\section{Replicator speed and fitness variance}\index{Fisher speed}
Substitute $p_i'=p_i(f_i-\bar f)$ into the length integrand:
\begin{equation}\label{eq:speedvariance}
\boxed{\|\p'\|_F^2=\sum_i\frac{p_i^2(f_i-\bar f)^2}{p_i}
=\Var_\p(f).}
\end{equation}
This identity holds for every smooth interior replicator trajectory. The fitness field can be constant, state-dependent, or cyclic. Equal fitness among all present types gives zero velocity and zero variance. Greater fitness variation corresponds to greater instantaneous squared speed through Fisher geometry.

\begin{example}[Checking the speed identity]
Let $p_A=1/4$, $p_B=3/4$, and $(f_A,f_B)=(2,0)$.
\begin{solution}
The mean fitness is $1/2$, the A velocity is $3/8$, and
\[
\Var(f)=\frac14\left(\frac32\right)^2+\frac34\left(-\frac12\right)^2=\frac34.
\]
The squared speed calculated from the Bernoulli metric is $(3/8)^2/[(1/4)(3/4)]=3/4$, as required.
\end{solution}
\end{example}

With constant fitnesses, Chapter 5 showed that $\bar f'=\Var(f)$. In that case the increase in mean fitness equals squared Fisher speed. For general state-dependent fitnesses, differentiating their values also contributes to $\bar f'$. The speed identity itself remains valid and will be combined with a potential function in Chapter 9.

\FloatBarrier
\section{Exercises}
\begin{exercises}
\item Derive the categorical Fisher metric from the quadratic KL expansion and verify positivity on tangent vectors.
\item Derive the Bernoulli information from both score values and show that the expected score is zero.
\item Compute $I(p)$ at $p=0.1,0.25,0.5,0.9$. Explain the symmetry between the first and last values.
\item Prove the coordinate transformation rule \eqref{eq:metrictransform} using the chain rule for scores.
\item Compute the exact Fisher lengths of the moves $0.10\to0.15$ and $0.50\to0.55$. Compare the first-order approximations based at the left endpoint.
\item For Bernoulli log-odds $L$, derive $I_L=p(1-p)$ directly from the scores rather than from the metric transformation law.
\item Find the constant-speed Bernoulli path from $1/4$ to $3/4$ over $0\leq t\leq1$ and calculate its speed.
\item Explain why a monotone change of curve parameter preserves Fisher length. What changes if the new parameter repeatedly reverses direction?
\item Derive the three-category matrix $G$ and verify both its determinant and the inverse in \eqref{eq:inversefisher}.
\item At $\p=(1/3,1/3,1/3)$, compare the Fisher and ambient Euclidean inner products of arbitrary tangent vectors. What happens to angles?
\item Verify every inner product in the worked angle example and obtain $\cos\alpha_F=1/7$.
\item Prove $I_N=NI$ for independent observations and identify the use of the zero-mean score.
\item For $N=100$ Bernoulli observations at $p=1/4$, compute the Cram\'er--Rao bound. Compare it with the empirical proportion variance.
\item Explain why an unbiasedness assertion only at one parameter value is insufficient for the derivative argument in the Cram\'er--Rao proof.
\item Derive \eqref{eq:speedvariance} and verify it at $\p=(0.2,0.3,0.5)$ with $f=(3,1,2)$.
\item For state-dependent fitnesses along a trajectory, show $\bar f'=\Var_\p(f)+\sum_i p_i\frac{d}{dt}f_i(\p(t))$.
\item \emph{Exploration.} Compare Fisher norms and angles at several simplex points. Keep the coordinate vectors fixed and explain the effect of moving the base distribution.
\item A one-dimensional metric is $ds^2=g(x)dx^2$. Find a vector $w$ such that $g(x)wv=\Phi'(x)v$ for every scalar $v$. Apply your result to the Bernoulli metric and $\Phi(x)=3x-2x^2$.
\end{exercises}

\chapter{Square-Root Coordinates and Spherical Geometry}
\label{ch:sphere}

The categorical Fisher metric has a concrete global model. Taking square roots of the probabilities sends the positive simplex to an orthant of a round sphere. The metric becomes ordinary spherical length, making exact distances and shortest paths accessible through Euclidean vector calculations.

Imagine drawing a straight segment between two points of the probability triangle. That segment is a natural mixture of the endpoint laws, but it need not be the shortest path for the Fisher metric. Square-root coordinates reveal why: the metric measures length on a curved spherical surface. The distinction between a straight mixture and a shortest geometric path will be visible in both drawings.

\section{The square-root isometry}
For a positive categorical law define
\begin{equation}\label{eq:sqrtmap}
s_i=2\sqrt{p_i}.
\end{equation}
Normalization gives $\sum_i s_i^2=4$, and positivity gives $s_i>0$. Thus $s$ belongs to the positive orthant of the radius-2 sphere in $\R^n$.

Differentiate the map:
\[
ds_i=\frac{dp_i}{\sqrt{p_i}},\qquad
\sum_i ds_i^2=\sum_i\frac{dp_i^2}{p_i}=ds_F^2.
\]
The Euclidean inner product on the image tangent spaces is exactly the Fisher inner product. The map is one-to-one, with inverse $p_i=s_i^2/4$.

\begin{proposition}[Square-root representation]\index{square-root coordinates}\index{isometry}
The positive categorical simplex with the Fisher metric is isometric to the positive orthant of the round sphere $S^{n-1}(2)$.
\end{proposition}
\begin{proof}
The map and its inverse are smooth on their positive domains, and the calculation above proves preservation of tangent inner products. For a tangent displacement, $\sum_i dp_i=0$ also gives $s\cdot ds=2\sum_i dp_i=0$, confirming that the image displacement is tangent to the sphere.
\end{proof}

An \emph{isometry} preserves the lengths and angles measured by the two metrics. It also preserves lengths of curves, since those lengths are integrals of tangent norms. The flat triangle remains useful as a probability diagram; the sphere displays the geometry supplied by Fisher information.

\begin{example}[A distribution and a tangent vector on the sphere]
Map $P=(1/2,1/3,1/6)$ to the sphere. Describe the image of a tangent vector $v$.
\begin{solution}
The image is $(\sqrt2,2/\sqrt3,2/\sqrt6)$, whose squared norm is $2+4/3+2/3=4$. The image velocity has components $v_i/\sqrt{p_i}$. Its dot product with the image point is $2\sum_i v_i=0$, so it lies in the tangent plane of the sphere.
\end{solution}
\end{example}

\begin{figure}[htbp]\centering\begin{tikzpicture}[x={(1.6cm,-.45cm)},y={(-1.2cm,-.6cm)},z={(0cm,1.65cm)},>=Stealth]
\draw[->,ink!50] (0,0,0)--(2.3,0,0) node[right]{$s_1$};\draw[->,ink!50] (0,0,0)--(0,2.3,0) node[left]{$s_2$};\draw[->,ink!50] (0,0,0)--(0,0,2.3) node[above]{$s_3$};
\draw[ink!24,line width=.45pt] plot coordinates {(0.3128689,0.0000000,1.9753767) (0.3127617,0.0081900,1.9753767) (0.3124402,0.0163743,1.9753767) (0.3119045,0.0245474,1.9753767) (0.3111550,0.0327037,1.9753767) (0.3101923,0.0408376,1.9753767) (0.3090170,0.0489435,1.9753767) (0.3076299,0.0570158,1.9753767) (0.3060320,0.0650491,1.9753767) (0.3042243,0.0730378,1.9753767) (0.3022082,0.0809764,1.9753767) (0.2999849,0.0888596,1.9753767) (0.2975560,0.0966818,1.9753767) (0.2949232,0.1044378,1.9753767) (0.2920883,0.1121222,1.9753767) (0.2890532,0.1197298,1.9753767) (0.2858200,0.1272553,1.9753767) (0.2823909,0.1346935,1.9753767) (0.2787683,0.1420395,1.9753767) (0.2749546,0.1492882,1.9753767) (0.2709524,0.1564345,1.9753767) (0.2667646,0.1634736,1.9753767) (0.2623940,0.1704006,1.9753767) (0.2578435,0.1772109,1.9753767) (0.2531163,0.1838997,1.9753767) (0.2482156,0.1904625,1.9753767) (0.2431448,0.1968948,1.9753767) (0.2379074,0.2031921,1.9753767) (0.2325069,0.2093502,1.9753767) (0.2269471,0.2153648,1.9753767) (0.2212317,0.2212317,1.9753767) (0.2153648,0.2269471,1.9753767) (0.2093502,0.2325069,1.9753767) (0.2031921,0.2379074,1.9753767) (0.1968948,0.2431448,1.9753767) (0.1904625,0.2482156,1.9753767) (0.1838997,0.2531163,1.9753767) (0.1772109,0.2578435,1.9753767) (0.1704006,0.2623940,1.9753767) (0.1634736,0.2667646,1.9753767) (0.1564345,0.2709524,1.9753767) (0.1492882,0.2749546,1.9753767) (0.1420395,0.2787683,1.9753767) (0.1346935,0.2823909,1.9753767) (0.1272553,0.2858200,1.9753767) (0.1197298,0.2890532,1.9753767) (0.1121222,0.2920883,1.9753767) (0.1044378,0.2949232,1.9753767) (0.0966818,0.2975560,1.9753767) (0.0888596,0.2999849,1.9753767) (0.0809764,0.3022082,1.9753767) (0.0730378,0.3042243,1.9753767) (0.0650491,0.3060320,1.9753767) (0.0570158,0.3076299,1.9753767) (0.0489435,0.3090170,1.9753767) (0.0408376,0.3101923,1.9753767) (0.0327037,0.3111550,1.9753767) (0.0245474,0.3119045,1.9753767) (0.0163743,0.3124402,1.9753767) (0.0081900,0.3127617,1.9753767) (0.0000000,0.3128689,1.9753767)};
\draw[ink!24,line width=.45pt] plot coordinates {(0.0000000,0.0000000,2.0000000) (0.0517093,0.0081900,1.9993146) (0.1033832,0.0163743,1.9972591) (0.1549863,0.0245474,1.9938347) (0.2064831,0.0327037,1.9890438) (0.2578384,0.0408376,1.9828897) (0.3090170,0.0489435,1.9753767) (0.3599838,0.0570158,1.9665098) (0.4107039,0.0650491,1.9562952) (0.4611425,0.0730378,1.9447398) (0.5112651,0.0809764,1.9318517) (0.5610373,0.0888596,1.9176395) (0.6104250,0.0966818,1.9021130) (0.6593943,0.1044378,1.8852830) (0.7079117,0.1121222,1.8671609) (0.7559439,0.1197298,1.8477591) (0.8034581,0.1272553,1.8270909) (0.8504216,0.1346935,1.8051706) (0.8968022,0.1420395,1.7820130) (0.9425683,0.1492882,1.7576342) (0.9876883,0.1564345,1.7320508) (1.0321315,0.1634736,1.7052803) (1.0758672,0.1704006,1.6773411) (1.1188657,0.1772109,1.6482524) (1.1610973,0.1838997,1.6180340) (1.2025331,0.1904625,1.5867067) (1.2431448,0.1968948,1.5542919) (1.2829045,0.2031921,1.5208119) (1.3217850,0.2093502,1.4862897) (1.3597596,0.2153648,1.4507487) (1.3968022,0.2212317,1.4142136) (1.4328876,0.2269471,1.3767092) (1.4679910,0.2325069,1.3382612) (1.5020882,0.2379074,1.2988961) (1.5351560,0.2431448,1.2586408) (1.5671717,0.2482156,1.2175229) (1.5981133,0.2531163,1.1755705) (1.6279597,0.2578435,1.1328125) (1.6566903,0.2623940,1.0892781) (1.6842855,0.2667646,1.0449971) (1.7107264,0.2709524,1.0000000) (1.7359948,0.2749546,0.9543175) (1.7600735,0.2787683,0.9079810) (1.7829459,0.2823909,0.8610222) (1.8045964,0.2858200,0.8134733) (1.8250101,0.2890532,0.7653669) (1.8441730,0.2920883,0.7167359) (1.8620720,0.2949232,0.6676137) (1.8786949,0.2975560,0.6180340) (1.8940301,0.2999849,0.5680307) (1.9080674,0.3022082,0.5176381) (1.9207969,0.3042243,0.4668907) (1.9322100,0.3060320,0.4158234) (1.9422988,0.3076299,0.3644711) (1.9510565,0.3090170,0.3128689) (1.9584771,0.3101923,0.2610524) (1.9645554,0.3111550,0.2090569) (1.9692873,0.3119045,0.1569182) (1.9726695,0.3124402,0.1046719) (1.9746998,0.3127617,0.0523539) (1.9753767,0.3128689,0.0000000)};
\draw[ink!24,line width=.45pt] plot coordinates {(0.6180340,0.0000000,1.9021130) (0.6178222,0.0161782,1.9021130) (0.6171870,0.0323454,1.9021130) (0.6161288,0.0484904,1.9021130) (0.6146483,0.0646021,1.9021130) (0.6127466,0.0806696,1.9021130) (0.6104250,0.0966818,1.9021130) (0.6076850,0.1126277,1.9021130) (0.6045285,0.1284965,1.9021130) (0.6009577,0.1442772,1.9021130) (0.5969750,0.1599590,1.9021130) (0.5925832,0.1755311,1.9021130) (0.5877853,0.1909830,1.9021130) (0.5825845,0.2063040,1.9021130) (0.5769844,0.2214836,1.9021130) (0.5709890,0.2365114,1.9021130) (0.5646021,0.2513771,1.9021130) (0.5578284,0.2660705,1.9021130) (0.5506723,0.2805816,1.9021130) (0.5431388,0.2949003,1.9021130) (0.5352331,0.3090170,1.9021130) (0.5269606,0.3229219,1.9021130) (0.5183269,0.3366054,1.9021130) (0.5093380,0.3500583,1.9021130) (0.5000000,0.3632713,1.9021130) (0.4903193,0.3762353,1.9021130) (0.4803026,0.3889414,1.9021130) (0.4699567,0.4013810,1.9021130) (0.4592888,0.4135455,1.9021130) (0.4483060,0.4254265,1.9021130) (0.4370160,0.4370160,1.9021130) (0.4254265,0.4483060,1.9021130) (0.4135455,0.4592888,1.9021130) (0.4013810,0.4699567,1.9021130) (0.3889414,0.4803026,1.9021130) (0.3762353,0.4903193,1.9021130) (0.3632713,0.5000000,1.9021130) (0.3500583,0.5093380,1.9021130) (0.3366054,0.5183269,1.9021130) (0.3229219,0.5269606,1.9021130) (0.3090170,0.5352331,1.9021130) (0.2949003,0.5431388,1.9021130) (0.2805816,0.5506723,1.9021130) (0.2660705,0.5578284,1.9021130) (0.2513771,0.5646021,1.9021130) (0.2365114,0.5709890,1.9021130) (0.2214836,0.5769844,1.9021130) (0.2063040,0.5825845,1.9021130) (0.1909830,0.5877853,1.9021130) (0.1755311,0.5925832,1.9021130) (0.1599590,0.5969750,1.9021130) (0.1442772,0.6009577,1.9021130) (0.1284965,0.6045285,1.9021130) (0.1126277,0.6076850,1.9021130) (0.0966818,0.6104250,1.9021130) (0.0806696,0.6127466,1.9021130) (0.0646021,0.6146483,1.9021130) (0.0484904,0.6161288,1.9021130) (0.0323454,0.6171870,1.9021130) (0.0161782,0.6178222,1.9021130) (0.0000000,0.6180340,1.9021130)};
\draw[ink!24,line width=.45pt] plot coordinates {(0.0000000,0.0000000,2.0000000) (0.0497915,0.0161782,1.9993146) (0.0995489,0.0323454,1.9972591) (0.1492381,0.0484904,1.9938347) (0.1988250,0.0646021,1.9890438) (0.2482756,0.0806696,1.9828897) (0.2975560,0.0966818,1.9753767) (0.3466326,0.1126277,1.9665098) (0.3954715,0.1284965,1.9562952) (0.4440395,0.1442772,1.9447398) (0.4923031,0.1599590,1.9318517) (0.5402293,0.1755311,1.9176395) (0.5877853,0.1909830,1.9021130) (0.6349384,0.2063040,1.8852830) (0.6816563,0.2214836,1.8671609) (0.7279071,0.2365114,1.8477591) (0.7736591,0.2513771,1.8270909) (0.8188808,0.2660705,1.8051706) (0.8635412,0.2805816,1.7820130) (0.9076099,0.2949003,1.7576342) (0.9510565,0.3090170,1.7320508) (0.9938513,0.3229219,1.7052803) (1.0359650,0.3366054,1.6773411) (1.0773687,0.3500583,1.6482524) (1.1180340,0.3632713,1.6180340) (1.1579330,0.3762353,1.5867067) (1.1970385,0.3889414,1.5542919) (1.2353236,0.4013810,1.5208119) (1.2727620,0.4135455,1.4862897) (1.3093282,0.4254265,1.4507487) (1.3449970,0.4370160,1.4142136) (1.3797440,0.4483060,1.3767092) (1.4135455,0.4592888,1.3382612) (1.4463781,0.4699567,1.2988961) (1.4782195,0.4803026,1.2586408) (1.5090477,0.4903193,1.2175229) (1.5388418,0.5000000,1.1755705) (1.5675812,0.5093380,1.1328125) (1.5952462,0.5183269,1.0892781) (1.6218180,0.5269606,1.0449971) (1.6472782,0.5352331,1.0000000) (1.6716095,0.5431388,0.9543175) (1.6947951,0.5506723,0.9079810) (1.7168192,0.5578284,0.8610222) (1.7376667,0.5646021,0.8134733) (1.7573233,0.5709890,0.7653669) (1.7757755,0.5769844,0.7167359) (1.7930107,0.5825845,0.6676137) (1.8090170,0.5877853,0.6180340) (1.8237835,0.5925832,0.5680307) (1.8373001,0.5969750,0.5176381) (1.8495575,0.6009577,0.4668907) (1.8605473,0.6045285,0.4158234) (1.8702620,0.6076850,0.3644711) (1.8786949,0.6104250,0.3128689) (1.8858402,0.6127466,0.2610524) (1.8916931,0.6146483,0.2090569) (1.8962495,0.6161288,0.1569182) (1.8995063,0.6171870,0.1046719) (1.9014612,0.6178222,0.0523539) (1.9021130,0.6180340,0.0000000)};
\draw[ink!24,line width=.45pt] plot coordinates {(0.9079810,0.0000000,1.7820130) (0.9076699,0.0237682,1.7820130) (0.9067366,0.0475201,1.7820130) (0.9051820,0.0712394,1.7820130) (0.9030070,0.0949099,1.7820130) (0.9002131,0.1185153,1.7820130) (0.8968022,0.1420395,1.7820130) (0.8927768,0.1654664,1.7820130) (0.8881394,0.1887799,1.7820130) (0.8828934,0.2119640,1.7820130) (0.8770423,0.2350028,1.7820130) (0.8705901,0.2578805,1.7820130) (0.8635412,0.2805816,1.7820130) (0.8559006,0.3030903,1.7820130) (0.8476733,0.3253913,1.7820130) (0.8388651,0.3474693,1.7820130) (0.8294819,0.3693091,1.7820130) (0.8195303,0.3908959,1.7820130) (0.8090170,0.4122147,1.7820130) (0.7979492,0.4332511,1.7820130) (0.7863346,0.4539905,1.7820130) (0.7741811,0.4744188,1.7820130) (0.7614969,0.4945219,1.7820130) (0.7482909,0.5142861,1.7820130) (0.7345721,0.5336978,1.7820130) (0.7203498,0.5527438,1.7820130) (0.7056338,0.5714110,1.7820130) (0.6904342,0.5896865,1.7820130) (0.6747614,0.6075579,1.7820130) (0.6586261,0.6250129,1.7820130) (0.6420395,0.6420395,1.7820130) (0.6250129,0.6586261,1.7820130) (0.6075579,0.6747614,1.7820130) (0.5896865,0.6904342,1.7820130) (0.5714110,0.7056338,1.7820130) (0.5527438,0.7203498,1.7820130) (0.5336978,0.7345721,1.7820130) (0.5142861,0.7482909,1.7820130) (0.4945219,0.7614969,1.7820130) (0.4744188,0.7741811,1.7820130) (0.4539905,0.7863346,1.7820130) (0.4332511,0.7979492,1.7820130) (0.4122147,0.8090170,1.7820130) (0.3908959,0.8195303,1.7820130) (0.3693091,0.8294819,1.7820130) (0.3474693,0.8388651,1.7820130) (0.3253913,0.8476733,1.7820130) (0.3030903,0.8559006,1.7820130) (0.2805816,0.8635412,1.7820130) (0.2578805,0.8705901,1.7820130) (0.2350028,0.8770423,1.7820130) (0.2119640,0.8828934,1.7820130) (0.1887799,0.8881394,1.7820130) (0.1654664,0.8927768,1.7820130) (0.1420395,0.8968022,1.7820130) (0.1185153,0.9002131,1.7820130) (0.0949099,0.9030070,1.7820130) (0.0712394,0.9051820,1.7820130) (0.0475201,0.9067366,1.7820130) (0.0237682,0.9076699,1.7820130) (0.0000000,0.9079810,1.7820130)};
\draw[ink!24,line width=.45pt] plot coordinates {(0.0000000,0.0000000,2.0000000) (0.0466477,0.0237682,1.9993146) (0.0932634,0.0475201,1.9972591) (0.1398151,0.0712394,1.9938347) (0.1862711,0.0949099,1.9890438) (0.2325994,0.1185153,1.9828897) (0.2787683,0.1420395,1.9753767) (0.3247461,0.1654664,1.9665098) (0.3705013,0.1887799,1.9562952) (0.4160027,0.2119640,1.9447398) (0.4612189,0.2350028,1.9318517) (0.5061191,0.2578805,1.9176395) (0.5506723,0.2805816,1.9021130) (0.5948482,0.3030903,1.8852830) (0.6386164,0.3253913,1.8671609) (0.6819469,0.3474693,1.8477591) (0.7248100,0.3693091,1.8270909) (0.7671764,0.3908959,1.8051706) (0.8090170,0.4122147,1.7820130) (0.8503031,0.4332511,1.7576342) (0.8910065,0.4539905,1.7320508) (0.9310993,0.4744188,1.7052803) (0.9705539,0.4945219,1.6773411) (1.0093433,0.5142861,1.6482524) (1.0474410,0.5336978,1.6180340) (1.0848208,0.5527438,1.5867067) (1.1214571,0.5714110,1.5542919) (1.1573249,0.5896865,1.5208119) (1.1923995,0.6075579,1.4862897) (1.2266568,0.6250129,1.4507487) (1.2600735,0.6420395,1.4142136) (1.2926266,0.6586261,1.3767092) (1.3242938,0.6747614,1.3382612) (1.3550534,0.6904342,1.2988961) (1.3848842,0.7056338,1.2586408) (1.4137660,0.7203498,1.2175229) (1.4416788,0.7345721,1.1755705) (1.4686036,0.7482909,1.1328125) (1.4945219,0.7614969,1.0892781) (1.5194159,0.7741811,1.0449971) (1.5432686,0.7863346,1.0000000) (1.5660636,0.7979492,0.9543175) (1.5877853,0.8090170,0.9079810) (1.6084188,0.8195303,0.8610222) (1.6279499,0.8294819,0.8134733) (1.6463654,0.8388651,0.7653669) (1.6636525,0.8476733,0.7167359) (1.6797994,0.8559006,0.6676137) (1.6947951,0.8635412,0.6180340) (1.7086293,0.8705901,0.5680307) (1.7212924,0.8770423,0.5176381) (1.7327759,0.8828934,0.4668907) (1.7430718,0.8881394,0.4158234) (1.7521731,0.8927768,0.3644711) (1.7600735,0.8968022,0.3128689) (1.7667677,0.9002131,0.2610524) (1.7722510,0.9030070,0.2090569) (1.7765197,0.9051820,0.1569182) (1.7795709,0.9067366,0.1046719) (1.7814024,0.9076699,0.0523539) (1.7820130,0.9079810,0.0000000)};
\draw[ink!24,line width=.45pt] plot coordinates {(1.1755705,0.0000000,1.6180340) (1.1751677,0.0307728,1.6180340) (1.1739594,0.0615246,1.6180340) (1.1719466,0.0922342,1.6180340) (1.1691306,0.1228806,1.6180340) (1.1655133,0.1534427,1.6180340) (1.1610973,0.1838997,1.6180340) (1.1558855,0.2142307,1.6180340) (1.1498815,0.2444149,1.6180340) (1.1430894,0.2744315,1.6180340) (1.1355139,0.3042600,1.6180340) (1.1271602,0.3338801,1.6180340) (1.1180340,0.3632713,1.6180340) (1.1081415,0.3924135,1.6180340) (1.0974896,0.4212868,1.6180340) (1.0860855,0.4498714,1.6180340) (1.0739371,0.4781476,1.6180340) (1.0610526,0.5060961,1.6180340) (1.0474410,0.5336978,1.6180340) (1.0331115,0.5609338,1.6180340) (1.0180739,0.5877853,1.6180340) (1.0023386,0.6142339,1.6180340) (0.9859164,0.6402616,1.6180340) (0.9688184,0.6658505,1.6180340) (0.9510565,0.6909830,1.6180340) (0.9326428,0.7156420,1.6180340) (0.9135899,0.7398105,1.6180340) (0.8939108,0.7634720,1.6180340) (0.8736191,0.7866102,1.6180340) (0.8527287,0.8092093,1.6180340) (0.8312539,0.8312539,1.6180340) (0.8092093,0.8527287,1.6180340) (0.7866102,0.8736191,1.6180340) (0.7634720,0.8939108,1.6180340) (0.7398105,0.9135899,1.6180340) (0.7156420,0.9326428,1.6180340) (0.6909830,0.9510565,1.6180340) (0.6658505,0.9688184,1.6180340) (0.6402616,0.9859164,1.6180340) (0.6142339,1.0023386,1.6180340) (0.5877853,1.0180739,1.6180340) (0.5609338,1.0331115,1.6180340) (0.5336978,1.0474410,1.6180340) (0.5060961,1.0610526,1.6180340) (0.4781476,1.0739371,1.6180340) (0.4498714,1.0860855,1.6180340) (0.4212868,1.0974896,1.6180340) (0.3924135,1.1081415,1.6180340) (0.3632713,1.1180340,1.6180340) (0.3338801,1.1271602,1.6180340) (0.3042600,1.1355139,1.6180340) (0.2744315,1.1430894,1.6180340) (0.2444149,1.1498815,1.6180340) (0.2142307,1.1558855,1.6180340) (0.1838997,1.1610973,1.6180340) (0.1534427,1.1655133,1.6180340) (0.1228806,1.1691306,1.6180340) (0.0922342,1.1719466,1.6180340) (0.0615246,1.1739594,1.6180340) (0.0307728,1.1751677,1.6180340) (0.0000000,1.1755705,1.6180340)};
\draw[ink!24,line width=.45pt] plot coordinates {(0.0000000,0.0000000,2.0000000) (0.0423552,0.0307728,1.9993146) (0.0846814,0.0615246,1.9972591) (0.1269495,0.0922342,1.9938347) (0.1691306,0.1228806,1.9890438) (0.2111958,0.1534427,1.9828897) (0.2531163,0.1838997,1.9753767) (0.2948633,0.2142307,1.9665098) (0.3364082,0.2444149,1.9562952) (0.3777225,0.2744315,1.9447398) (0.4187780,0.3042600,1.9318517) (0.4595465,0.3338801,1.9176395) (0.5000000,0.3632713,1.9021130) (0.5401108,0.3924135,1.8852830) (0.5798515,0.4212868,1.8671609) (0.6191948,0.4498714,1.8477591) (0.6581137,0.4781476,1.8270909) (0.6965816,0.5060961,1.8051706) (0.7345721,0.5336978,1.7820130) (0.7720591,0.5609338,1.7576342) (0.8090170,0.5877853,1.7320508) (0.8454204,0.6142339,1.7052803) (0.8812445,0.6402616,1.6773411) (0.9164645,0.6658505,1.6482524) (0.9510565,0.6909830,1.6180340) (0.9849967,0.7156420,1.5867067) (1.0182618,0.7398105,1.5542919) (1.0508290,0.7634720,1.5208119) (1.0826761,0.7866102,1.4862897) (1.1137811,0.8092093,1.4507487) (1.1441228,0.8312539,1.4142136) (1.1736804,0.8527287,1.3767092) (1.2024336,0.8736191,1.3382612) (1.2303627,0.8939108,1.2988961) (1.2574486,0.9135899,1.2586408) (1.2836727,0.9326428,1.2175229) (1.3090170,0.9510565,1.1755705) (1.3334642,0.9688184,1.1328125) (1.3569975,0.9859164,1.0892781) (1.3796008,1.0023386,1.0449971) (1.4012585,1.0180739,1.0000000) (1.4219560,1.0331115,0.9543175) (1.4416788,1.0474410,0.9079810) (1.4604137,1.0610526,0.8610222) (1.4781476,1.0739371,0.8134733) (1.4948685,1.0860855,0.7653669) (1.5105649,1.0974896,0.7167359) (1.5252260,1.1081415,0.6676137) (1.5388418,1.1180340,0.6180340) (1.5514029,1.1271602,0.5680307) (1.5629008,1.1355139,0.5176381) (1.5733276,1.1430894,0.4668907) (1.5826761,1.1498815,0.4158234) (1.5909399,1.1558855,0.3644711) (1.5981133,1.1610973,0.3128689) (1.6041915,1.1655133,0.2610524) (1.6091702,1.1691306,0.2090569) (1.6130461,1.1719466,0.1569182) (1.6158165,1.1739594,0.1046719) (1.6174795,1.1751677,0.0523539) (1.6180340,1.1755705,0.0000000)};
\draw[ink!24,line width=.45pt] plot coordinates {(1.4142136,0.0000000,1.4142136) (1.4137289,0.0370198,1.4142136) (1.4122754,0.0740142,1.4142136) (1.4098540,0.1109579,1.4142136) (1.4064664,0.1478256,1.4142136) (1.4021148,0.1845919,1.4142136) (1.3968022,0.2212317,1.4142136) (1.3905324,0.2577200,1.4142136) (1.3833096,0.2940315,1.4142136) (1.3751387,0.3301416,1.4142136) (1.3660254,0.3660254,1.4142136) (1.3559759,0.4016584,1.4142136) (1.3449970,0.4370160,1.4142136) (1.3330964,0.4720742,1.4142136) (1.3202821,0.5068088,1.4142136) (1.3065630,0.5411961,1.4142136) (1.2919484,0.5752125,1.4142136) (1.2764484,0.6088346,1.4142136) (1.2600735,0.6420395,1.4142136) (1.2428351,0.6748044,1.4142136) (1.2247449,0.7071068,1.4142136) (1.2058153,0.7389246,1.4142136) (1.1860593,0.7702359,1.4142136) (1.1654904,0.8010194,1.4142136) (1.1441228,0.8312539,1.4142136) (1.1219711,0.8609187,1.4142136) (1.0990504,0.8899934,1.4142136) (1.0753764,0.9184582,1.4142136) (1.0509655,0.9462936,1.4142136) (1.0258343,0.9734804,1.4142136) (1.0000000,1.0000000,1.4142136) (0.9734804,1.0258343,1.4142136) (0.9462936,1.0509655,1.4142136) (0.9184582,1.0753764,1.4142136) (0.8899934,1.0990504,1.4142136) (0.8609187,1.1219711,1.4142136) (0.8312539,1.1441228,1.4142136) (0.8010194,1.1654904,1.4142136) (0.7702359,1.1860593,1.4142136) (0.7389246,1.2058153,1.4142136) (0.7071068,1.2247449,1.4142136) (0.6748044,1.2428351,1.4142136) (0.6420395,1.2600735,1.4142136) (0.6088346,1.2764484,1.4142136) (0.5752125,1.2919484,1.4142136) (0.5411961,1.3065630,1.4142136) (0.5068088,1.3202821,1.4142136) (0.4720742,1.3330964,1.4142136) (0.4370160,1.3449970,1.4142136) (0.4016584,1.3559759,1.4142136) (0.3660254,1.3660254,1.4142136) (0.3301416,1.3751387,1.4142136) (0.2940315,1.3833096,1.4142136) (0.2577200,1.3905324,1.4142136) (0.2212317,1.3968022,1.4142136) (0.1845919,1.4021148,1.4142136) (0.1478256,1.4064664,1.4142136) (0.1109579,1.4098540,1.4142136) (0.0740142,1.4122754,1.4142136) (0.0370198,1.4137289,1.4142136) (0.0000000,1.4142136,1.4142136)};
\draw[ink!24,line width=.45pt] plot coordinates {(0.0000000,0.0000000,2.0000000) (0.0370198,0.0370198,1.9993146) (0.0740142,0.0740142,1.9972591) (0.1109579,0.1109579,1.9938347) (0.1478256,0.1478256,1.9890438) (0.1845919,0.1845919,1.9828897) (0.2212317,0.2212317,1.9753767) (0.2577200,0.2577200,1.9665098) (0.2940315,0.2940315,1.9562952) (0.3301416,0.3301416,1.9447398) (0.3660254,0.3660254,1.9318517) (0.4016584,0.4016584,1.9176395) (0.4370160,0.4370160,1.9021130) (0.4720742,0.4720742,1.8852830) (0.5068088,0.5068088,1.8671609) (0.5411961,0.5411961,1.8477591) (0.5752125,0.5752125,1.8270909) (0.6088346,0.6088346,1.8051706) (0.6420395,0.6420395,1.7820130) (0.6748044,0.6748044,1.7576342) (0.7071068,0.7071068,1.7320508) (0.7389246,0.7389246,1.7052803) (0.7702359,0.7702359,1.6773411) (0.8010194,0.8010194,1.6482524) (0.8312539,0.8312539,1.6180340) (0.8609187,0.8609187,1.5867067) (0.8899934,0.8899934,1.5542919) (0.9184582,0.9184582,1.5208119) (0.9462936,0.9462936,1.4862897) (0.9734804,0.9734804,1.4507487) (1.0000000,1.0000000,1.4142136) (1.0258343,1.0258343,1.3767092) (1.0509655,1.0509655,1.3382612) (1.0753764,1.0753764,1.2988961) (1.0990504,1.0990504,1.2586408) (1.1219711,1.1219711,1.2175229) (1.1441228,1.1441228,1.1755705) (1.1654904,1.1654904,1.1328125) (1.1860593,1.1860593,1.0892781) (1.2058153,1.2058153,1.0449971) (1.2247449,1.2247449,1.0000000) (1.2428351,1.2428351,0.9543175) (1.2600735,1.2600735,0.9079810) (1.2764484,1.2764484,0.8610222) (1.2919484,1.2919484,0.8134733) (1.3065630,1.3065630,0.7653669) (1.3202821,1.3202821,0.7167359) (1.3330964,1.3330964,0.6676137) (1.3449970,1.3449970,0.6180340) (1.3559759,1.3559759,0.5680307) (1.3660254,1.3660254,0.5176381) (1.3751387,1.3751387,0.4668907) (1.3833096,1.3833096,0.4158234) (1.3905324,1.3905324,0.3644711) (1.3968022,1.3968022,0.3128689) (1.4021148,1.4021148,0.2610524) (1.4064664,1.4064664,0.2090569) (1.4098540,1.4098540,0.1569182) (1.4122754,1.4122754,0.1046719) (1.4137289,1.4137289,0.0523539) (1.4142136,1.4142136,0.0000000)};
\draw[ink!24,line width=.45pt] plot coordinates {(1.6180340,0.0000000,1.1755705) (1.6174795,0.0423552,1.1755705) (1.6158165,0.0846814,1.1755705) (1.6130461,0.1269495,1.1755705) (1.6091702,0.1691306,1.1755705) (1.6041915,0.2111958,1.1755705) (1.5981133,0.2531163,1.1755705) (1.5909399,0.2948633,1.1755705) (1.5826761,0.3364082,1.1755705) (1.5733276,0.3777225,1.1755705) (1.5629008,0.4187780,1.1755705) (1.5514029,0.4595465,1.1755705) (1.5388418,0.5000000,1.1755705) (1.5252260,0.5401108,1.1755705) (1.5105649,0.5798515,1.1755705) (1.4948685,0.6191948,1.1755705) (1.4781476,0.6581137,1.1755705) (1.4604137,0.6965816,1.1755705) (1.4416788,0.7345721,1.1755705) (1.4219560,0.7720591,1.1755705) (1.4012585,0.8090170,1.1755705) (1.3796008,0.8454204,1.1755705) (1.3569975,0.8812445,1.1755705) (1.3334642,0.9164645,1.1755705) (1.3090170,0.9510565,1.1755705) (1.2836727,0.9849967,1.1755705) (1.2574486,1.0182618,1.1755705) (1.2303627,1.0508290,1.1755705) (1.2024336,1.0826761,1.1755705) (1.1736804,1.1137811,1.1755705) (1.1441228,1.1441228,1.1755705) (1.1137811,1.1736804,1.1755705) (1.0826761,1.2024336,1.1755705) (1.0508290,1.2303627,1.1755705) (1.0182618,1.2574486,1.1755705) (0.9849967,1.2836727,1.1755705) (0.9510565,1.3090170,1.1755705) (0.9164645,1.3334642,1.1755705) (0.8812445,1.3569975,1.1755705) (0.8454204,1.3796008,1.1755705) (0.8090170,1.4012585,1.1755705) (0.7720591,1.4219560,1.1755705) (0.7345721,1.4416788,1.1755705) (0.6965816,1.4604137,1.1755705) (0.6581137,1.4781476,1.1755705) (0.6191948,1.4948685,1.1755705) (0.5798515,1.5105649,1.1755705) (0.5401108,1.5252260,1.1755705) (0.5000000,1.5388418,1.1755705) (0.4595465,1.5514029,1.1755705) (0.4187780,1.5629008,1.1755705) (0.3777225,1.5733276,1.1755705) (0.3364082,1.5826761,1.1755705) (0.2948633,1.5909399,1.1755705) (0.2531163,1.5981133,1.1755705) (0.2111958,1.6041915,1.1755705) (0.1691306,1.6091702,1.1755705) (0.1269495,1.6130461,1.1755705) (0.0846814,1.6158165,1.1755705) (0.0423552,1.6174795,1.1755705) (0.0000000,1.6180340,1.1755705)};
\draw[ink!24,line width=.45pt] plot coordinates {(0.0000000,0.0000000,2.0000000) (0.0307728,0.0423552,1.9993146) (0.0615246,0.0846814,1.9972591) (0.0922342,0.1269495,1.9938347) (0.1228806,0.1691306,1.9890438) (0.1534427,0.2111958,1.9828897) (0.1838997,0.2531163,1.9753767) (0.2142307,0.2948633,1.9665098) (0.2444149,0.3364082,1.9562952) (0.2744315,0.3777225,1.9447398) (0.3042600,0.4187780,1.9318517) (0.3338801,0.4595465,1.9176395) (0.3632713,0.5000000,1.9021130) (0.3924135,0.5401108,1.8852830) (0.4212868,0.5798515,1.8671609) (0.4498714,0.6191948,1.8477591) (0.4781476,0.6581137,1.8270909) (0.5060961,0.6965816,1.8051706) (0.5336978,0.7345721,1.7820130) (0.5609338,0.7720591,1.7576342) (0.5877853,0.8090170,1.7320508) (0.6142339,0.8454204,1.7052803) (0.6402616,0.8812445,1.6773411) (0.6658505,0.9164645,1.6482524) (0.6909830,0.9510565,1.6180340) (0.7156420,0.9849967,1.5867067) (0.7398105,1.0182618,1.5542919) (0.7634720,1.0508290,1.5208119) (0.7866102,1.0826761,1.4862897) (0.8092093,1.1137811,1.4507487) (0.8312539,1.1441228,1.4142136) (0.8527287,1.1736804,1.3767092) (0.8736191,1.2024336,1.3382612) (0.8939108,1.2303627,1.2988961) (0.9135899,1.2574486,1.2586408) (0.9326428,1.2836727,1.2175229) (0.9510565,1.3090170,1.1755705) (0.9688184,1.3334642,1.1328125) (0.9859164,1.3569975,1.0892781) (1.0023386,1.3796008,1.0449971) (1.0180739,1.4012585,1.0000000) (1.0331115,1.4219560,0.9543175) (1.0474410,1.4416788,0.9079810) (1.0610526,1.4604137,0.8610222) (1.0739371,1.4781476,0.8134733) (1.0860855,1.4948685,0.7653669) (1.0974896,1.5105649,0.7167359) (1.1081415,1.5252260,0.6676137) (1.1180340,1.5388418,0.6180340) (1.1271602,1.5514029,0.5680307) (1.1355139,1.5629008,0.5176381) (1.1430894,1.5733276,0.4668907) (1.1498815,1.5826761,0.4158234) (1.1558855,1.5909399,0.3644711) (1.1610973,1.5981133,0.3128689) (1.1655133,1.6041915,0.2610524) (1.1691306,1.6091702,0.2090569) (1.1719466,1.6130461,0.1569182) (1.1739594,1.6158165,0.1046719) (1.1751677,1.6174795,0.0523539) (1.1755705,1.6180340,0.0000000)};
\draw[ink!24,line width=.45pt] plot coordinates {(1.7820130,0.0000000,0.9079810) (1.7814024,0.0466477,0.9079810) (1.7795709,0.0932634,0.9079810) (1.7765197,0.1398151,0.9079810) (1.7722510,0.1862711,0.9079810) (1.7667677,0.2325994,0.9079810) (1.7600735,0.2787683,0.9079810) (1.7521731,0.3247461,0.9079810) (1.7430718,0.3705013,0.9079810) (1.7327759,0.4160027,0.9079810) (1.7212924,0.4612189,0.9079810) (1.7086293,0.5061191,0.9079810) (1.6947951,0.5506723,0.9079810) (1.6797994,0.5948482,0.9079810) (1.6636525,0.6386164,0.9079810) (1.6463654,0.6819469,0.9079810) (1.6279499,0.7248100,0.9079810) (1.6084188,0.7671764,0.9079810) (1.5877853,0.8090170,0.9079810) (1.5660636,0.8503031,0.9079810) (1.5432686,0.8910065,0.9079810) (1.5194159,0.9310993,0.9079810) (1.4945219,0.9705539,0.9079810) (1.4686036,1.0093433,0.9079810) (1.4416788,1.0474410,0.9079810) (1.4137660,1.0848208,0.9079810) (1.3848842,1.1214571,0.9079810) (1.3550534,1.1573249,0.9079810) (1.3242938,1.1923995,0.9079810) (1.2926266,1.2266568,0.9079810) (1.2600735,1.2600735,0.9079810) (1.2266568,1.2926266,0.9079810) (1.1923995,1.3242938,0.9079810) (1.1573249,1.3550534,0.9079810) (1.1214571,1.3848842,0.9079810) (1.0848208,1.4137660,0.9079810) (1.0474410,1.4416788,0.9079810) (1.0093433,1.4686036,0.9079810) (0.9705539,1.4945219,0.9079810) (0.9310993,1.5194159,0.9079810) (0.8910065,1.5432686,0.9079810) (0.8503031,1.5660636,0.9079810) (0.8090170,1.5877853,0.9079810) (0.7671764,1.6084188,0.9079810) (0.7248100,1.6279499,0.9079810) (0.6819469,1.6463654,0.9079810) (0.6386164,1.6636525,0.9079810) (0.5948482,1.6797994,0.9079810) (0.5506723,1.6947951,0.9079810) (0.5061191,1.7086293,0.9079810) (0.4612189,1.7212924,0.9079810) (0.4160027,1.7327759,0.9079810) (0.3705013,1.7430718,0.9079810) (0.3247461,1.7521731,0.9079810) (0.2787683,1.7600735,0.9079810) (0.2325994,1.7667677,0.9079810) (0.1862711,1.7722510,0.9079810) (0.1398151,1.7765197,0.9079810) (0.0932634,1.7795709,0.9079810) (0.0466477,1.7814024,0.9079810) (0.0000000,1.7820130,0.9079810)};
\draw[ink!24,line width=.45pt] plot coordinates {(0.0000000,0.0000000,2.0000000) (0.0237682,0.0466477,1.9993146) (0.0475201,0.0932634,1.9972591) (0.0712394,0.1398151,1.9938347) (0.0949099,0.1862711,1.9890438) (0.1185153,0.2325994,1.9828897) (0.1420395,0.2787683,1.9753767) (0.1654664,0.3247461,1.9665098) (0.1887799,0.3705013,1.9562952) (0.2119640,0.4160027,1.9447398) (0.2350028,0.4612189,1.9318517) (0.2578805,0.5061191,1.9176395) (0.2805816,0.5506723,1.9021130) (0.3030903,0.5948482,1.8852830) (0.3253913,0.6386164,1.8671609) (0.3474693,0.6819469,1.8477591) (0.3693091,0.7248100,1.8270909) (0.3908959,0.7671764,1.8051706) (0.4122147,0.8090170,1.7820130) (0.4332511,0.8503031,1.7576342) (0.4539905,0.8910065,1.7320508) (0.4744188,0.9310993,1.7052803) (0.4945219,0.9705539,1.6773411) (0.5142861,1.0093433,1.6482524) (0.5336978,1.0474410,1.6180340) (0.5527438,1.0848208,1.5867067) (0.5714110,1.1214571,1.5542919) (0.5896865,1.1573249,1.5208119) (0.6075579,1.1923995,1.4862897) (0.6250129,1.2266568,1.4507487) (0.6420395,1.2600735,1.4142136) (0.6586261,1.2926266,1.3767092) (0.6747614,1.3242938,1.3382612) (0.6904342,1.3550534,1.2988961) (0.7056338,1.3848842,1.2586408) (0.7203498,1.4137660,1.2175229) (0.7345721,1.4416788,1.1755705) (0.7482909,1.4686036,1.1328125) (0.7614969,1.4945219,1.0892781) (0.7741811,1.5194159,1.0449971) (0.7863346,1.5432686,1.0000000) (0.7979492,1.5660636,0.9543175) (0.8090170,1.5877853,0.9079810) (0.8195303,1.6084188,0.8610222) (0.8294819,1.6279499,0.8134733) (0.8388651,1.6463654,0.7653669) (0.8476733,1.6636525,0.7167359) (0.8559006,1.6797994,0.6676137) (0.8635412,1.6947951,0.6180340) (0.8705901,1.7086293,0.5680307) (0.8770423,1.7212924,0.5176381) (0.8828934,1.7327759,0.4668907) (0.8881394,1.7430718,0.4158234) (0.8927768,1.7521731,0.3644711) (0.8968022,1.7600735,0.3128689) (0.9002131,1.7667677,0.2610524) (0.9030070,1.7722510,0.2090569) (0.9051820,1.7765197,0.1569182) (0.9067366,1.7795709,0.1046719) (0.9076699,1.7814024,0.0523539) (0.9079810,1.7820130,0.0000000)};
\draw[ink!24,line width=.45pt] plot coordinates {(1.9021130,0.0000000,0.6180340) (1.9014612,0.0497915,0.6180340) (1.8995063,0.0995489,0.6180340) (1.8962495,0.1492381,0.6180340) (1.8916931,0.1988250,0.6180340) (1.8858402,0.2482756,0.6180340) (1.8786949,0.2975560,0.6180340) (1.8702620,0.3466326,0.6180340) (1.8605473,0.3954715,0.6180340) (1.8495575,0.4440395,0.6180340) (1.8373001,0.4923031,0.6180340) (1.8237835,0.5402293,0.6180340) (1.8090170,0.5877853,0.6180340) (1.7930107,0.6349384,0.6180340) (1.7757755,0.6816563,0.6180340) (1.7573233,0.7279071,0.6180340) (1.7376667,0.7736591,0.6180340) (1.7168192,0.8188808,0.6180340) (1.6947951,0.8635412,0.6180340) (1.6716095,0.9076099,0.6180340) (1.6472782,0.9510565,0.6180340) (1.6218180,0.9938513,0.6180340) (1.5952462,1.0359650,0.6180340) (1.5675812,1.0773687,0.6180340) (1.5388418,1.1180340,0.6180340) (1.5090477,1.1579330,0.6180340) (1.4782195,1.1970385,0.6180340) (1.4463781,1.2353236,0.6180340) (1.4135455,1.2727620,0.6180340) (1.3797440,1.3093282,0.6180340) (1.3449970,1.3449970,0.6180340) (1.3093282,1.3797440,0.6180340) (1.2727620,1.4135455,0.6180340) (1.2353236,1.4463781,0.6180340) (1.1970385,1.4782195,0.6180340) (1.1579330,1.5090477,0.6180340) (1.1180340,1.5388418,0.6180340) (1.0773687,1.5675812,0.6180340) (1.0359650,1.5952462,0.6180340) (0.9938513,1.6218180,0.6180340) (0.9510565,1.6472782,0.6180340) (0.9076099,1.6716095,0.6180340) (0.8635412,1.6947951,0.6180340) (0.8188808,1.7168192,0.6180340) (0.7736591,1.7376667,0.6180340) (0.7279071,1.7573233,0.6180340) (0.6816563,1.7757755,0.6180340) (0.6349384,1.7930107,0.6180340) (0.5877853,1.8090170,0.6180340) (0.5402293,1.8237835,0.6180340) (0.4923031,1.8373001,0.6180340) (0.4440395,1.8495575,0.6180340) (0.3954715,1.8605473,0.6180340) (0.3466326,1.8702620,0.6180340) (0.2975560,1.8786949,0.6180340) (0.2482756,1.8858402,0.6180340) (0.1988250,1.8916931,0.6180340) (0.1492381,1.8962495,0.6180340) (0.0995489,1.8995063,0.6180340) (0.0497915,1.9014612,0.6180340) (0.0000000,1.9021130,0.6180340)};
\draw[ink!24,line width=.45pt] plot coordinates {(0.0000000,0.0000000,2.0000000) (0.0161782,0.0497915,1.9993146) (0.0323454,0.0995489,1.9972591) (0.0484904,0.1492381,1.9938347) (0.0646021,0.1988250,1.9890438) (0.0806696,0.2482756,1.9828897) (0.0966818,0.2975560,1.9753767) (0.1126277,0.3466326,1.9665098) (0.1284965,0.3954715,1.9562952) (0.1442772,0.4440395,1.9447398) (0.1599590,0.4923031,1.9318517) (0.1755311,0.5402293,1.9176395) (0.1909830,0.5877853,1.9021130) (0.2063040,0.6349384,1.8852830) (0.2214836,0.6816563,1.8671609) (0.2365114,0.7279071,1.8477591) (0.2513771,0.7736591,1.8270909) (0.2660705,0.8188808,1.8051706) (0.2805816,0.8635412,1.7820130) (0.2949003,0.9076099,1.7576342) (0.3090170,0.9510565,1.7320508) (0.3229219,0.9938513,1.7052803) (0.3366054,1.0359650,1.6773411) (0.3500583,1.0773687,1.6482524) (0.3632713,1.1180340,1.6180340) (0.3762353,1.1579330,1.5867067) (0.3889414,1.1970385,1.5542919) (0.4013810,1.2353236,1.5208119) (0.4135455,1.2727620,1.4862897) (0.4254265,1.3093282,1.4507487) (0.4370160,1.3449970,1.4142136) (0.4483060,1.3797440,1.3767092) (0.4592888,1.4135455,1.3382612) (0.4699567,1.4463781,1.2988961) (0.4803026,1.4782195,1.2586408) (0.4903193,1.5090477,1.2175229) (0.5000000,1.5388418,1.1755705) (0.5093380,1.5675812,1.1328125) (0.5183269,1.5952462,1.0892781) (0.5269606,1.6218180,1.0449971) (0.5352331,1.6472782,1.0000000) (0.5431388,1.6716095,0.9543175) (0.5506723,1.6947951,0.9079810) (0.5578284,1.7168192,0.8610222) (0.5646021,1.7376667,0.8134733) (0.5709890,1.7573233,0.7653669) (0.5769844,1.7757755,0.7167359) (0.5825845,1.7930107,0.6676137) (0.5877853,1.8090170,0.6180340) (0.5925832,1.8237835,0.5680307) (0.5969750,1.8373001,0.5176381) (0.6009577,1.8495575,0.4668907) (0.6045285,1.8605473,0.4158234) (0.6076850,1.8702620,0.3644711) (0.6104250,1.8786949,0.3128689) (0.6127466,1.8858402,0.2610524) (0.6146483,1.8916931,0.2090569) (0.6161288,1.8962495,0.1569182) (0.6171870,1.8995063,0.1046719) (0.6178222,1.9014612,0.0523539) (0.6180340,1.9021130,0.0000000)};
\draw[ink!24,line width=.45pt] plot coordinates {(1.9753767,0.0000000,0.3128689) (1.9746998,0.0517093,0.3128689) (1.9726695,0.1033832,0.3128689) (1.9692873,0.1549863,0.3128689) (1.9645554,0.2064831,0.3128689) (1.9584771,0.2578384,0.3128689) (1.9510565,0.3090170,0.3128689) (1.9422988,0.3599838,0.3128689) (1.9322100,0.4107039,0.3128689) (1.9207969,0.4611425,0.3128689) (1.9080674,0.5112651,0.3128689) (1.8940301,0.5610373,0.3128689) (1.8786949,0.6104250,0.3128689) (1.8620720,0.6593943,0.3128689) (1.8441730,0.7079117,0.3128689) (1.8250101,0.7559439,0.3128689) (1.8045964,0.8034581,0.3128689) (1.7829459,0.8504216,0.3128689) (1.7600735,0.8968022,0.3128689) (1.7359948,0.9425683,0.3128689) (1.7107264,0.9876883,0.3128689) (1.6842855,1.0321315,0.3128689) (1.6566903,1.0758672,0.3128689) (1.6279597,1.1188657,0.3128689) (1.5981133,1.1610973,0.3128689) (1.5671717,1.2025331,0.3128689) (1.5351560,1.2431448,0.3128689) (1.5020882,1.2829045,0.3128689) (1.4679910,1.3217850,0.3128689) (1.4328876,1.3597596,0.3128689) (1.3968022,1.3968022,0.3128689) (1.3597596,1.4328876,0.3128689) (1.3217850,1.4679910,0.3128689) (1.2829045,1.5020882,0.3128689) (1.2431448,1.5351560,0.3128689) (1.2025331,1.5671717,0.3128689) (1.1610973,1.5981133,0.3128689) (1.1188657,1.6279597,0.3128689) (1.0758672,1.6566903,0.3128689) (1.0321315,1.6842855,0.3128689) (0.9876883,1.7107264,0.3128689) (0.9425683,1.7359948,0.3128689) (0.8968022,1.7600735,0.3128689) (0.8504216,1.7829459,0.3128689) (0.8034581,1.8045964,0.3128689) (0.7559439,1.8250101,0.3128689) (0.7079117,1.8441730,0.3128689) (0.6593943,1.8620720,0.3128689) (0.6104250,1.8786949,0.3128689) (0.5610373,1.8940301,0.3128689) (0.5112651,1.9080674,0.3128689) (0.4611425,1.9207969,0.3128689) (0.4107039,1.9322100,0.3128689) (0.3599838,1.9422988,0.3128689) (0.3090170,1.9510565,0.3128689) (0.2578384,1.9584771,0.3128689) (0.2064831,1.9645554,0.3128689) (0.1549863,1.9692873,0.3128689) (0.1033832,1.9726695,0.3128689) (0.0517093,1.9746998,0.3128689) (0.0000000,1.9753767,0.3128689)};
\draw[ink!24,line width=.45pt] plot coordinates {(0.0000000,0.0000000,2.0000000) (0.0081900,0.0517093,1.9993146) (0.0163743,0.1033832,1.9972591) (0.0245474,0.1549863,1.9938347) (0.0327037,0.2064831,1.9890438) (0.0408376,0.2578384,1.9828897) (0.0489435,0.3090170,1.9753767) (0.0570158,0.3599838,1.9665098) (0.0650491,0.4107039,1.9562952) (0.0730378,0.4611425,1.9447398) (0.0809764,0.5112651,1.9318517) (0.0888596,0.5610373,1.9176395) (0.0966818,0.6104250,1.9021130) (0.1044378,0.6593943,1.8852830) (0.1121222,0.7079117,1.8671609) (0.1197298,0.7559439,1.8477591) (0.1272553,0.8034581,1.8270909) (0.1346935,0.8504216,1.8051706) (0.1420395,0.8968022,1.7820130) (0.1492882,0.9425683,1.7576342) (0.1564345,0.9876883,1.7320508) (0.1634736,1.0321315,1.7052803) (0.1704006,1.0758672,1.6773411) (0.1772109,1.1188657,1.6482524) (0.1838997,1.1610973,1.6180340) (0.1904625,1.2025331,1.5867067) (0.1968948,1.2431448,1.5542919) (0.2031921,1.2829045,1.5208119) (0.2093502,1.3217850,1.4862897) (0.2153648,1.3597596,1.4507487) (0.2212317,1.3968022,1.4142136) (0.2269471,1.4328876,1.3767092) (0.2325069,1.4679910,1.3382612) (0.2379074,1.5020882,1.2988961) (0.2431448,1.5351560,1.2586408) (0.2482156,1.5671717,1.2175229) (0.2531163,1.5981133,1.1755705) (0.2578435,1.6279597,1.1328125) (0.2623940,1.6566903,1.0892781) (0.2667646,1.6842855,1.0449971) (0.2709524,1.7107264,1.0000000) (0.2749546,1.7359948,0.9543175) (0.2787683,1.7600735,0.9079810) (0.2823909,1.7829459,0.8610222) (0.2858200,1.8045964,0.8134733) (0.2890532,1.8250101,0.7653669) (0.2920883,1.8441730,0.7167359) (0.2949232,1.8620720,0.6676137) (0.2975560,1.8786949,0.6180340) (0.2999849,1.8940301,0.5680307) (0.3022082,1.9080674,0.5176381) (0.3042243,1.9207969,0.4668907) (0.3060320,1.9322100,0.4158234) (0.3076299,1.9422988,0.3644711) (0.3090170,1.9510565,0.3128689) (0.3101923,1.9584771,0.2610524) (0.3111550,1.9645554,0.2090569) (0.3119045,1.9692873,0.1569182) (0.3124402,1.9726695,0.1046719) (0.3127617,1.9746998,0.0523539) (0.3128689,1.9753767,0.0000000)};
\draw[ink!65,line width=.75pt] plot coordinates {(0.0000000,0.0000000,2.0000000) (0.0523539,0.0000000,1.9993146) (0.1046719,0.0000000,1.9972591) (0.1569182,0.0000000,1.9938347) (0.2090569,0.0000000,1.9890438) (0.2610524,0.0000000,1.9828897) (0.3128689,0.0000000,1.9753767) (0.3644711,0.0000000,1.9665098) (0.4158234,0.0000000,1.9562952) (0.4668907,0.0000000,1.9447398) (0.5176381,0.0000000,1.9318517) (0.5680307,0.0000000,1.9176395) (0.6180340,0.0000000,1.9021130) (0.6676137,0.0000000,1.8852830) (0.7167359,0.0000000,1.8671609) (0.7653669,0.0000000,1.8477591) (0.8134733,0.0000000,1.8270909) (0.8610222,0.0000000,1.8051706) (0.9079810,0.0000000,1.7820130) (0.9543175,0.0000000,1.7576342) (1.0000000,0.0000000,1.7320508) (1.0449971,0.0000000,1.7052803) (1.0892781,0.0000000,1.6773411) (1.1328125,0.0000000,1.6482524) (1.1755705,0.0000000,1.6180340) (1.2175229,0.0000000,1.5867067) (1.2586408,0.0000000,1.5542919) (1.2988961,0.0000000,1.5208119) (1.3382612,0.0000000,1.4862897) (1.3767092,0.0000000,1.4507487) (1.4142136,0.0000000,1.4142136) (1.4507487,0.0000000,1.3767092) (1.4862897,0.0000000,1.3382612) (1.5208119,0.0000000,1.2988961) (1.5542919,0.0000000,1.2586408) (1.5867067,0.0000000,1.2175229) (1.6180340,0.0000000,1.1755705) (1.6482524,0.0000000,1.1328125) (1.6773411,0.0000000,1.0892781) (1.7052803,0.0000000,1.0449971) (1.7320508,0.0000000,1.0000000) (1.7576342,0.0000000,0.9543175) (1.7820130,0.0000000,0.9079810) (1.8051706,0.0000000,0.8610222) (1.8270909,0.0000000,0.8134733) (1.8477591,0.0000000,0.7653669) (1.8671609,0.0000000,0.7167359) (1.8852830,0.0000000,0.6676137) (1.9021130,0.0000000,0.6180340) (1.9176395,0.0000000,0.5680307) (1.9318517,0.0000000,0.5176381) (1.9447398,0.0000000,0.4668907) (1.9562952,0.0000000,0.4158234) (1.9665098,0.0000000,0.3644711) (1.9753767,0.0000000,0.3128689) (1.9828897,0.0000000,0.2610524) (1.9890438,0.0000000,0.2090569) (1.9938347,0.0000000,0.1569182) (1.9972591,0.0000000,0.1046719) (1.9993146,0.0000000,0.0523539) (2.0000000,0.0000000,0.0000000)};
\draw[ink!65,line width=.75pt] plot coordinates {(0.0000000,0.0000000,2.0000000) (0.0000000,0.0523539,1.9993146) (0.0000000,0.1046719,1.9972591) (0.0000000,0.1569182,1.9938347) (0.0000000,0.2090569,1.9890438) (0.0000000,0.2610524,1.9828897) (0.0000000,0.3128689,1.9753767) (0.0000000,0.3644711,1.9665098) (0.0000000,0.4158234,1.9562952) (0.0000000,0.4668907,1.9447398) (0.0000000,0.5176381,1.9318517) (0.0000000,0.5680307,1.9176395) (0.0000000,0.6180340,1.9021130) (0.0000000,0.6676137,1.8852830) (0.0000000,0.7167359,1.8671609) (0.0000000,0.7653669,1.8477591) (0.0000000,0.8134733,1.8270909) (0.0000000,0.8610222,1.8051706) (0.0000000,0.9079810,1.7820130) (0.0000000,0.9543175,1.7576342) (0.0000000,1.0000000,1.7320508) (0.0000000,1.0449971,1.7052803) (0.0000000,1.0892781,1.6773411) (0.0000000,1.1328125,1.6482524) (0.0000000,1.1755705,1.6180340) (0.0000000,1.2175229,1.5867067) (0.0000000,1.2586408,1.5542919) (0.0000000,1.2988961,1.5208119) (0.0000000,1.3382612,1.4862897) (0.0000000,1.3767092,1.4507487) (0.0000000,1.4142136,1.4142136) (0.0000000,1.4507487,1.3767092) (0.0000000,1.4862897,1.3382612) (0.0000000,1.5208119,1.2988961) (0.0000000,1.5542919,1.2586408) (0.0000000,1.5867067,1.2175229) (0.0000000,1.6180340,1.1755705) (0.0000000,1.6482524,1.1328125) (0.0000000,1.6773411,1.0892781) (0.0000000,1.7052803,1.0449971) (0.0000000,1.7320508,1.0000000) (0.0000000,1.7576342,0.9543175) (0.0000000,1.7820130,0.9079810) (0.0000000,1.8051706,0.8610222) (0.0000000,1.8270909,0.8134733) (0.0000000,1.8477591,0.7653669) (0.0000000,1.8671609,0.7167359) (0.0000000,1.8852830,0.6676137) (0.0000000,1.9021130,0.6180340) (0.0000000,1.9176395,0.5680307) (0.0000000,1.9318517,0.5176381) (0.0000000,1.9447398,0.4668907) (0.0000000,1.9562952,0.4158234) (0.0000000,1.9665098,0.3644711) (0.0000000,1.9753767,0.3128689) (0.0000000,1.9828897,0.2610524) (0.0000000,1.9890438,0.2090569) (0.0000000,1.9938347,0.1569182) (0.0000000,1.9972591,0.1046719) (0.0000000,1.9993146,0.0523539) (0.0000000,2.0000000,0.0000000)};
\draw[ink!65,line width=.75pt] plot coordinates {(2.0000000,0.0000000,0.0000000) (1.9993146,0.0523539,0.0000000) (1.9972591,0.1046719,0.0000000) (1.9938347,0.1569182,0.0000000) (1.9890438,0.2090569,0.0000000) (1.9828897,0.2610524,0.0000000) (1.9753767,0.3128689,0.0000000) (1.9665098,0.3644711,0.0000000) (1.9562952,0.4158234,0.0000000) (1.9447398,0.4668907,0.0000000) (1.9318517,0.5176381,0.0000000) (1.9176395,0.5680307,0.0000000) (1.9021130,0.6180340,0.0000000) (1.8852830,0.6676137,0.0000000) (1.8671609,0.7167359,0.0000000) (1.8477591,0.7653669,0.0000000) (1.8270909,0.8134733,0.0000000) (1.8051706,0.8610222,0.0000000) (1.7820130,0.9079810,0.0000000) (1.7576342,0.9543175,0.0000000) (1.7320508,1.0000000,0.0000000) (1.7052803,1.0449971,0.0000000) (1.6773411,1.0892781,0.0000000) (1.6482524,1.1328125,0.0000000) (1.6180340,1.1755705,0.0000000) (1.5867067,1.2175229,0.0000000) (1.5542919,1.2586408,0.0000000) (1.5208119,1.2988961,0.0000000) (1.4862897,1.3382612,0.0000000) (1.4507487,1.3767092,0.0000000) (1.4142136,1.4142136,0.0000000) (1.3767092,1.4507487,0.0000000) (1.3382612,1.4862897,0.0000000) (1.2988961,1.5208119,0.0000000) (1.2586408,1.5542919,0.0000000) (1.2175229,1.5867067,0.0000000) (1.1755705,1.6180340,0.0000000) (1.1328125,1.6482524,0.0000000) (1.0892781,1.6773411,0.0000000) (1.0449971,1.7052803,0.0000000) (1.0000000,1.7320508,0.0000000) (0.9543175,1.7576342,0.0000000) (0.9079810,1.7820130,0.0000000) (0.8610222,1.8051706,0.0000000) (0.8134733,1.8270909,0.0000000) (0.7653669,1.8477591,0.0000000) (0.7167359,1.8671609,0.0000000) (0.6676137,1.8852830,0.0000000) (0.6180340,1.9021130,0.0000000) (0.5680307,1.9176395,0.0000000) (0.5176381,1.9318517,0.0000000) (0.4668907,1.9447398,0.0000000) (0.4158234,1.9562952,0.0000000) (0.3644711,1.9665098,0.0000000) (0.3128689,1.9753767,0.0000000) (0.2610524,1.9828897,0.0000000) (0.2090569,1.9890438,0.0000000) (0.1569182,1.9938347,0.0000000) (0.1046719,1.9972591,0.0000000) (0.0523539,1.9993146,0.0000000) (0.0000000,2.0000000,0.0000000)};
\draw[teal,line width=1.65pt] plot coordinates {(1.7888544,0.6324555,0.6324555) (1.7832225,0.6367176,0.6439784) (1.7775092,0.6409507,0.6554718) (1.7717147,0.6451544,0.6669352) (1.7658394,0.6493287,0.6783683) (1.7598834,0.6534734,0.6897703) (1.7538470,0.6575882,0.7011409) (1.7477305,0.6616730,0.7124794) (1.7415342,0.6657276,0.7237854) (1.7352584,0.6697517,0.7350583) (1.7289034,0.6737453,0.7462977) (1.7224694,0.6777081,0.7575030) (1.7159568,0.6816400,0.7686737) (1.7093657,0.6855408,0.7798093) (1.7026967,0.6894102,0.7909093) (1.6959499,0.6932481,0.8019732) (1.6891256,0.6970544,0.8130005) (1.6822242,0.7008289,0.8239906) (1.6752460,0.7045714,0.8349431) (1.6681913,0.7082817,0.8458575) (1.6610604,0.7119596,0.8567333) (1.6538536,0.7156051,0.8675699) (1.6465714,0.7192178,0.8783669) (1.6392139,0.7227977,0.8891238) (1.6317816,0.7263446,0.8998401) (1.6242748,0.7298584,0.9105153) (1.6166938,0.7333388,0.9211490) (1.6090390,0.7367857,0.9317405) (1.6013107,0.7401990,0.9422895) (1.5935092,0.7435785,0.9527955) (1.5856350,0.7469240,0.9632580) (1.5776885,0.7502354,0.9736765) (1.5696698,0.7535126,0.9840505) (1.5615795,0.7567553,0.9943796) (1.5534179,0.7599635,1.0046633) (1.5451853,0.7631370,1.0149011) (1.5368822,0.7662756,1.0250926) (1.5285089,0.7693793,1.0352373) (1.5200657,0.7724478,1.0453346) (1.5115532,0.7754810,1.0553843) (1.5029717,0.7784789,1.0653857) (1.4943215,0.7814412,1.0753385) (1.4856031,0.7843677,1.0852422) (1.4768168,0.7872585,1.0950964) (1.4679631,0.7901133,1.1049005) (1.4590424,0.7929321,1.1146542) (1.4500550,0.7957146,1.1243570) (1.4410015,0.7984608,1.1340084) (1.4318821,0.8011706,1.1436080) (1.4226973,0.8038437,1.1531555) (1.4134476,0.8064802,1.1626502) (1.4041333,0.8090798,1.1720919) (1.3947549,0.8116425,1.1814801) (1.3853129,0.8141681,1.1908143) (1.3758075,0.8166565,1.2000941) (1.3662394,0.8191076,1.2093191) (1.3566088,0.8215214,1.2184889) (1.3469163,0.8238976,1.2276031) (1.3371623,0.8262362,1.2366612) (1.3273472,0.8285371,1.2456628) (1.3174715,0.8308001,1.2546075) (1.3075357,0.8330252,1.2634950) (1.2975401,0.8352122,1.2723247) (1.2874853,0.8373611,1.2810964) (1.2773717,0.8394718,1.2898095) (1.2671998,0.8415442,1.2984638) (1.2569700,0.8435781,1.3070587) (1.2466828,0.8455735,1.3155940) (1.2363386,0.8475302,1.3240692) (1.2259380,0.8494483,1.3324840) (1.2154815,0.8513276,1.3408378) (1.2049694,0.8531680,1.3491305) (1.1944023,0.8549695,1.3573615) (1.1837807,0.8567319,1.3655306) (1.1731050,0.8584552,1.3736373) (1.1623757,0.8601392,1.3816813) (1.1515934,0.8617840,1.3896622) (1.1407584,0.8633895,1.3975796) (1.1298714,0.8649555,1.4054332) (1.1189328,0.8664821,1.4132227) (1.1079431,0.8679690,1.4209476) (1.0969027,0.8694163,1.4286076) (1.0858123,0.8708240,1.4362023) (1.0746724,0.8721918,1.4437315) (1.0634833,0.8735199,1.4511948) (1.0522457,0.8748080,1.4585918) (1.0409600,0.8760562,1.4659222) (1.0296268,0.8772644,1.4731856) (1.0182466,0.8784325,1.4803818) (1.0068199,0.8795605,1.4875104) (0.9953472,0.8806484,1.4945710) (0.9838290,0.8816960,1.5015634) (0.9722659,0.8827034,1.5084872) (0.9606585,0.8836704,1.5153422) (0.9490071,0.8845972,1.5221279) (0.9373124,0.8854835,1.5288441) (0.9255750,0.8863294,1.5354906) (0.9137952,0.8871348,1.5420669) (0.9019737,0.8878997,1.5485728) (0.8901111,0.8886240,1.5550079) (0.8782077,0.8893078,1.5613721) (0.8662643,0.8899509,1.5676650) (0.8542814,0.8905535,1.5738862) (0.8422594,0.8911153,1.5800356) (0.8301989,0.8916365,1.5861129) (0.8181006,0.8921169,1.5921177) (0.8059649,0.8925566,1.5980498) (0.7937923,0.8929556,1.6039090) (0.7815836,0.8933138,1.6096949) (0.7693391,0.8936311,1.6154073) (0.7570595,0.8939077,1.6210459) (0.7447454,0.8941435,1.6266105) (0.7323972,0.8943384,1.6321008) (0.7200156,0.8944925,1.6375166) (0.7076011,0.8946057,1.6428577) (0.6951543,0.8946781,1.6481237) (0.6826757,0.8947096,1.6533144) (0.6701660,0.8947003,1.6584296) (0.6576257,0.8946501,1.6634691) (0.6450553,0.8945591,1.6684327) (0.6324555,0.8944272,1.6733201)};
\fill[copper] (1.7888544,0.6324555,0.6324555) circle(2.2pt);
\fill[copper] (0.6324555,0.8944272,1.6733201) circle(2.2pt);
\end{tikzpicture}\caption{An orthographic view of the positive octant of the radius-2 sphere. The teal curve is the great-circle image of a categorical Fisher geodesic; orange points mark its endpoint laws.}\end{figure}

\section{Why shortest spherical paths are great-circle arcs}
A great circle is the intersection of a sphere with a plane through its center. Two non-antipodal points determine a shorter great-circle arc, with length equal to radius times central angle.

The minimizing property follows from a useful length estimate. Let $c(t)$ lie on a radius-$R$ sphere, let $e_0$ be the unit vector toward its starting point, and define the central angle $\theta(t)=\arccos(e_0\cdot c(t)/R)$. At regular interior angles,
\[
|\theta'|=\frac{|e_0\cdot c'|}{R\sin\theta}\leq\frac{\|c'\|}{R}.
\]
For the inequality, project $e_0$ into the tangent plane at $c(t)$; that projection has norm $\sin\theta$, while $c'$ is tangent. Cauchy--Schwarz then gives the bound. Integrating, and taking endpoint limits when necessary, shows that every connecting curve has length at least $R$ times the endpoint central angle. The shorter great-circle arc attains this lower bound.

For square-root images of categorical laws, the endpoint angle is at most $\pi/2$. Their shorter arc has nonnegative linear-combination coefficients, as the explicit formula below shows. It therefore remains in the spherical orthant. The spherical shortest path is also an admissible path in the statistical family.

\section{The Bhattacharyya coefficient and Fisher distance}\index{Fisher distance!categorical}
On the sphere, the shortest distance between two points is determined by the angle between their radius vectors. The square-root map lets us calculate that angle directly from the original probabilities. Its cosine is a familiar statistical overlap: the sum of the geometric means of corresponding category weights.
For laws $P=(p_i)$ and $R=(r_i)$, define their \emph{Bhattacharyya coefficient}\index{Bhattacharyya coefficient}
\[
B(P,R)=\sum_i\sqrt{p_i r_i}.
\]
Cauchy--Schwarz applied to the unit vectors $\sqrt P$ and $\sqrt R$ gives $B\leq1$, with equality exactly when $P=R$. Nonnegativity gives $B\geq0$. At the boundary $B=0$ exactly when their supports are disjoint.

The square-root image vectors $s=2\sqrt P$ and $r=2\sqrt R$ satisfy
\[
\cos\alpha=\frac{s\cdot r}{4}=B(P,R).
\]
Multiplying the central angle by the radius gives
\begin{equation}\label{eq:sphericaldistance}
\boxed{d_F(P,R)=2\arccos B(P,R).}
\end{equation}
The formula is symmetric, vanishes precisely at equality, and satisfies the triangle inequality because it is a shortest-path distance.

\begin{example}[Computing a categorical distance]
For $P=(0.5,0.3,0.2)$ and $R=(0.2,0.5,0.3)$, express the coefficient and distance.
\begin{solution}
The coefficient is
\[
B=\sqrt{0.10}+\sqrt{0.15}+\sqrt{0.06}.
\]
Numerically $B\approx0.948475$, so $\alpha\approx0.322408$ radians and $d_F\approx0.644817$. Each term compares the mass assigned to the same category in the two laws; the sum is their square-root-vector overlap.
\end{solution}
\end{example}

\section{Constructing the constant-speed geodesic}
Let $P\ne R$ be positive laws and $\alpha=\arccos B(P,R)$. Their spherical interpolation is
\begin{equation}\label{eq:slerp}
s(t)=\frac{\sin((1-t)\alpha)}{\sin\alpha}s(0)
+\frac{\sin(t\alpha)}{\sin\alpha}s(1),\qquad 0\leq t\leq1,
\end{equation}
where $s(0)=2\sqrt P$ and $s(1)=2\sqrt R$. The probability path is
\begin{equation}\label{eq:sphericalpath}
\gamma_i(t)=\frac{s_i(t)^2}{4}.
\end{equation}
It is a \emph{geodesic}\index{geodesic!categorical}: a shortest path in the Fisher geometry, parameterized with constant speed.

To verify its properties, choose orthonormal vectors
\[
e_0=\frac{s(0)}2,\qquad
e_1=\frac{s(1)/2-\cos\alpha\,e_0}{\sin\alpha}.
\]
The interpolation becomes
\[
s(t)=2\bigl[\cos(t\alpha)e_0+\sin(t\alpha)e_1\bigr].
\]
Thus its squared norm is 4, its speed is $2\alpha$, and its length over $[0,1]$ is $2\alpha$. The sine coefficients in \eqref{eq:slerp} are positive between the endpoints, so every coordinate stays positive. If $P=R$, the appropriate geodesic is the constant path.

\begin{example}[The Fisher midpoint]
Find the probability midpoint $\gamma(1/2)$ between two positive laws.
\begin{solution}
At the midpoint both coefficients in \eqref{eq:slerp} equal $1/[2\cos(\alpha/2)]$. Consequently
\begin{equation}\label{eq:fishermidpoint}
\gamma_i(1/2)=\frac{(\sqrt{p_i}+\sqrt{r_i})^2}{2(1+B(P,R))}.
\end{equation}
The numerators sum to $2+2B$, confirming normalization. This midpoint generally differs from the ordinary average $(P+R)/2$.
\end{solution}
\end{example}

\section{Mixture paths and geometric paths}
The mixture $M(t)=(1-t)P+tR$ traces a straight segment in probability coordinates. Its Fisher length is
\begin{equation}\label{eq:mixturelength}
L_F(M)=\int_0^1\sqrt{\sum_i\frac{(r_i-p_i)^2}{(1-t)p_i+tr_i}}\dd t.
\end{equation}
The shortest-path property implies $L_F(M)\geq d_F(P,R)$. Equality holds when the mixture traces the same minimizing arc, with no backtracking. For general three-outcome endpoints its spherical image bends away from the great circle, and the inequality is strict.

\begin{figure}[htbp]\centering\begin{tikzpicture}[scale=6]
\fill[soft] (0,0)--(1,0)--(.5,.8660254)--cycle;
\draw[ink!75] (0,0)--(1,0)--(.5,.8660254)--cycle;
\node[below left,font=\small] at (0,0){$A$};\node[below right,font=\small] at (1,0){$B$};\node[above,font=\small] at (.5,.8660254){$C$};
\draw[copper,dashed,line width=1.1pt] plot coordinates {(0.1500000,0.0866025) (0.5500000,0.6062178)};
\draw[teal,line width=1.25pt] plot coordinates {(0.1500000,0.0866025) (0.1531909,0.0897869) (0.1564098,0.0930205) (0.1596564,0.0963026) (0.1629299,0.0996327) (0.1662298,0.1030101) (0.1695554,0.1064342) (0.1729062,0.1099044) (0.1762815,0.1134202) (0.1796807,0.1169807) (0.1831032,0.1205854) (0.1865484,0.1242337) (0.1900157,0.1279248) (0.1935044,0.1316581) (0.1970138,0.1354329) (0.2005434,0.1392485) (0.2040925,0.1431042) (0.2076604,0.1469993) (0.2112465,0.1509331) (0.2148501,0.1549049) (0.2184706,0.1589139) (0.2221073,0.1629595) (0.2257596,0.1670408) (0.2294268,0.1711572) (0.2331082,0.1753078) (0.2368031,0.1794920) (0.2405109,0.1837090) (0.2442309,0.1879579) (0.2479623,0.1922381) (0.2517047,0.1965487) (0.2554571,0.2008889) (0.2592190,0.2052580) (0.2629897,0.2096551) (0.2667685,0.2140795) (0.2705547,0.2185303) (0.2743476,0.2230068) (0.2781464,0.2275081) (0.2819506,0.2320334) (0.2857595,0.2365818) (0.2895722,0.2411526) (0.2933882,0.2457448) (0.2972067,0.2503578) (0.3010270,0.2549905) (0.3048485,0.2596422) (0.3086704,0.2643121) (0.3124921,0.2689992) (0.3163128,0.2737027) (0.3201318,0.2784218) (0.3239485,0.2831555) (0.3277621,0.2879031) (0.3315720,0.2926637) (0.3353775,0.2974363) (0.3391778,0.3022202) (0.3429722,0.3070144) (0.3467602,0.3118180) (0.3505409,0.3166303) (0.3543137,0.3214503) (0.3580780,0.3262771) (0.3618329,0.3311098) (0.3655779,0.3359477) (0.3693122,0.3407897) (0.3730352,0.3456350) (0.3767461,0.3504827) (0.3804444,0.3553319) (0.3841293,0.3601818) (0.3878002,0.3650315) (0.3914563,0.3698800) (0.3950971,0.3747265) (0.3987218,0.3795701) (0.4023298,0.3844099) (0.4059204,0.3892451) (0.4094931,0.3940747) (0.4130470,0.3988979) (0.4165816,0.4037137) (0.4200962,0.4085214) (0.4235903,0.4133200) (0.4270631,0.4181086) (0.4305140,0.4228864) (0.4339423,0.4276526) (0.4373476,0.4324061) (0.4407290,0.4371462) (0.4440861,0.4418720) (0.4474182,0.4465827) (0.4507247,0.4512773) (0.4540050,0.4559551) (0.4572585,0.4606151) (0.4604846,0.4652565) (0.4636827,0.4698785) (0.4668522,0.4744802) (0.4699926,0.4790608) (0.4731032,0.4836194) (0.4761835,0.4881553) (0.4792330,0.4926675) (0.4822511,0.4971553) (0.4852372,0.5016178) (0.4881908,0.5060542) (0.4911113,0.5104638) (0.4939983,0.5148457) (0.4968512,0.5191990) (0.4996694,0.5235231) (0.5024524,0.5278171) (0.5051998,0.5320803) (0.5079111,0.5363118) (0.5105857,0.5405108) (0.5132232,0.5446767) (0.5158230,0.5488087) (0.5183847,0.5529060) (0.5209079,0.5569678) (0.5233921,0.5609934) (0.5258367,0.5649821) (0.5282415,0.5689332) (0.5306059,0.5728459) (0.5329294,0.5767195) (0.5352118,0.5805533) (0.5374525,0.5843466) (0.5396512,0.5880987) (0.5418074,0.5918090) (0.5439208,0.5954767) (0.5459909,0.5991011) (0.5480174,0.6026817) (0.5500000,0.6062178)};
\fill[ink] (0.1500000,0.0866025) circle(.009);\node[below,font=\small] at (0.1500000,0.0866025){$P$};
\fill[ink] (0.5500000,0.6062178) circle(.009);\node[above,font=\small] at (0.5500000,0.6062178){$R$};
\node[font=\small,copper] at (.8,.16){mixture};\node[font=\small,teal] at (.18,.43){geodesic};
\end{tikzpicture}\caption{The straight mixture and Fisher geodesic from $P=(.8,.1,.1)$ to $R=(.1,.2,.7)$ in barycentric coordinates. Their square-root images trace different spherical routes.}\end{figure}

A useful special case fixes the conditional proportions of a group of categories. For example, let
\[
P=(a,(1-a)r,(1-a)(1-r)),
\]
and let $R$ have the same form with $a$ replaced by $b$, keeping $r$ fixed. Their square-root vectors lie in the plane spanned by $(1,0,0)$ and $(0,\sqrt r,\sqrt{1-r})$. Varying only the first probability traces a great-circle arc in that plane. A monotone mixture therefore has the geodesic's image, although its speed in spherical angle is generally variable.

Path image and parameterization are separate features. Two parameterizations can trace the same shortest route at different speeds. A constant-speed geodesic fixes both the route and its speed over the chosen parameter interval.

\section{Bernoulli angles and arc coordinates}
For $P_p=(p,1-p)$, the image is
\[
s(p)=(2\sqrt p,2\sqrt{1-p}).
\]
Writing $p=\sin^2\varphi$ gives $s=(2\sin\varphi,2\cos\varphi)$, where $0<\varphi<\pi/2$. The ordinary central angle measured from the second coordinate axis is
\[
\varphi=\arcsin\sqrt p.
\]
Since the radius is 2, the arc coordinate is $\theta=2\varphi=2\arcsin\sqrt p$. This is the length coordinate of Chapter 7. The central angle and arc length differ by the radius factor.

\begin{example}[A Bernoulli midpoint]
Find the constant-speed midpoint between $p=1/4$ and $p=3/4$.
\begin{solution}
The central angles are $\pi/6$ and $\pi/3$. Their midpoint is $\pi/4$, so the midpoint probability is $\sin^2(\pi/4)=1/2$. The total arc length is $2(\pi/3-\pi/6)=\pi/3$.
\end{solution}
\end{example}

Every monotone Bernoulli probability path has the same quarter-circle route between its endpoints. This one-dimensional case explains why a linear probability mixture can have minimizing length while failing to have constant Fisher speed.

\section{Chord length and Hellinger distance}
The straight Euclidean chord between square-root vectors has length
\begin{equation}\label{eq:chord}
\|2\sqrt P-2\sqrt R\|^2
=4\sum_i(\sqrt{p_i}-\sqrt{r_i})^2=8(1-B).
\end{equation}
The \emph{Hellinger distance}\index{Hellinger distance}, in a common normalization, is
\[
h(P,R)=\frac1{\sqrt2}\|\sqrt P-\sqrt R\|,
\qquad h(P,R)^2=1-B(P,R).
\]
It is a metric because it is a scaled Euclidean distance under an injective map. The chord is $2\sqrt2\,h$, while the intrinsic Fisher distance is $2\arccos(1-h^2)$.

\begin{figure}[htbp]
\centering
\begin{tikzpicture}[scale=1.45,>=Stealth]
\draw[ink!50] (0,0)--(2.15,0);\draw[ink!50] (0,0)--(0,2.15);
\draw[ink!45] (2,0) arc(0:90:2);
\coordinate (P) at ({2*cos(15)},{2*sin(15)});\coordinate (R) at ({2*cos(65)},{2*sin(65)});
\draw[ink!40] (0,0)--(P);\draw[ink!40] (0,0)--(R);
\draw[copper,thick] (P)--(R);
\draw[teal,line width=1.6pt] (P) arc(15:65:2);
\fill[ink] (P) circle(1.5pt);\fill[ink] (R) circle(1.5pt);
\node[right] at (P){$s_P$};\node[above] at (R){$s_R$};
\node[font=\small,anchor=west] at (1.72,1.45){arc: $2\alpha$};
\node[font=\small,anchor=east] at (1.2,.9){chord};
\draw[->] (.47,.126) arc(15:65:.49);\node[font=\small] at (.52,.44){$\alpha$};
\end{tikzpicture}
\caption{Arc and chord lengths on the radius-2 circle. The arc stays within the statistical manifold; the interior of the chord lies inside the sphere.}
\end{figure}

For small central angle, the chord $4\sin(\alpha/2)$ and arc $2\alpha$ agree to first order. Expanding square roots gives $h(P,P+v)^2=\tfrac18\sum_i v_i^2/p_i+O(\|v\|^3)$, consistent with the same local Fisher geometry after accounting for this normalization.

\section{Initial velocity and local consistency}
The explicit geodesic formula also specifies the initial direction in probability coordinates. Differentiate \eqref{eq:slerp} at $t=0$, then use $\gamma_i'=s_i s_i'/2$. For distinct endpoints,
\begin{equation}\label{eq:geodesicinitial}
\gamma_i'(0)=\frac{2\alpha}{\sin\alpha}
\left(\sqrt{p_ir_i}-B(P,R)p_i\right).
\end{equation}
The components sum to zero because both summed terms equal $B$. The direction is a tangent vector based at $P$, and its Fisher norm is $2\alpha$. Thus an endpoint pair determines both a shortest route and the exact initial velocity for traversing it over a unit interval.

For nearby laws, the global formula recovers the local Fisher metric. Set $R=P+h$ with $\sum_i h_i=0$. Expanding the square root in each overlap term gives
\begin{align*}
B(P,P+h)
&=\sum_i p_i\sqrt{1+h_i/p_i}\\
&=1-\frac18\sum_i\frac{h_i^2}{p_i}+O(\|h\|^3).
\end{align*}
Since $\cos\alpha=1-\alpha^2/2+O(\alpha^4)$, the squared distance satisfies
\[
d_F(P,P+h)^2=\sum_i\frac{h_i^2}{p_i}+O(\|h\|^3).
\]
The arc, chord, and local metric calculations therefore agree to their appropriate leading orders. Their global differences arise as the displacement grows.

\begin{example}[A geodesic with a persistent category symmetry]
Connect $P=(1/3,1/3,1/3)$ to $R=(1/2,1/4,1/4)$ with constant Fisher speed.
\begin{solution}
Both endpoints have equal second and third coordinates, so their square-root vectors lie in the plane $s_2=s_3$. The connecting great circle remains in that plane. Every intermediate law has the form $(a,(1-a)/2,(1-a)/2)$, a Bernoulli split between A and the combined BC group.

Let $\varphi_0=\arcsin\sqrt{1/3}$ and $\varphi_1=\arcsin\sqrt{1/2}$. The constant-speed path is
\[
a(t)=\sin^2\bigl((1-t)\varphi_0+t\varphi_1\bigr),\qquad
\gamma(t)=\left(a(t),\frac{1-a(t)}2,\frac{1-a(t)}2\right).
\]
Its total length is $2(\varphi_1-\varphi_0)$. Grouping the two symmetric categories reduces this calculation to the one-dimensional circular geometry without changing the Fisher length.
\end{solution}
\end{example}

The probability coordinates of a constant-speed geodesic generally change at nonconstant rates. It is the angle on the sphere, rather than each probability separately, that advances uniformly. The construction specifies a probability model at each intermediate position and therefore permits the numerical readings, entropy, or likelihood of that model to be evaluated along the path.

\section{Boundary completion and curvature}\index{curvature!positive}\index{spherical geometry}
The positive spherical orthant has boundary points at finite distance. Adding its nonnegative boundary produces the \emph{metric completion}\index{metric completion}: every Cauchy sequence then has a limit. The square-root inverse identifies these added points with laws having zero probabilities.

\begin{example}[Two pure distributions]
Take $P=(1,0,0)$ and $R=(0,1,0)$ in the completion.
\begin{solution}
Their square-root vectors $(2,0,0)$ and $(0,2,0)$ have central angle $\pi/2$. Hence $d_F(P,R)=\pi$. Their ordinary probability-vector distance is $\sqrt2$, their square-root chord length is $2\sqrt2$, and both directed KL divergences are infinite. The Fisher distance remains finite because a finite-length spherical arc connects the limiting states.
\end{solution}
\end{example}

For three outcomes the image is an open patch of a 2-sphere of radius 2. A round radius-$R$ sphere has Gaussian curvature $1/R^2$, so the categorical Fisher surface has constant curvature $1/4$. For $n\geq3$, the corresponding higher-dimensional sphere has the same sectional curvature in every two-dimensional tangent plane. A Bernoulli family is one-dimensional; it has arc length geometry, and Gaussian curvature is not defined for it.

Shortest paths answer a geometric question about prescribed endpoints. The next chapter uses the same metric to define motion driven by a scalar objective.

\FloatBarrier
\section{Exercises}
\begin{exercises}
\item Prove the square-root map preserves tangent inner products and maps the normalization tangent space to the spherical tangent space.
\item Map $P=(1/4,1/4,1/2)$ to $S^2(2)$ and verify the norm directly.
\item For $P=(0.5,0.3,0.2)$ and $R=(0.2,0.5,0.3)$, compute $B$, the central angle, Fisher distance, and chord length numerically.
\item Prove $0\leq B(P,R)\leq1$ and characterize equality at both ends.
\item Derive the geodesic midpoint formula \eqref{eq:fishermidpoint}. Verify its coordinates sum to one.
\item Use the orthonormal representation to prove that \eqref{eq:slerp} has constant speed $2\alpha$.
\item For Bernoulli probabilities $1/4$ and $3/4$, compute the central-angle separation, Fisher distance, and square-root chord length.
\item Explain why the shorter spherical arc stays in the positive orthant when both endpoints are positive.
\item Derive the mixture-length formula \eqref{eq:mixturelength}. Explain why its value is at least the Fisher distance.
\item Verify that the fixed-conditional-proportion family in the text lies on a great circle. Explain why its monotone mixture paths have minimizing length.
\item Compute the geodesic from $(1,0,0)$ to $(0,1,0)$ in the metric completion, including its probability parameterization.
\item Derive $h^2=1-B$, the chord formula, and $d_F=2\arccos(1-h^2)$.
\item Show that the distance between any two categorical laws in the completion is at most $\pi$. Characterize pairs attaining this bound.
\item Explain why divergence of a coordinate metric coefficient does not by itself imply infinite boundary distance. Use the Bernoulli integral as evidence.
\item \emph{Exploration.} Compare the mixture and Fisher paths from $(0.8,0.1,0.1)$ to $(0.1,0.2,0.7)$. Numerically integrate the mixture length and compare midpoints.
\item Expand $(\sqrt{p_i+v_i}-\sqrt{p_i})^2$ through quadratic order and obtain the local Hellinger approximation.
\item Prove directly from the spherical distance formula that simultaneous permutation of the outcome labels preserves Fisher distance.
\item \emph{Exploration.} Move two distributions toward different vertices. Track $B$, chord length, Fisher distance, and both KL divergences. Identify which limits depend on support.
\end{exercises}

\chapter{Natural Gradients and Potential-Driven Evolution}
\label{ch:natural}

A scalar objective assigns a value to every statistical state. Its differential specifies how that value changes along a velocity. A metric converts the differential into a gradient vector, selecting the direction of steepest increase for the metric's definition of unit length. With the Fisher metric, the resulting natural gradient has the replicator form.

The familiar coexistence equation from Chapter 1 provides a guide. Its equilibrium attracted nearby proportions, and later chapters described that motion using entropy and Fisher length. We can now ask which scalar function increases as the state moves and how the metric converts that function into the observed velocity. The answer connects the geometry to the original dynamics.

\section{The metric definition of a gradient}\index{gradient flow}
Let $\Phi$ be a differentiable function on a domain with metric $g$. Its differential $d\Phi_p$ is the linear function on velocities that gives directional derivatives. The metric gradient is the unique vector satisfying
\begin{equation}\label{eq:gradientdefinition}
g_p(\grad_g\Phi,v)=d\Phi_p(v)
\quad\text{for every tangent vector }v.
\end{equation}
Uniqueness follows from positive definiteness. In coordinates, if $d\Phi(v)=(\nabla\Phi)^Tv$ and $g(u,v)=u^TGv$, then
\begin{equation}\label{eq:metricgradient}
\grad_g\Phi=G^{-1}\nabla\Phi.
\end{equation}
The coordinate derivatives form a covector, a linear rule for evaluating velocities. The inverse metric turns this rule into a tangent vector.

Cauchy--Schwarz shows that among all vectors with $g(v,v)=1$, the largest value of $d\Phi(v)$ is $\|\grad_g\Phi\|_g$. It is attained in the normalized gradient direction when the gradient is nonzero. Thus the defining identity recovers the geometric meaning of steepest ascent.

For the Fisher metric the gradient is called the \emph{natural gradient}\index{natural gradient}. The associated continuous ascent equation is $p'=\grad_F\Phi$. Optimization of a loss function instead uses the negative natural gradient.

\section{The Bernoulli gradient and the original equation}
In one dimension, the metric is $ds^2=g(x)dx^2$, so \eqref{eq:gradientdefinition} gives gradient component $\Phi'(x)/g(x)$. The Bernoulli metric therefore gives
\begin{equation}\label{eq:binarynatural}
(\grad_F\Phi)^x=x(1-x)\Phi'(x).
\end{equation}
Every smooth scalar advantage $\Delta(x)$ has an antiderivative $\Phi(x)=\int^x\Delta(s)ds$. The corresponding binary replicator equation is its Fisher-gradient ascent.

\begin{example}[A potential for stable coexistence]
Take $\Phi(x)=3x-2x^2$.
\begin{solution}
Its derivative is $3-4x$. Hence natural-gradient ascent is
\[
x'=x(1-x)(3-4x),
\]
the recurring coexistence equation. The potential has its unique interior maximum at $x=3/4$, with value $9/8$. Along the flow,
\[
\dot\Phi=x(1-x)(3-4x)^2\geq0.
\]
The endpoints are equilibria of the smooth extended replicator field; the gradient is defined on the positive interval.
\end{solution}
\end{example}

\begin{figure}[htbp]
\centering
\begin{tikzpicture}
\begin{axis}[width=.86\textwidth,height=5cm,xmin=0,xmax=1,ymin=0,ymax=1.25,xlabel={probability $x$},ylabel={potential $\Phi(x)$},grid=major,xtick={0,.25,.5,.75,1}]
\addplot[teal,domain=0:1,samples=80]{3*x-2*x^2};
\addplot[only marks,ink] coordinates {(.75,1.125)};
\end{axis}
\end{tikzpicture}
\caption{The potential generating the coexistence model under the Bernoulli Fisher metric. Interior trajectories move toward its maximum.}
\end{figure}

The Fisher multiplier changes the rate of ascent at different probabilities. Near an interior equilibrium, it rescales the ordinary derivative field by a positive factor. Near the boundary it becomes small, reflecting the local length cost of moving a rare category's probability.

\section{Deriving the simplex natural gradient}
Let $\Phi(\p)$ be smooth on the positive simplex, and choose a smooth extension to its normalization neighborhood. Write $f_i=\partial_i\Phi$ and $\bar f=\sum_i p_if_i$. Consider
\begin{equation}\label{eq:simplexnatural}
\boxed{(\grad_F\Phi)_i=p_i\left(\partial_i\Phi-\sum_jp_j\partial_j\Phi\right).}
\end{equation}
The components sum to zero. For every tangent vector $v$,
\begin{align*}
g_\p(\grad_F\Phi,v)
&=\sum_i(f_i-\bar f)v_i\\
&=\sum_if_iv_i-\bar f\sum_iv_i=d\Phi_\p(v).
\end{align*}
It therefore satisfies the gradient definition.

The extension does not affect the result. Two extensions agreeing on the simplex have the same directional derivatives on its tangent plane. Their ambient gradients differ there by a multiple of $\ones$, which cancels under the weighted mean subtraction. The resulting vector depends only on the function restricted to the simplex.

\begin{example}[A linear objective]
Let $\Phi=2p_A+p_B$ and $\p=(1/3,1/3,1/3)$. Find the natural gradient and its ascent rate.
\begin{solution}
The ambient derivative vector is $(2,1,0)$, with probability-weighted mean 1. The natural gradient is $(1/3,0,-1/3)$. The directional derivative is $2(1/3)+0=2/3$. Its squared Fisher norm is $3(1/9+1/9)=2/3$.
\end{solution}
\end{example}

For comparison within the same constrained problem, the Euclidean tangent gradient of this objective is $(1,0,-1)$, obtained by subtracting the arithmetic mean of its ambient derivative components. Following that constant vector from the uniform point reaches the boundary in finite time. The Fisher gradient adjusts with the probabilities, and its constant-fitness solution remains positive at every finite time.

\section{Ascent, variance, and symmetric interactions}\index{potential}\index{symmetric payoff matrix}
Natural-gradient ascent gives a particularly direct relation between the objective and the velocity. The rate at which the objective increases is the squared Fisher speed. For replicator equations that speed is already known from Chapter 7: it is fitness variance. Thus one identity can be read as a statement about optimization, motion, or variation among the types.
Along natural-gradient ascent,
\begin{equation}\label{eq:ascentidentity}
\boxed{\dot\Phi=d\Phi(\grad_F\Phi)
=\|\grad_F\Phi\|_F^2=\Var_\p(f)\geq0.}
\end{equation}
The first equality is the chain rule, the second is the gradient definition, and the third is the replicator speed identity. At an interior point, ascent vanishes exactly when all fitness components agree, which is the equilibrium condition.

The scalar potential supplies a monotone quantity for the dynamics. Near a strict local maximum, $\Phi(q)-\Phi(p)$ is a natural Lyapunov candidate. Compact sublevels and isolated critical points permit the stability arguments of Chapter 6 to be applied.

If $M$ is symmetric, the quadratic potential
\begin{equation}\label{eq:symmetricpotential}
\Phi(\p)=\frac12\p^TM\p
\end{equation}
has ambient derivative $M\p$. Thus every symmetric-matrix replicator system is natural-gradient ascent for this potential. Its mean fitness is $\p^TM\p=2\Phi$, so in this class $\bar f'=2\Var_\p(f)$.

\begin{example}[A concave interaction potential]
Let $M=-I$ for three types. Find the dynamics and their velocity at $(0.6,0.3,0.1)$.
\begin{solution}
The potential is $\Phi=-\tfrac12\sum_i p_i^2$, with fitnesses $f_i=-p_i$. Writing $S=\sum_i p_i^2$, the replicator equation is
\[
p_i'=p_i(S-p_i).
\]
At the given state $S=0.46$, so $\p'=(-0.084,0.048,0.036)$. Its ascent rate is $0.0324$. The uniform law uniquely maximizes this strictly concave potential. A KL Lyapunov function can establish convergence of every interior trajectory, as an exercise below shows.
\end{solution}
\end{example}

\section{Stability at a nondegenerate potential maximum}\index{Hessian}
Suppose $q$ is an interior critical point of a smooth potential, expressed in independent coordinates. Write $G(q)$ for its metric matrix and $H(q)$ for the coordinate Hessian of the potential. The gradient field is $F=G^{-1}\nabla\Phi$. Differentiating at a critical point eliminates all terms that multiply $\nabla\Phi(q)=0$, leaving
\[
DF(q)=G(q)^{-1}H(q).
\]
Although this product need not be symmetric in the displayed coordinates, it is similar to a symmetric matrix. The positive-definite square root of $G$ gives
\[
G^{1/2}(G^{-1}H)G^{-1/2}=G^{-1/2}HG^{-1/2}.
\]
The right-hand side is symmetric because $H$ is symmetric. The linearized gradient field therefore has real eigenvalues.

If the Hessian is negative definite, then so is $G^{-1/2}HG^{-1/2}$: its quadratic form on $v$ is the Hessian quadratic form on $G^{-1/2}v$. All eigenvalues are negative, giving a locally asymptotically stable equilibrium. A positive-definite Hessian gives repulsion under ascent. An indefinite nonsingular Hessian produces a saddle with increasing and decreasing linearized directions.

\begin{example}[The local rate at the coexistence maximum]
Apply this calculation to $\Phi(x)=3x-2x^2$ at $q=3/4$.
\begin{solution}
The inverse Bernoulli metric is $q(1-q)=3/16$, and the Hessian is the scalar $-4$. Their product is $-3/4$, reproducing the local rate obtained from the original replicator equation. The potential determines the sign of attraction, while the metric also determines its rate.
\end{solution}
\end{example}

This criterion distinguishes two ways in which trajectories can approach a stable state. A nondegenerate gradient maximum has real linearized decay rates. The cyclic family with $s<1$ has complex linearized eigenvalues and spiraling decay. Its equilibrium attracts, but its rotating linearized motion cannot be the gradient of a smooth potential for a positive-definite metric near that critical point. Convergence by itself is insufficient to identify gradient structure.

The Hessian test is local. It determines behavior near a nondegenerate critical point, while a global convergence statement also needs control of other critical points and of trajectories approaching the boundary. The explicit KL and entropy flows below permit both local and global behavior to be calculated directly.

\section{Linear potentials and exponential evolution}
For $\Phi=\sum_i c_ip_i$, the fitnesses are constants $c_i$. Its natural-gradient flow is the exponential-normalization solution
\[
p_i(t)=\frac{p_i(0)e^{c_it}}{\sum_jp_j(0)e^{c_jt}}.
\]
The formula from Chapter 5 therefore describes an exact optimization flow. The potential is mean constant fitness, and its rate of increase is the variance of the fixed fitness values.

This flow generally approaches the boundary where the highest-fitness types are supported. Positivity at finite time and convergence to a boundary optimum are compatible: a solution can have a positive coordinate at every finite time while that coordinate tends to zero asymptotically.

\section{A KL potential gives mixture contraction}\index{KL divergence!gradient flow}
Fix a positive target $\q$ and choose
\[
\Phi(\p)=-D(\q\|\p).
\]
Its ambient partial derivatives are $f_i=q_i/p_i$, whose weighted mean is 1. The natural-gradient equation simplifies to
\begin{equation}\label{eq:klgradient}
\boxed{p_i'=q_i-p_i.}
\end{equation}
The exact solution is
\begin{equation}\label{eq:klcontraction}
\p(t)=\q+e^{-t}(\p(0)-\q).
\end{equation}
It follows a mixture segment and converges exponentially in probability coordinates. Its geometric route is generally different from the Fisher geodesic connecting the initial law to the target.

The decay identity is
\begin{equation}\label{eq:klgradientdecay}
\frac{d}{dt}D(\q\|\p(t))=-\sum_i\frac{(q_i-p_i)^2}{p_i}.
\end{equation}
For example, direct differentiation gives $1-\sum_iq_i^2/p_i$, and expanding the squared expression proves the equality.

\begin{example}[A target-seeking flow]
Take $\q=(0.3,0.4,0.3)$ and $\p(0)=(0.65,0.2,0.15)$. Find $\p(1)$.
\begin{solution}
The initial displacement is $(0.35,-0.2,-0.15)$. Therefore
\[
\p(1)=(0.3+0.35e^{-1},\ 0.4-0.2e^{-1},\ 0.3-0.15e^{-1}).
\]
Each displacement coordinate contracts by $e^{-1}$, and their sum remains zero.
\end{solution}
\end{example}

The original replicator fitnesses $q_i/p_i$ are defined only in the interior and become singular at a boundary where $q_i>0$. Consequently the smooth-fitness invariant-face theorem from Chapter 3 does not apply to them at that boundary. The simplified vector field $\q-\p$ has a smooth extension that can introduce previously absent types. Specifying the domain and regularity resolves this boundary behavior.

\section{The reverse KL potential and entropy ascent}\index{entropy!gradient flow}
Reverse the directed comparison and let $\Phi(\p)=-D(\p\|\q)$, with the same positive fixed target. Its derivatives are $-\ln(p_i/q_i)-1$, and the flow is
\[
p_i'=p_i\left[-\ln\frac{p_i}{q_i}+D(\p\|\q)\right].
\]
Subtracting the equations for logarithmic ratios gives
\[
\frac{d}{dt}\ln\frac{p_i/q_i}{p_j/q_j}
=-\ln\frac{p_i/q_i}{p_j/q_j}.
\]
Every such log-ratio decays exponentially. For positive initial laws, normalization gives the solution
\begin{equation}\label{eq:reverseklflow}
p_i(t)=\frac{q_i^{1-e^{-t}}p_i(0)^{e^{-t}}}
{\sum_jq_j^{1-e^{-t}}p_j(0)^{e^{-t}}}.
\end{equation}
It converges to $\q$ by a normalized geometric interpolation. The two KL orientations therefore generate different target-seeking paths under the same metric.

For the uniform target $U=(1/n,\ldots,1/n)$, $-D(\p\|U)=H(\p)-\ln n$. Adding a constant to a potential leaves its gradient unchanged, so this is entropy ascent. Directly,
\begin{equation}\label{eq:entropyflow}
p_i'=p_i[-\ln p_i-H(\p)],\qquad
p_i(t)=\frac{p_i(0)^{e^{-t}}}{\sum_jp_j(0)^{e^{-t}}}.
\end{equation}
Entropy increases along this particular dynamical system because it is the generating potential. All initially positive laws approach the uniform distribution.

\begin{example}[Entropy ascent from unequal probabilities]
Let $\p(0)=(1/2,1/3,1/6)$. Find the entropy-ascent state at $t=\ln2$.
\begin{solution}
Here $e^{-t}=1/2$, so the probabilities are proportional to $(\sqrt{1/2},\sqrt{1/3},\sqrt{1/6})$. Normalizing this vector gives the state. At later times the exponent decreases toward zero, and the three positive weights approach equality.
\end{solution}
\end{example}

\section{Cyclic fields and the potential condition}
For the balanced cyclic fitnesses,
\[
f_A=p_B-p_C,\quad f_B=p_C-p_A,\quad f_C=p_A-p_B,
\]
Chapter 3 established nonstationary periodic interior trajectories. If this field were the Fisher gradient of a single-valued potential, then over one period $T$,
\[
0=\Phi(\p(T))-\Phi(\p(0))
=\int_0^T\|\p'(t)\|_F^2\dd t>0,
\]
a contradiction. The Fisher metric still measures its speed and angles; its dynamics circulate without being potential ascent.

Multivariable calculus also supplies a local test. In coordinates $x=p_A,y=p_B$, the differential that a putative potential would need is
\[
d\Phi=(f_A-f_C)dx+(f_B-f_C)dy.
\]
For the balanced cyclic field, this is $(3y-1)dx+(1-3x)dy$. Its cross derivatives are $3$ and $-3$, so it cannot be a differential of a twice differentiable scalar potential on any open region. For general smooth fields on the convex simplex interior, equality of the corresponding cross derivatives is the integrability condition for such a potential.

\begin{figure}[htbp]\centering\begin{tikzpicture}
\begin{scope}[xshift=0.0cm]\begin{scope}[scale=5.7]
\fill[soft] (0,0)--(1,0)--(.5,.8660254)--cycle;
\draw[ink!75] (0,0)--(1,0)--(.5,.8660254)--cycle;
\node[font=\small] at (.5,1.08){Quadratic potential};
\begin{scope}\clip (0,0)--(1,0)--(.5,.8660254)--cycle;
\draw[ink!22] (.5,.2886751) circle(0.08);
\draw[ink!22] (.5,.2886751) circle(0.16);
\draw[ink!22] (.5,.2886751) circle(0.24);
\draw[ink!22] (.5,.2886751) circle(0.32);
\end{scope}
\draw[->,>=Stealth,ink!35,line width=.4pt] (0.5000000,0.6495191)--(0.5000000,0.6175191);
\draw[->,>=Stealth,ink!35,line width=.4pt] (0.5625000,0.5412659)--(0.5647117,0.5093424);
\draw[->,>=Stealth,ink!35,line width=.4pt] (0.6250000,0.4330127)--(0.6161248,0.4022681);
\draw[->,>=Stealth,ink!35,line width=.4pt] (0.6875000,0.3247595)--(0.6564368,0.3324457);
\draw[->,>=Stealth,ink!35,line width=.4pt] (0.7500000,0.2165064)--(0.7234593,0.2343835);
\draw[->,>=Stealth,ink!35,line width=.4pt] (0.8125000,0.1082532)--(0.7847872,0.1242532);
\draw[->,>=Stealth,ink!35,line width=.4pt] (0.4375000,0.5412659)--(0.4352883,0.5093424);
\draw[->,>=Stealth,ink!35,line width=.4pt] (0.5000000,0.4330127)--(0.5000000,0.4010127);
\draw[->,>=Stealth,ink!35,line width=.4pt] (0.5625000,0.3247595)--(0.5347872,0.3087595);
\draw[->,>=Stealth,ink!35,line width=.4pt] (0.6250000,0.2165064)--(0.5972872,0.2325064);
\draw[->,>=Stealth,ink!35,line width=.4pt] (0.6875000,0.1082532)--(0.6587476,0.1222995);
\draw[->,>=Stealth,ink!35,line width=.4pt] (0.3750000,0.4330127)--(0.3838752,0.4022681);
\draw[->,>=Stealth,ink!35,line width=.4pt] (0.4375000,0.3247595)--(0.4652128,0.3087595);
\draw[->,>=Stealth,ink!35,line width=.4pt] (0.5000000,0.2165064)--(0.5000000,0.2485064);
\draw[->,>=Stealth,ink!35,line width=.4pt] (0.5625000,0.1082532)--(0.5403120,0.1313116);
\draw[->,>=Stealth,ink!35,line width=.4pt] (0.3125000,0.3247595)--(0.3435632,0.3324457);
\draw[->,>=Stealth,ink!35,line width=.4pt] (0.3750000,0.2165064)--(0.4027128,0.2325064);
\draw[->,>=Stealth,ink!35,line width=.4pt] (0.4375000,0.1082532)--(0.4596880,0.1313116);
\draw[->,>=Stealth,ink!35,line width=.4pt] (0.2500000,0.2165064)--(0.2765407,0.2343835);
\draw[->,>=Stealth,ink!35,line width=.4pt] (0.3125000,0.1082532)--(0.3412524,0.1222995);
\draw[->,>=Stealth,ink!35,line width=.4pt] (0.1875000,0.1082532)--(0.2152128,0.1242532);
\draw[teal,line width=1.1pt] plot coordinates {(0.2750000,0.1299038) (0.2831183,0.1342287) (0.2910410,0.1384956) (0.2987586,0.1426986) (0.3062639,0.1468326) (0.3135513,0.1508932) (0.3206171,0.1548769) (0.3274591,0.1587805) (0.3340765,0.1626016) (0.3404698,0.1663385) (0.3466406,0.1699898) (0.3525914,0.1735546) (0.3583254,0.1770324) (0.3638468,0.1804230) (0.3691599,0.1837266) (0.3742699,0.1869436) (0.3791821,0.1900748) (0.3839022,0.1931211) (0.3884360,0.1960834) (0.3927894,0.1989630) (0.3969685,0.2017613) (0.4009794,0.2044796) (0.4048281,0.2071195) (0.4085207,0.2096826) (0.4120630,0.2121705) (0.4154609,0.2145848) (0.4187200,0.2169273) (0.4218460,0.2191997) (0.4248442,0.2214036) (0.4277198,0.2235409) (0.4304780,0.2256132) (0.4331237,0.2276223) (0.4356616,0.2295697) (0.4380962,0.2314573) (0.4404320,0.2332866) (0.4426731,0.2350593) (0.4448236,0.2367769) (0.4468875,0.2384411) (0.4488683,0.2400533) (0.4507697,0.2416151) (0.4525951,0.2431280) (0.4543478,0.2445933) (0.4560308,0.2460125) (0.4576472,0.2473869) (0.4591999,0.2487180) (0.4606914,0.2500069) (0.4621245,0.2512550) (0.4635016,0.2524635) (0.4648251,0.2536336) (0.4660973,0.2547665) (0.4673203,0.2558634) (0.4684962,0.2569252) (0.4696270,0.2579532) (0.4707145,0.2589483) (0.4717606,0.2599116) (0.4727669,0.2608440) (0.4737352,0.2617466) (0.4746669,0.2626203) (0.4755636,0.2634659) (0.4764267,0.2642843) (0.4772575,0.2650764) (0.4780574,0.2658431) (0.4788277,0.2665851) (0.4795694,0.2673032) (0.4802838,0.2679981) (0.4809719,0.2686707) (0.4816348,0.2693215) (0.4822734,0.2699513) (0.4828888,0.2705608) (0.4834818,0.2711506) (0.4840533,0.2717213) (0.4846042,0.2722735) (0.4851352,0.2728079) (0.4856471,0.2733250) (0.4861406,0.2738253) (0.4866165,0.2743094) (0.4870754,0.2747778) (0.4875180,0.2752310) (0.4879448,0.2756696) (0.4883565,0.2760938) (0.4887537,0.2765044) (0.4891368,0.2769015) (0.4895065,0.2772858) (0.4898632,0.2776576) (0.4902073,0.2780173) (0.4905394,0.2783653) (0.4908600,0.2787020) (0.4911693,0.2790277) (0.4914679,0.2793429) (0.4917561,0.2796478) (0.4920343,0.2799427) (0.4923029,0.2802281) (0.4925622,0.2805041) (0.4928125,0.2807712) (0.4930542,0.2810295) (0.4932876,0.2812795) (0.4935130,0.2815213) (0.4937307,0.2817552) (0.4939408,0.2819814) (0.4941439,0.2822003) (0.4943399,0.2824121) (0.4945293,0.2826169) (0.4947122,0.2828151) (0.4948889,0.2830068) (0.4950596,0.2831922) (0.4952245,0.2833716) (0.4953838,0.2835451) (0.4955377,0.2837130) (0.4956864,0.2838754) (0.4958300,0.2840325) (0.4959688,0.2841844) (0.4961029,0.2843314) (0.4962325,0.2844735) (0.4963578,0.2846111) (0.4964788,0.2847441) (0.4965957,0.2848728) (0.4967087,0.2849973) (0.4968179,0.2851177) (0.4969235,0.2852342) (0.4970255,0.2853468) (0.4971241,0.2854558) (0.4972194,0.2855612) (0.4973115,0.2856632) (0.4974005,0.2857618) (0.4974866,0.2858572) (0.4975697,0.2859495) (0.4976501,0.2860388) (0.4977278,0.2861251) (0.4978030,0.2862086) (0.4978756,0.2862894) (0.4979458,0.2863676) (0.4980136,0.2864432) (0.4980792,0.2865163) (0.4981427,0.2865870) (0.4982040,0.2866554) (0.4982633,0.2867216) (0.4983206,0.2867856) (0.4983760,0.2868475) (0.4984295,0.2869073) (0.4984813,0.2869653) (0.4985314,0.2870213) (0.4985798,0.2870755) (0.4986266,0.2871279) (0.4986719,0.2871786) (0.4987156,0.2872276) (0.4987579,0.2872750) (0.4987989,0.2873209) (0.4988384,0.2873653) (0.4988767,0.2874082) (0.4989136,0.2874497) (0.4989494,0.2874899) (0.4989840,0.2875287) (0.4990174,0.2875663) (0.4990497,0.2876026) (0.4990810,0.2876378) (0.4991112,0.2876718) (0.4991405,0.2877047) (0.4991687,0.2877365) (0.4991961,0.2877672) (0.4992225,0.2877970) (0.4992481,0.2878258) (0.4992728,0.2878536) (0.4992967,0.2878805) (0.4993198,0.2879066) (0.4993422,0.2879317) (0.4993638,0.2879561) (0.4993847,0.2879797) (0.4994049,0.2880025) (0.4994245,0.2880245) (0.4994434,0.2880458) (0.4994617,0.2880665) (0.4994794,0.2880864) (0.4994965,0.2881057) (0.4995130,0.2881244) (0.4995290,0.2881424) (0.4995445,0.2881599) (0.4995594,0.2881768) (0.4995739,0.2881931) (0.4995879,0.2882089) (0.4996014,0.2882242) (0.4996145,0.2882390) (0.4996272,0.2882533) (0.4996394,0.2882671) (0.4996512,0.2882805) (0.4996627,0.2882934) (0.4996738,0.2883059) (0.4996845,0.2883180) (0.4996948,0.2883297) (0.4997048,0.2883411) (0.4997145,0.2883520) (0.4997239,0.2883626) (0.4997330,0.2883728) (0.4997417,0.2883828) (0.4997502,0.2883923) (0.4997584,0.2884016) (0.4997663,0.2884106) (0.4997740,0.2884192) (0.4997814,0.2884276) (0.4997886,0.2884357) (0.4997955,0.2884436) (0.4998022,0.2884512) (0.4998087,0.2884585) (0.4998150,0.2884656) (0.4998211,0.2884725) (0.4998269,0.2884791) (0.4998326,0.2884856) (0.4998381,0.2884918) (0.4998434,0.2884978) (0.4998485,0.2885036) (0.4998535,0.2885092) (0.4998583,0.2885147) (0.4998630,0.2885199) (0.4998675,0.2885250) (0.4998718,0.2885299) (0.4998760,0.2885347) (0.4998801,0.2885393) (0.4998840,0.2885438) (0.4998878,0.2885481) (0.4998915,0.2885522) (0.4998951,0.2885563) (0.4998985,0.2885602) (0.4999018,0.2885639) (0.4999050,0.2885676) (0.4999082,0.2885711) (0.4999112,0.2885745) (0.4999141,0.2885778) (0.4999169,0.2885810) (0.4999196,0.2885841) (0.4999223,0.2885871) (0.4999248,0.2885900) (0.4999273,0.2885927) (0.4999297,0.2885954) (0.4999320,0.2885981) (0.4999342,0.2886006) (0.4999364,0.2886030) (0.4999384,0.2886054) (0.4999405,0.2886077) (0.4999424,0.2886099) (0.4999443,0.2886120) (0.4999461,0.2886141) (0.4999479,0.2886161) (0.4999496,0.2886180) (0.4999513,0.2886199) (0.4999529,0.2886217) (0.4999544,0.2886235) (0.4999559,0.2886252) (0.4999573,0.2886268) (0.4999587,0.2886284) (0.4999601,0.2886299) (0.4999614,0.2886314) (0.4999627,0.2886328) (0.4999639,0.2886342) (0.4999651,0.2886356) (0.4999662,0.2886369) (0.4999673,0.2886381) (0.4999684,0.2886393) (0.4999694,0.2886405) (0.4999704,0.2886416) (0.4999714,0.2886427) (0.4999723,0.2886438) (0.4999733,0.2886448) (0.4999741,0.2886458) (0.4999750,0.2886468) (0.4999758,0.2886477) (0.4999766,0.2886486) (0.4999774,0.2886495) (0.4999781,0.2886503) (0.4999788,0.2886511) (0.4999795,0.2886519) (0.4999802,0.2886527) (0.4999808,0.2886534) (0.4999815,0.2886541) (0.4999821,0.2886548) (0.4999827,0.2886555) (0.4999832,0.2886561) (0.4999838,0.2886568) (0.4999843,0.2886574) (0.4999848,0.2886579) (0.4999853,0.2886585) (0.4999858,0.2886590) (0.4999863,0.2886596) (0.4999867,0.2886601) (0.4999872,0.2886606) (0.4999876,0.2886611) (0.4999880,0.2886615) (0.4999884,0.2886620) (0.4999888,0.2886624) (0.4999891,0.2886628) (0.4999895,0.2886632) (0.4999898,0.2886636) (0.4999902,0.2886640) (0.4999905,0.2886643) (0.4999908,0.2886647) (0.4999911,0.2886650) (0.4999914,0.2886654) (0.4999917,0.2886657) (0.4999919,0.2886660) (0.4999922,0.2886663) (0.4999925,0.2886666) (0.4999927,0.2886669) (0.4999929,0.2886671) (0.4999932,0.2886674) (0.4999934,0.2886677) (0.4999936,0.2886679) (0.4999938,0.2886681) (0.4999940,0.2886684) (0.4999942,0.2886686) (0.4999944,0.2886688) (0.4999946,0.2886690) (0.4999948,0.2886692) (0.4999949,0.2886694) (0.4999951,0.2886696) (0.4999953,0.2886698) (0.4999954,0.2886700) (0.4999956,0.2886701) (0.4999957,0.2886703) (0.4999959,0.2886704) (0.4999960,0.2886706) (0.4999961,0.2886707) (0.4999963,0.2886709) (0.4999964,0.2886710) (0.4999965,0.2886712) (0.4999966,0.2886713) (0.4999967,0.2886714) (0.4999968,0.2886715) (0.4999969,0.2886717) (0.4999970,0.2886718) (0.4999971,0.2886719) (0.4999972,0.2886720) (0.4999973,0.2886721) (0.4999974,0.2886722) (0.4999975,0.2886723) (0.4999976,0.2886724) (0.4999977,0.2886725) (0.4999977,0.2886726) (0.4999978,0.2886726) (0.4999979,0.2886727) (0.4999979,0.2886728) (0.4999980,0.2886729) (0.4999981,0.2886730) (0.4999981,0.2886730) (0.4999982,0.2886731) (0.4999983,0.2886732) (0.4999983,0.2886732) (0.4999984,0.2886733) (0.4999984,0.2886734) (0.4999985,0.2886734) (0.4999985,0.2886735) (0.4999986,0.2886735) (0.4999986,0.2886736) (0.4999987,0.2886736)};
\fill[copper] (0.2750000,0.1299038) circle(.009);
\end{scope}\end{scope}
\begin{scope}[xshift=7.4cm]\begin{scope}[scale=5.7]
\fill[soft] (0,0)--(1,0)--(.5,.8660254)--cycle;
\draw[ink!75] (0,0)--(1,0)--(.5,.8660254)--cycle;
\node[font=\small] at (.5,1.08){Balanced cyclic model};
\draw[->,>=Stealth,ink!35,line width=.4pt] (0.5000000,0.6495191)--(0.5320000,0.6495191);
\draw[->,>=Stealth,ink!35,line width=.4pt] (0.5625000,0.5412659)--(0.5876429,0.5214710);
\draw[->,>=Stealth,ink!35,line width=.4pt] (0.6250000,0.4330127)--(0.6437266,0.4070644);
\draw[->,>=Stealth,ink!35,line width=.4pt] (0.6875000,0.3247595)--(0.7006086,0.2955677);
\draw[->,>=Stealth,ink!35,line width=.4pt] (0.7500000,0.2165064)--(0.7545714,0.1848346);
\draw[->,>=Stealth,ink!35,line width=.4pt] (0.8125000,0.1082532)--(0.7965000,0.0805404);
\draw[->,>=Stealth,ink!35,line width=.4pt] (0.4375000,0.5412659)--(0.4626429,0.5610607);
\draw[->,>=Stealth,ink!35,line width=.4pt] (0.5000000,0.4330127)--(0.5320000,0.4330127);
\draw[->,>=Stealth,ink!35,line width=.4pt] (0.5625000,0.3247595)--(0.5785000,0.2970467);
\draw[->,>=Stealth,ink!35,line width=.4pt] (0.6250000,0.2165064)--(0.6090000,0.1887935);
\draw[->,>=Stealth,ink!35,line width=.4pt] (0.6875000,0.1082532)--(0.6577857,0.0963763);
\draw[->,>=Stealth,ink!35,line width=.4pt] (0.3750000,0.4330127)--(0.3937266,0.4589610);
\draw[->,>=Stealth,ink!35,line width=.4pt] (0.4375000,0.3247595)--(0.4535000,0.3524723);
\draw[->,>=Stealth,ink!35,line width=.4pt] (0.5000000,0.2165064)--(0.4680000,0.2165064);
\draw[->,>=Stealth,ink!35,line width=.4pt] (0.5625000,0.1082532)--(0.5306648,0.1050096);
\draw[->,>=Stealth,ink!35,line width=.4pt] (0.3125000,0.3247595)--(0.3256086,0.3539514);
\draw[->,>=Stealth,ink!35,line width=.4pt] (0.3750000,0.2165064)--(0.3590000,0.2442192);
\draw[->,>=Stealth,ink!35,line width=.4pt] (0.4375000,0.1082532)--(0.4056648,0.1114967);
\draw[->,>=Stealth,ink!35,line width=.4pt] (0.2500000,0.2165064)--(0.2545714,0.2481781);
\draw[->,>=Stealth,ink!35,line width=.4pt] (0.3125000,0.1082532)--(0.2827857,0.1201301);
\draw[->,>=Stealth,ink!35,line width=.4pt] (0.1875000,0.1082532)--(0.1715000,0.1359660);
\draw[teal,line width=1.1pt] plot coordinates {(0.2750000,0.1299038) (0.2687861,0.1359693) (0.2633953,0.1424833) (0.2588263,0.1494582) (0.2550746,0.1569049) (0.2521329,0.1648330) (0.2499910,0.1732496) (0.2486364,0.1821597) (0.2480539,0.1915650) (0.2482262,0.2014642) (0.2491336,0.2118520) (0.2507542,0.2227190) (0.2530639,0.2340513) (0.2560364,0.2458302) (0.2596432,0.2580320) (0.2638539,0.2706276) (0.2686360,0.2835827) (0.2739548,0.2968578) (0.2797744,0.3104083) (0.2860570,0.3241845) (0.2927636,0.3381324) (0.2998546,0.3521938) (0.3072893,0.3663073) (0.3150274,0.3804085) (0.3230284,0.3944310) (0.3312529,0.4083070) (0.3396626,0.4219685) (0.3482206,0.4353475) (0.3568925,0.4483772) (0.3656462,0.4609926) (0.3744523,0.4731313) (0.3832848,0.4847336) (0.3921209,0.4957435) (0.4009413,0.5061090) (0.4097304,0.5157821) (0.4184761,0.5247190) (0.4271700,0.5328806) (0.4358070,0.5402320) (0.4443852,0.5467428) (0.4529057,0.5523867) (0.4613725,0.5571418) (0.4697916,0.5609897) (0.4781711,0.5639162) (0.4865209,0.5659106) (0.4948517,0.5669656) (0.5031752,0.5670775) (0.5115033,0.5662458) (0.5198475,0.5644736) (0.5282191,0.5617671) (0.5366280,0.5581364) (0.5450828,0.5535949) (0.5535902,0.5481598) (0.5621545,0.5418522) (0.5707776,0.5346974) (0.5794583,0.5267249) (0.5881924,0.5179683) (0.5969722,0.5084660) (0.6057865,0.4982605) (0.6146204,0.4873987) (0.6234557,0.4759319) (0.6322703,0.4639151) (0.6410390,0.4514071) (0.6497331,0.4384698) (0.6583211,0.4251676) (0.6667692,0.4115670) (0.6750410,0.3977356) (0.6830985,0.3837417) (0.6909025,0.3696531) (0.6984128,0.3555366) (0.7055889,0.3414571) (0.7123903,0.3274769) (0.7187770,0.3136550) (0.7247097,0.3000464) (0.7301506,0.2867018) (0.7350630,0.2736670) (0.7394122,0.2609829) (0.7431650,0.2486850) (0.7462905,0.2368035) (0.7487598,0.2253635) (0.7505460,0.2143848) (0.7516244,0.2038824) (0.7519728,0.1938668) (0.7515709,0.1843441) (0.7504007,0.1753167) (0.7484467,0.1667834) (0.7456955,0.1587402) (0.7421366,0.1511801) (0.7377618,0.1440944) (0.7325660,0.1374722) (0.7265472,0.1313016) (0.7197065,0.1255693) (0.7120491,0.1202615) (0.7035838,0.1153640) (0.6943240,0.1108622) (0.6842874,0.1067420) (0.6734970,0.1029891) (0.6619809,0.0995902) (0.6497724,0.0965323) (0.6369106,0.0938032) (0.6234402,0.0913916) (0.6094113,0.0892873) (0.5948793,0.0874811) (0.5799050,0.0859648) (0.5645535,0.0847314) (0.5488941,0.0837752) (0.5329997,0.0830916) (0.5169458,0.0826775) (0.5008096,0.0825308) (0.4846695,0.0826507) (0.4686038,0.0830380) (0.4526900,0.0836944) (0.4370037,0.0846231) (0.4216182,0.0858284) (0.4066030,0.0873161) (0.3920239,0.0890930) (0.3779419,0.0911669) (0.3644132,0.0935472) (0.3514883,0.0962440) (0.3392125,0.0992685) (0.3276253,0.1026326) (0.3167608,0.1063494) (0.3066476,0.1104322) (0.2973089,0.1148950) (0.2887633,0.1197523) (0.2810244,0.1250183) (0.2741016,0.1307074) (0.2680003,0.1368334) (0.2627220,0.1434097) (0.2582651,0.1504485) (0.2546247,0.1579606) (0.2517930,0.1659551) (0.2497596,0.1744390) (0.2485115,0.1834167) (0.2480333,0.1928898) (0.2483073,0.2028562) (0.2493136,0.2133102) (0.2510299,0.2242417) (0.2534319,0.2356363) (0.2564929,0.2474746) (0.2601842,0.2597320) (0.2644750,0.2723789) (0.2693325,0.2853803) (0.2747220,0.2986958) (0.2806071,0.3122802) (0.2869500,0.3260832) (0.2937116,0.3400501) (0.3008520,0.3541224) (0.3083306,0.3682380) (0.3161070,0.3823324) (0.3241409,0.3963389) (0.3323930,0.4101896) (0.3408250,0.4238164) (0.3494006,0.4371516) (0.3580856,0.4501284) (0.3668480,0.4626823) (0.3756593,0.4747510) (0.3844935,0.4862755) (0.3933285,0.4972003) (0.4021455,0.5074738) (0.4109293,0.5170487) (0.4196683,0.5258818) (0.4283547,0.5339345) (0.4369836,0.5411727) (0.4455538,0.5475664) (0.4540667,0.5530900) (0.4625266,0.5577219) (0.4709397,0.5614446) (0.4793146,0.5642442) (0.4876611,0.5661105) (0.4959902,0.5670366) (0.5043137,0.5670194) (0.5126432,0.5660588) (0.5209906,0.5641584) (0.5293667,0.5613249) (0.5377814,0.5575687) (0.5462431,0.5529039) (0.5547580,0.5473481) (0.5633304,0.5409231) (0.5719615,0.5336546) (0.5806498,0.5255728) (0.5893907,0.5167119)};
\fill[copper] (0.2750000,0.1299038) circle(.009);
\end{scope}\end{scope}
\end{tikzpicture}\caption{Natural-gradient ascent for $\Phi=-\tfrac12\sum_i p_i^2$ and balanced cyclic dynamics from the same initial law. Left contours are potential levels. Short arrows show field directions with normalized drawing lengths.}\end{figure}

For numerical motion under a smooth fitness field, a useful finite update is
\[
p_i^{\mathrm{new}}=\frac{p_i e^{\eta f_i(\p)}}{\sum_jp_j e^{\eta f_j(\p)}}.
\]
It stays positive and normalized for finite scores. Expanding in small $\eta$ yields $p_i^{\mathrm{new}}=p_i+\eta p_i(f_i-\bar f)+O(\eta^2)$, so it approximates a replicator step. It is exact for a constant-fitness time step. For a varying fitness field, preservation of the simplex alone does not establish a particular numerical error or monotonicity guarantee.

\FloatBarrier
\section{Exercises}
\begin{exercises}
\item Derive the coordinate gradient formula $G^{-1}\nabla\Phi$ from the defining inner-product identity.
\item Use Cauchy--Schwarz to prove that the normalized metric gradient maximizes directional increase among unit tangent vectors.
\item For $\Phi(x)=2x-(5/2)x^2$, derive the Bernoulli natural-gradient flow and classify its equilibria.
\item Verify \eqref{eq:simplexnatural}, including the tangent constraint and independence from the chosen ambient extension.
\item At $\p=(1/2,1/3,1/6)$ and $\Phi=2p_A+p_B$, compute the natural gradient and its ascent rate.
\item For a symmetric payoff matrix, prove that $\tfrac12\p^TM\p$ is a potential and that mean fitness increases at rate $2\Var_\p(f)$.
\item For $M=-I$ with $n$ types, use the multi-strategy KL derivative from Chapter 6 and a uniform reference to prove global interior convergence to the uniform law.
\item Derive the solution \eqref{eq:klcontraction} and verify the exact divergence decay identity \eqref{eq:klgradientdecay}.
\item Compute the target-seeking state in the worked example at $t=1$ numerically. Determine the time needed to reduce every coordinate displacement by a factor of 10.
\item Explain why the smooth extension $p'=q-p$ can introduce a type that was absent initially without contradicting the smooth replicator invariant-face theorem.
\item Derive the reverse-KL gradient and the normalized geometric solution \eqref{eq:reverseklflow}.
\item Derive the entropy gradient and show that its log-probability ratios decay exponentially.
\item For entropy ascent from $(1/2,1/3,1/6)$, find the normalized state at $t=\ln2$ to four decimal places.
\item Prove that no nonstationary periodic trajectory can be a gradient-ascent trajectory for any positive-definite Riemannian metric and a single-valued potential.
\item In reduced simplex coordinates, derive the differential coefficients $f_A-f_C$ and $f_B-f_C$. Apply the cross-derivative test to balanced cyclic competition.
\item Expand the exponential finite-step update through first order in $\eta$ and recover the replicator field.
\item \emph{Exploration.} Compare a Fisher geodesic to the mixture-contraction trajectory with the same initial and target laws. Measure route length separately from the rate of potential improvement.
\item \emph{Exploration.} Compare linear-potential ascent, entropy ascent, and balanced cyclic motion from the same positive initial state. Track a justified potential or conserved quantity for each system.
\end{exercises}

\chapter{Gaussian Families and Hyperbolic Geometry}
\label{ch:gaussian}

A finite-dimensional statistical family can describe observations on a continuous outcome space. The normal distributions form a two-parameter example: a point specifies a mean and a standard deviation, and therefore an entire density curve. Its Fisher metric turns this parameter half-plane into a model of hyperbolic geometry.

Changing the mean moves a Gaussian bell curve horizontally; changing its standard deviation alters its width and height. A path in the parameter half-plane therefore represents a continuous deformation of density curves. The Fisher metric measures the statistical length of that deformation. We will find that geodesics can become wider on their way between two distributions with the same standard deviation, a distinctive feature of this geometry.

\section{Continuous probability densities}\index{probability density}\index{Gaussian family}
A probability density $\rho(x)$ on $\R$ is nonnegative and satisfies $\int_{\R}\rho(x)dx=1$. Probabilities are areas:
\[
\Pr(a\leq X\leq b)=\int_a^b\rho(x)\dd x.
\]
The density's value is probability per unit of the observation coordinate. A single point has probability zero under a continuous density, while intervals can have positive probability. A density can exceed 1 on a short interval while its total area remains 1.

Expectation is defined by replacing a probability-weighted sum with an integral:
\[
\E[g(X)]=\int_{\R}g(x)\rho(x)\dd x,
\]
whenever the integral exists. The definitions of mean, variance, and score products then take the same form as in the finite case.

The normal, or Gaussian, law $N(\mu,\sigma^2)$ has density
\begin{equation}\label{eq:gaussiandensity}
\rho(x;\mu,\sigma)=\frac1{\sigma\sqrt{2\pi}}
\exp\left[-\frac{(x-\mu)^2}{2\sigma^2}\right],\qquad \sigma>0.
\end{equation}
The substitution $z=(x-\mu)/\sigma$ reduces its normalization to the Gaussian integral $\int_{\R}e^{-z^2/2}dz=\sqrt{2\pi}$. Its mean is $\mu$ and standard deviation is $\sigma$. The parameter $\sigma^2$ in the distribution's name is its variance; the coordinate used initially in this chapter is $\sigma$.

\begin{figure}[htbp]
\centering
\begin{tikzpicture}
\begin{axis}[width=.9\textwidth,height=5.3cm,xmin=-6,xmax=6,ymin=0,ymax=.46,xlabel={observed value $x$},ylabel={density},grid=major,legend pos=north west]
\addplot[teal,domain=-6:6,samples=160]{exp(-x^2/2)/sqrt(2*pi)};\addlegendentry{$N(0,1)$}
\addplot[copper,domain=-6:6,samples=160]{exp(-(x-2)^2/2)/sqrt(2*pi)};\addlegendentry{$N(2,1)$}
\addplot[ink,dashed,domain=-6:6,samples=160]{exp(-x^2/8)/(2*sqrt(2*pi))};\addlegendentry{$N(0,4)$}
\end{axis}
\end{tikzpicture}
\caption{Three density curves represented by points in the Gaussian parameter family. Each curve has total area one.}
\end{figure}

\section{Gaussian moments needed for information}
Let $Z=(X-\mu)/\sigma$ be standard normal, with density $\phi(z)=e^{-z^2/2}/\sqrt{2\pi}$. Symmetry gives zero for odd moments. Since $\phi'(z)=-z\phi(z)$, integration by parts gives, for $k\geq2$,
\[
\int_{\R}z^k\phi(z)\dd z
=-\int_{\R}z^{k-1}\phi'(z)\dd z
=(k-1)\int_{\R}z^{k-2}\phi(z)\dd z.
\]
The boundary term vanishes because the Gaussian exponential decays faster than any polynomial. Starting from $\E[1]=1$, the formula gives $\E[Z^2]=1$ and $\E[Z^4]=3$. Therefore, with $Y=X-\mu$,
\begin{equation}\label{eq:gaussianmoments}
\E[Y]=0,\quad\E[Y^2]=\sigma^2,\quad
\E[Y^3]=0,\quad\E[Y^4]=3\sigma^4.
\end{equation}
These moments are sufficient for every Fisher matrix entry below.

\section{Scores and the Gaussian Fisher matrix}\index{Fisher information matrix!Gaussian}
The log density is
\[
\ell(x;\mu,\sigma)=-\ln\sigma-\frac12\ln(2\pi)
-\frac{(x-\mu)^2}{2\sigma^2}.
\]
Its score components are
\begin{equation}\label{eq:gaussianscores}
S_\mu=\frac{x-\mu}{\sigma^2},\qquad
S_\sigma=-\frac1\sigma+\frac{(x-\mu)^2}{\sigma^3}.
\end{equation}
Equivalently, $S_\mu=Z/\sigma$ and $S_\sigma=(Z^2-1)/\sigma$. Their means vanish by the moment identities. Their expected products give
\begin{align*}
G_{\mu\mu}&=\frac{\E[Z^2]}{\sigma^2}=\frac1{\sigma^2},\\
G_{\mu\sigma}&=\frac{\E[Z^3-Z]}{\sigma^2}=0,\\
G_{\sigma\sigma}&=\frac{\E[Z^4-2Z^2+1]}{\sigma^2}=\frac2{\sigma^2}.
\end{align*}
Hence
\begin{equation}\label{eq:gaussianmatrix}
\boxed{G(\mu,\sigma)=\begin{pmatrix}\sigma^{-2}&0\\0&2\sigma^{-2}\end{pmatrix},
\qquad ds_F^2=\frac{d\mu^2+2d\sigma^2}{\sigma^2}.}
\end{equation}
The off-diagonal entry vanishes because an odd score is orthogonal, in the expectation inner product, to an even score. A small mean change and a small scale change therefore give Fisher-orthogonal directions in these coordinates.

\begin{example}[Local costs at a Gaussian law]
At $N(0,4)$, compare small changes $(d\mu,d\sigma)=(h,0)$ and $(0,h)$.
\begin{solution}
Here $\sigma=2$. The squared lengths are $h^2/4$ and $h^2/2$. Their lengths are $|h|/2$ and $|h|/\sqrt2$. The scale direction is $\sqrt2$ times as long for equal numerical increments. Both costs grow if the base Gaussian becomes narrower.
\end{solution}
\end{example}

\begin{example}[Using variance as a coordinate]
Rewrite the metric using $v=\sigma^2$.
\begin{solution}
Since $d\sigma=dv/(2\sqrt v)$,
\[
ds_F^2=\frac{d\mu^2}{v}+\frac{dv^2}{2v^2}.
\]
Thus the Fisher matrix in $(\mu,v)$ coordinates is $\diag(1/v,1/(2v^2))$. The different matrix expresses the same metric under the coordinate transformation.
\end{solution}
\end{example}

\section{The hyperbolic half-plane}\index{hyperbolic geometry}
Set
\[
u=\frac{\mu}{\sqrt2},\qquad v=\sigma>0.
\]
Then $d\mu^2=2du^2$, giving
\begin{equation}\label{eq:hyperbolicmetric}
\boxed{ds_F^2=2\frac{du^2+dv^2}{v^2}.}
\end{equation}
The metric $(du^2+dv^2)/v^2$ is the Poincar\'e upper-half-plane metric. The Gaussian Fisher metric is twice that metric, so lengths are $\sqrt2$ times the corresponding hyperbolic lengths.

At height $v$, the metric multiplies ordinary Euclidean lengths by $\sqrt2/v$. This factor changes with height and is the same in every tangent direction at a point. Thus the half-plane drawing preserves angles, and horizontal motion has lower length cost at larger $v$.

The half-plane has useful isometries. The transformation
\[
(u,v)\longmapsto(au+b,av),\qquad a>0,
\]
multiplies both numerator differentials and denominator height by $a$, leaving the metric unchanged. In Gaussian coordinates these include common translations of the means and common positive rescaling of both means and standard deviations.

\section{Vertical and circular geodesics}\index{geodesic!Gaussian}
The half-plane metric makes movement less costly, per unit of Euclidean displacement, at greater height. A path between two means can therefore reduce its total Fisher length by rising to a larger standard deviation before coming back down. Vertical paths and circular arcs describe exactly how shortest routes balance this effect.
When two states have the same mean, a vertical path connects them. Every path obeys the lower bound
\[
L_F=\sqrt2\int\frac{\sqrt{u'^2+v'^2}}v\dd t
\geq\sqrt2\left|\int\frac{v'}v\dd t\right|
=\sqrt2\left|\ln\frac{v_2}{v_1}\right|.
\]
A monotone vertical path attains it. Therefore
\begin{equation}\label{eq:verticaldistance}
d_F((\mu,\sigma_1),(\mu,\sigma_2))
=\sqrt2\left|\ln\frac{\sigma_2}{\sigma_1}\right|.
\end{equation}
Constant-speed vertical travel has $v(t)=v_1\exp[t\ln(v_2/v_1)]$ on $[0,1]$.

For nonvertical geodesics, one can use the Euler--Lagrange equations for energy. Omitting a constant factor that does not affect the equations, the energy integrand is $(u'^2+v'^2)/(2v^2)$. Its equations are
\begin{equation}\label{eq:hyperbolicgeoeq}
u''-\frac{2u'v'}v=0,\qquad
v''+\frac{u'^2-v'^2}v=0.
\end{equation}
The first gives $u'/v^2=C$. On a nonvertical geodesic $C\ne0$. Constant metric speed also gives $(u'^2+v'^2)/v^2=k^2$. These identities yield
\[
\frac{du}{dv}=\frac{Cv}{\pm\sqrt{k^2-C^2v^2}},
\]
whose integral is a circle $(u-c)^2+v^2=R^2$. The center lies on $v=0$. Thus the geodesics are vertical lines and semicircles orthogonal to the boundary.

For endpoints with $u_1\ne u_2$, the circle center and radius are
\begin{equation}\label{eq:circlecenter}
c=\frac{u_2^2+v_2^2-u_1^2-v_1^2}{2(u_2-u_1)},\qquad
R=\sqrt{(u_1-c)^2+v_1^2}.
\end{equation}
Write $u=c+R\cos\varphi$, $v=R\sin\varphi$, with $0<\varphi<\pi$. Its length element is $\sqrt2\,|d\varphi|/\sin\varphi$. Since $\int d\varphi/\sin\varphi=\ln\tan(\varphi/2)$, the parameter $w=\ln\tan(\varphi/2)$ is an arc coordinate up to the factor $\sqrt2$. Linear interpolation of $w$ gives constant Fisher speed along the arc.

The circular arcs have the same global minimizing property as vertical lines. One way to verify it is to apply a half-plane isometry that maps the circle to a vertical line, where the preceding lower bound applies. In complex notation $z=u+iv$, real translations and $z\mapsto-1/z$ preserve $|dz|/\operatorname{Im}z$. Translating one boundary endpoint of the circle to zero and applying this inversion sends the circle to a vertical line. This also establishes uniqueness of the connecting minimizing arc.

\section{Exact Gaussian distance}\index{Fisher distance!Gaussian}
Once the shape of a minimizing path is known, its length can be expressed entirely in terms of its endpoints. This removes the need to construct the particular circle for every distance calculation. The result is a closed formula involving the differences of means and standard deviations, together with their scales.
Integrating along the connecting semicircle gives
\[
\frac{d_F}{\sqrt2}
=\left|\ln\frac{\tan(\varphi_2/2)}{\tan(\varphi_1/2)}\right|.
\]
Using the half-angle identities together with $u_j=c+R\cos\varphi_j$ and $v_j=R\sin\varphi_j$ simplifies this to
\[
\cosh\left(\frac{d_F}{\sqrt2}\right)
=1+\frac{(u_1-u_2)^2+(v_1-v_2)^2}{2v_1v_2}.
\]
The same expression holds for a vertical pair by \eqref{eq:verticaldistance}. Returning to Gaussian coordinates gives
\begin{equation}\label{eq:gaussiandistance}
\boxed{d_F=\sqrt2\,\arcosh\left(1+
\frac{(\mu_1-\mu_2)^2+2(\sigma_1-\sigma_2)^2}{4\sigma_1\sigma_2}\right).}
\end{equation}
For $z\geq1$, $\arcosh z=\ln(z+\sqrt{z^2-1})$, making the formula directly computable with logarithms and square roots.

\begin{example}[Equal distance for equal scale ratios]
Compare the distances from $N(0,1)$ to $N(0,4)$ and from $N(0,4)$ to $N(0,16)$.
\begin{solution}
The standard deviations are $1,2,4$. Each pair has ratio 2, so each distance is $\sqrt2\ln2\approx0.980258$. Along a fixed-mean geodesic, a scale ratio determines distance.
\end{solution}
\end{example}

\begin{example}[A geodesic that broadens the distribution]
Connect $N(-2,1)$ to $N(2,1)$ in the full Gaussian family.
\begin{solution}
The half-plane endpoints are $(-\sqrt2,1)$ and $(\sqrt2,1)$. Their connecting circle has center 0 and radius $\sqrt3$. The midpoint reaches standard deviation $\sqrt3$. The distance is
\[
d_F=\sqrt2\,\arcosh5\approx3.242.
\]
The horizontal path keeping $\sigma=1$ has length 4. Moving through broader intermediate laws reduces the total intrinsic length.
\end{solution}
\end{example}

\begin{figure}[htbp]
\centering
\begin{tikzpicture}
\begin{axis}[width=.9\textwidth,height=6cm,axis equal image,xmin=-2,xmax=2,ymin=0,ymax=2.15,xlabel={$u=\mu/\sqrt2$},ylabel={$v=\sigma$},grid=major,legend pos=north east]
\addplot[teal,domain=35.2643897:144.7356103,samples=100] ({sqrt(3)*cos(x)},{sqrt(3)*sin(x)});\addlegendentry{Fisher geodesic}
\addplot[copper,dashed] coordinates {(-1.41421356,1)(1.41421356,1)};\addlegendentry{fixed-scale path}
\addplot[only marks,ink] coordinates {(-1.41421356,1)(1.41421356,1)(0,1.73205081)};
\end{axis}
\end{tikzpicture}
\caption{The connecting Gaussian geodesic rises to a larger standard deviation. Equal scales on the two drawing axes make the orthogonal-circle geometry visible.}
\end{figure}

For a general straight parameter interpolation, the Fisher length can also be written explicitly. Let $A=\sqrt{(\mu_2-\mu_1)^2+2(\sigma_2-\sigma_1)^2}$. Its length is
\begin{equation}\label{eq:gaussianstraightlength}
L_{\mathrm{straight}}=
\begin{cases}
\displaystyle A\,\frac{\ln(\sigma_2/\sigma_1)}{\sigma_2-\sigma_1},&\sigma_2\ne\sigma_1,\\[7pt]
\displaystyle\frac{|\mu_2-\mu_1|}{\sigma_1},&\sigma_2=\sigma_1.
\end{cases}
\end{equation}
The logarithm and scale difference have the same sign. The expression follows by integrating $A/[(1-t)\sigma_1+t\sigma_2]$. It is at least the full-family geodesic distance, with equality for fixed-mean vertical travel.

\section{Following a Gaussian geodesic in time}
For the equal-scale endpoints $N(-2,1)$ and $N(2,1)$, the circular arc admits a particularly explicit parameterization. Set $a=\arcosh\sqrt3$ and $w=a(2t-1)$ for $0\leq t\leq1$. Define
\begin{equation}\label{eq:gaussianexplicitpath}
\mu(t)=\sqrt6\tanh w,\qquad
\sigma(t)=\sqrt3\sech w,
\end{equation}
where $\sech w=1/\cosh w$. The identity $\tanh^2w+\sech^2w=1$ shows that $(u,v)=(\mu/\sqrt2,\sigma)$ lies on the circle of radius $\sqrt3$ centered at zero.

At $t=0$, $\cosh a=\sqrt3$ and $\tanh a=\sqrt{2/3}$, so the parameters are $(-2,1)$. At $t=1$ they are $(2,1)$. At $t=1/2$, they are $(0,\sqrt3)$, giving the midpoint law $N(0,3)$.

These three points of the half-plane correspond to three visibly different density curves. At the endpoints the distributions have the same spread and different centers. At the midpoint the mean lies halfway between those centers, and the increased standard deviation produces a lower, broader bell curve. The change of width is an essential part of the shortest statistical route.

\begin{figure}[htbp]
\centering
\begin{tikzpicture}
\begin{axis}[
  width=.9\textwidth,height=5.8cm,
  xmin=-5.5,xmax=5.5,ymin=0,ymax=.50,
  xlabel={observed value $x$},ylabel={probability density},
  grid=major
]
\addplot[teal,domain=-5.5:5.5,samples=150]{exp(-((x+2)^2)/2)/sqrt(2*pi)};
\addplot[ink,densely dashed,domain=-5.5:5.5,samples=150]{exp(-(x^2)/6)/sqrt(6*pi)};
\addplot[copper,domain=-5.5:5.5,samples=150]{exp(-((x-2)^2)/2)/sqrt(2*pi)};
\node[font=\scriptsize,text=teal,anchor=south] at (axis cs:-2,.412){$N(-2,1)$};
\node[font=\scriptsize,text=ink,anchor=south] at (axis cs:0,.24){$N(0,3)$};
\node[font=\scriptsize,text=copper,anchor=south] at (axis cs:2,.412){$N(2,1)$};
\end{axis}
\end{tikzpicture}
\caption{Densities along the Gaussian Fisher geodesic pictured in the preceding half-plane diagram. The geometric midpoint is a broader single Gaussian, rather than the arithmetic mixture of the endpoint densities.}
\end{figure}

Differentiating the two parameter functions and substituting into the Fisher speed gives
\[
\frac{\mu'^2+2\sigma'^2}{\sigma^2}=8a^2.
\]
The speed is the constant $2\sqrt2a$. Its length over $[0,1]$ is the same constant. Since $\cosh(2a)=2\cosh^2a-1=5$, this agrees with $\sqrt2\arcosh5$ from the endpoint distance formula.

The interpolation parameter is a choice of clock along the geometric path. A physical or learning process can follow a different time law, and an objective-driven natural-gradient process can follow a different route. Specifying the family, the metric, and the rule of motion determines which of these descriptions is being used.

\begin{example}[Two intermediate probability laws]
Find the mean and variance of the equal mixture of $N(-2,1)$ and $N(2,1)$, and compare it with the Gaussian geodesic midpoint.
\begin{solution}
The mixture first selects either Gaussian with probability $1/2$ and then draws an observation from that selected law. Its mean is zero. Averaging component variances contributes 1, and the separation of the component means contributes 4, so its variance is 5. Its density is the arithmetic average of the two endpoint densities.

The Fisher geodesic midpoint within the Gaussian family is instead the single Gaussian $N(0,3)$, with variance 3. Each construction specifies a legitimate intermediate law. The mixture construction is linear in densities, while the Gaussian geodesic is shortest among paths constrained to remain Gaussian.
\end{solution}
\end{example}

The mixture variance identity from Chapter 4 applies here by the same expansion of first and second moments. It is the distribution family, rather than the continuous or discrete nature of the outcome space, that determines which interpolations remain inside the model. This provides a concrete reason to state the allowed family whenever a statistical distance or connecting curve is calculated.

\section{Gaussian relative entropy}\index{KL divergence!Gaussian}
For densities, relative entropy is
\[
D(P\|R)=\int_{\R}\rho_P(x)\ln\frac{\rho_P(x)}{\rho_R(x)}\dd x.
\]
For $P=N(\mu_1,\sigma_1^2)$ and $R=N(\mu_2,\sigma_2^2)$, the log ratio is
\[
\ln\frac{\sigma_2}{\sigma_1}
-\frac{(x-\mu_1)^2}{2\sigma_1^2}
+\frac{(x-\mu_2)^2}{2\sigma_2^2}.
\]
Under $P$, the two squared expectations are $\sigma_1^2$ and $\sigma_1^2+(\mu_1-\mu_2)^2$. Therefore
\begin{equation}\label{eq:gaussiankl}
\boxed{D(P\|R)=\ln\frac{\sigma_2}{\sigma_1}
+\frac{\sigma_1^2+(\mu_1-\mu_2)^2}{2\sigma_2^2}-\frac12.}
\end{equation}

\begin{example}[Directed comparison and intrinsic length]
Compare $P=N(0,1)$ and $R=N(1,4)$.
\begin{solution}
The directed divergences are
\[
D(P\|R)=\ln2-\frac14,\qquad
D(R\|P)=2-\ln2.
\]
The Fisher distance is $\sqrt2\,\arcosh(11/8)$ and is unchanged upon reversing the endpoints. The divergence and distance agree locally through their second-order relationship, while the displayed formulas describe distinct global comparisons.
\end{solution}
\end{example}

Expanding \eqref{eq:gaussiankl} for $\mu_2=\mu+d\mu$ and $\sigma_2=\sigma+d\sigma$ yields
\[
D(P\|R)=\frac{d\mu^2+2d\sigma^2}{2\sigma^2}
+O(\|(d\mu,d\sigma)\|^3),
\]
recovering the Gaussian Fisher matrix independently from relative entropy.

\section{Curvature, boundary, and completeness}\index{curvature!negative}\index{completeness}
The standard hyperbolic half-plane has curvature $-1$. Scaling a metric by a positive constant $c$ scales lengths by $\sqrt c$ and curvature by $1/c$. Equation \eqref{eq:hyperbolicmetric} therefore gives Gaussian Fisher curvature $-1/2$.

An explicit calculation is available from the curvature formula for a conformal planar metric $ds^2=e^{2\omega}(du^2+dv^2)$:
\[
K=-e^{-2\omega}(\omega_{uu}+\omega_{vv}).
\]
Here $\omega=\tfrac12\ln2-\ln v$, so $\omega_{uu}=0$, $\omega_{vv}=1/v^2$, and $K=-1/2$. The sign describes the intrinsic surface geometry; the half-plane's Euclidean drawing is a coordinate model of that geometry.

The length bound used for vertical geodesics implies that every path approaching $\sigma=0$ or $\sigma=\infty$ has infinite length from a fixed positive scale. The lower boundary is therefore infinitely far away. This differs from the finite-distance categorical boundary.

The Gaussian family is metrically complete. To see this directly, fix a base point and a distance bound. Formula \eqref{eq:gaussiandistance} bounds $\sigma$ above and away from zero, and then bounds $\mu$. The closed metric ball is consequently a closed bounded parameter set contained in $\sigma>0$, hence compact. Every Cauchy sequence stays in such a ball and has an interior limit. Degenerate zero-variance laws are not added by completion of this metric.

\section{The chosen family determines the available paths}
Distance is always measured using curves allowed in the statistical family. For the subfamily with a fixed scale $\sigma_0$, the only coordinate is $\mu$ and the intrinsic distance is $|\mu_2-\mu_1|/\sigma_0$. In the full Gaussian family a shorter curve can change scale along the way, as the broadening example demonstrates.

The square-root idea also extends to densities. The map $\rho\mapsto2\sqrt\rho$ has squared norm $\int4\rho=4$ in the function inner product $\ip{f}{g}=\int fg$. Its differential again produces the Fisher metric. All densities together therefore have a spherical square-root representation. The Gaussian laws occupy a particular two-dimensional surface inside that much larger sphere. A shortest path among all densities will generally leave the Gaussian surface, so its distance need not equal the intrinsic Gaussian distance.

For example, the square-root overlap of two Gaussians is
\[
B(P,R)=\sqrt{\frac{2\sigma_1\sigma_2}{\sigma_1^2+\sigma_2^2}}
\exp\left[-\frac{(\mu_1-\mu_2)^2}{4(\sigma_1^2+\sigma_2^2)}\right].
\]
The full-density spherical distance is $2\arccos B$, bounded by $\pi$. The intrinsic Gaussian distance can be arbitrarily large. The metric is induced from the same local Fisher construction; the allowed connecting curves determine the global distance.

\begin{center}
\begin{tabular}{@{}L{.26\textwidth}L{.36\textwidth}L{.27\textwidth}@{}}
\toprule
Family & Fisher geometry & Boundary and curvature\\\midrule
Bernoulli interior & Radius-2 quarter-circle arc & Finite-distance ends; one-dimensional\\[5pt]
Three-outcome categorical interior & Positive octant of $S^2(2)$ & Finite-distance faces; $K=1/4$\\[5pt]
Full Gaussian family & Scaled hyperbolic half-plane & Complete; $K=-1/2$\\
\bottomrule
\end{tabular}
\end{center}
The common construction begins with a probability model, compares nearby predictions by relative entropy, and extracts an inner product from the leading quadratic term. The resulting metric supports both geometric questions about paths and dynamical questions about gradients. Population proportions, statistical observations, and classical geometries are linked by these explicit calculations.

\FloatBarrier
\section{Exercises}
\begin{exercises}
\item Using $\phi'=-z\phi$, derive the second and fourth standard-normal moments. Explain why the first and third vanish.
\item Derive both Gaussian score components and verify their expectations are zero.
\item Compute every entry of the Gaussian Fisher matrix using the standardized score expressions.
\item Use the coordinate transformation law to obtain the Fisher matrix in $(\mu,\sigma^2)$ coordinates.
\item At $\sigma=3$, calculate the Fisher norms of parameter velocities $(1,0)$, $(0,1)$, and $(1,1)$.
\item Verify directly that $(u,v)\mapsto(au+b,av)$ preserves the Gaussian half-plane metric for $a>0$.
\item Find the Fisher distance from $N(2,9)$ to $N(2,36)$ and give a constant-speed connecting path over $[0,1]$.
\item Derive the two Euler--Lagrange equations \eqref{eq:hyperbolicgeoeq} from the energy integrand.
\item For endpoints $(-\sqrt2,1)$ and $(\sqrt2,1)$ in the half-plane, compute the circle center, radius, Fisher distance, and maximum height.
\item Verify the Gaussian distance formula for equal means and reduce it to the logarithmic scale-ratio formula.
\item Compute both directed divergences between $N(0,1)$ and $N(1,4)$ numerically, together with their Fisher distance.
\item Derive the straight-interpolation length \eqref{eq:gaussianstraightlength}, treating equal endpoint scales separately.
\item Use $w=\ln\tan(\varphi/2)$ to give an explicit constant-speed circle-arc interpolation between two nonvertical half-plane endpoints.
\item Expand Gaussian KL through quadratic order at equal endpoints and recover the Fisher line element.
\item Prove that every curve tending to $\sigma=0$ has infinite Fisher length, allowing its mean to vary along the curve.
\item Explain why the full Gaussian intrinsic distance can exceed the full-density spherical distance while the Gaussian metric is induced by the same square-root construction.
\item \emph{Exploration.} In the Gaussian demonstration, compare geodesics and straight parameter paths for several endpoint pairs. Relate broader intermediate densities to the half-plane length factor.
\item Derive the Gaussian overlap coefficient $B(P,R)$ by completing the square in $\int\sqrt{\rho_P\rho_R}dx$. Apply it to $N(-2,1)$ and $N(2,1)$, and compare the full-density spherical distance with the intrinsic Gaussian distance.
\end{exercises}

\FloatBarrier\backmatter\chapter*{Selected Solutions}\addcontentsline{toc}{chapter}{Selected Solutions}\markboth{Selected Solutions}{}
\section*{Chapter 1}
\begin{answers}
\item[1.1] The totals are $90$ and $240$; A's proportions are $2/5$ and $3/10$. Its count doubles, while its share decreases by $1/10$.
\item[1.3] $x'=0.3(0.7)(0.6)=0.126$ per unit time. The total per-capita rate is $0.3(0.8)+0.7(0.2)=0.38$ per unit time.
\item[1.5] The roots are $0,2/5,1$. The velocity is positive on $(0,2/5)$ and negative on $(2/5,1)$. The interior root attracts, with $F'(2/5)=-6/5$. Both endpoints repel from the interior.
\item[1.7] $x(t)=1/(1+4e^{-2t})$. At $x=4/5$, $4e^{-2t}=1/4$, giving $t=\ln4$.
\item[1.9] Differentiate the product: $F'=(1-2x)\Delta+x(1-x)\Delta'$. Evaluate at $\Delta(e)=0$.
\item[1.11] At the endpoints the vector field points inward: $F(0)=0.1>0$ and $F(1)=-0.9<0$. The explicit solution $x(t)=0.1+(x_0-0.1)e^{-t}$ stays in $[0,1]$ for $t\geq0$. The endpoints are not equilibria of this field.
\item[1.13] Fitness differences are unchanged, so the composition equation is unchanged. The total per-capita rate increases by $c(x)$.
\item[1.15] The equilibrium is $e=a/b$. The derivative is $-be(1-e)$, so a small displacement decays approximately as $\exp[-be(1-e)t]$.
\item[1.17] Since $F'(e)<0$, a sufficiently small punctured neighborhood satisfies $(x-e)F(x)<0$. A solution converging to $e$ eventually lies there. Then $\frac{d}{dt}(x-e)^2=2(x-e)F(x)<0$ away from equilibrium. Stability additionally controls every sufficiently close initial condition for all future times.
\end{answers}

\section*{Chapter 2}
\begin{answers}
\item[2.1] $f_A=5-4x$, $f_B=2+x$, and $\Delta=3-5x$. The equilibria are $0,3/5,1$; the interior one attracts.
\item[2.3] Both endpoint advantages are positive: $b-d=1$ and $a-c=2$. Thus $x'>0$ throughout $(0,1)$ and every interior trajectory tends to 1.
\item[2.5] $\Delta=2-4x$, so $x^*=1/2$ and $F'(x^*)=-1$. Small perturbations decay approximately as $e^{-t}$.
\item[2.7] The necessary and sufficient conditions are $a=c$ and $b=d$: the rows agree.
\item[2.9] Scaling multiplies $F'$ by the same constant. Positive scaling preserves all phase arrows; negative scaling reverses them and interchanges attracting and repelling hyperbolic equilibria.
\item[2.11] The velocity is positive throughout $(0,1)$. Thus 0 repels and 1 attracts. At 0 the derivative vanishes, while at 1 it equals $-1$.
\item[2.13] At $x=0.35$, $(f_A,f_B)=(2.05,2.65)$ and $x'=-0.1365$. At $x=0.65$, the fitnesses are $(2.95,2.35)$ and $x'=0.1365$. Initial motion takes each state farther from the threshold $1/2$.
\item[2.17] At $x^*=3/4$, both fitnesses equal $5/2$, so their weighted mean is $5/2$. Adding 10 to every payoff gives mean fitness $25/2$ and the same equilibrium composition.
\end{answers}

\section*{Chapter 3}
\begin{answers}
\item[3.1] The displayed position is $(0.65,0.15\sqrt3)$. The AB edge lacks C, the BC edge lacks A, and the CA edge lacks B.
\item[3.3] $\bar f=5/4$ and $\p'=(3/8,-5/16,-1/16)$. The three components sum to zero.
\item[3.5] Differentiating the exponential gives $p_i'=p_i(f_i-\bar f)$. Smooth fitnesses on a neighborhood of the compact simplex are bounded there, as is mean fitness, so their finite-time integrals are finite.
\item[3.7] Antisymmetry gives $\p^TM_1\p=0$. Interior stationarity additionally requires $M_1\p=0$, which gives the uniform state. For example, at $(1/2,1/3,1/6)$ the velocity is nonzero.
\item[3.9] Expanding the squared differences gives $\sum_i p_i^2-Q$. Since $1=\sum_i p_i^2+2Q$, this equals $1-3Q$. It vanishes exactly when all coordinates agree.
\item[3.11] At $s=1$, the trace is zero, determinant is $1/3$, and eigenvalues are $\pm i/\sqrt3$. Purely imaginary linearized eigenvalues alone do not establish nonlinear periodicity. The conserved product confines trajectories to smooth closed invariant curves.
\item[3.13] The real part is $0.05$, and the imaginary-part magnitude is $\sqrt3(2.3)/6\approx0.66395$. Small perturbations grow while circulating; their linearized period is about $9.46$.
\item[3.17] Subtract $p_j'/p_j$ from $p_i'/p_i$. The two mean-fitness terms cancel. If the two fitnesses agree throughout a time interval, their probability ratio stays constant there.
\end{answers}

\section*{Chapter 4}
\begin{answers}
\item[4.1] The mean is $5/4$, second moment is $11/4$, and variance is $19/16$.
\item[4.3] The mean is $1/2$, second moment is $9/2$, and variance is $17/4$.
\item[4.5] Since $X^2=X$, the two required expectations are both $p$. Variance is $p-p^2$, maximized at $p=1/2$ with value $1/4$. At either endpoint $X$ is constant.
\item[4.7] The derivative is $0(-0.1)+1(0.05)+3(0.05)=0.20$. Probability is transferred away from the zero reading to both positive readings.
\item[4.9] The derivative is $2.16-2.88t$, so the maximum occurs at $t=3/4$. Its value is $1.65$.
\item[4.11] The empirical law is $(0.475,0.3375,0.1875)$. Its mean is $0.9$, compared with $5/6$ for the specified model.
\item[4.13] The standard deviations are $\sqrt{0.24/100}\approx0.04899$ and $\sqrt{0.24/400}\approx0.02449$.
\item[4.17] Since $p(1-p)\leq1/4$, the probability is at most $1/(4N(0.05)^2)=100/N$. Taking $N\geq1000$ makes this at most $0.1$. This general bound is sufficient and can be conservative.
\end{answers}

\section*{Chapter 5}
\begin{answers}
\item[5.1] The values in nats are $\ln2,\ln4,\ln10$; in bits they are $1,2,\log_210\approx3.3219$.
\item[5.3] The entropy is $(3/2)\ln2\approx1.0397$ nats, or $1.5$ bits, below $\ln3$ nats.
\item[5.5] The nonzero terms contribute $-k(1/k)\ln(1/k)=\ln k$. Exactly the pure laws have entropy zero.
\item[5.7] $L(t)=-\ln3+1.2t$. The target $x=1/2$ corresponds to $L=0$, so $t=\ln3/1.2$.
\item[5.9] Normalize the weights $(0.2e,0.3,0.5e^{-1})$. To four decimal places the state is
\[
(0.5291,0.2919,0.1790).
\]
\item[5.11] $Z'/Z=\sum_i p_ic_i$. Differentiating again gives $\sum_i p_ic_i^2-(\sum_i p_ic_i)^2$. This is a nonnegative variance.
\item[5.15] The state increases to 1. Entropy increases until $x=1/2$, then decreases, tending to zero.
\item[5.17] For example, start from $(0.05,0.05,0.90)$ with constant fitnesses $(1,0,-1)$. Initially selection transfers weight away from the dominant third type, increasing entropy; eventually the uniquely fittest first type dominates and entropy tends to zero. The derivative formula can verify the initial sign.
\end{answers}

\section*{Chapter 6}
\begin{answers}
\item[6.1] The values are $0.130812$ and $0.143841$ nats, respectively.
\item[6.3] If $A=(1/2,1/2,0)$ and $U=(1/3,1/3,1/3)$, then $D(A\|U)=\ln(3/2)$. The reverse divergence is infinite because the third predictive probability is zero.
\item[6.5] $V'(x)=(x-2/5)/[x(1-x)]$. The chain rule gives $\dot V=(x-2/5)(2-5x)=-5(x-2/5)^2$.
\item[6.7] The second derivative along the mixture is $\sum_iq_i(q_i-p_i)^2/M_i(t)^2>0$ for distinct endpoints. Since the first derivative at $t=1$ is zero, it is negative for $t<1$.
\item[6.9] Differentiation gives $\bar f-\sum_iq_if_i$. It is nonpositive when $\sum_iq_if_i\geq\sum_ip_if_i$ at the current state.
\item[6.11] The value and first derivative at $h=0$ are zero. The second derivative is $1/[q(1-q)]$, producing $h^2/[2q(1-q)]$. The first derivative cancels because the two displaced probabilities still sum to one.
\item[6.15] Integrate the displayed derivative. On the segment from $q$ to $x$, $1/[s(1-s)]\geq4$. Accounting for the orientation of the integral on either side of $q$ gives $V(x)\geq2(x-q)^2$.
\item[6.16] On the compact interval, $1/[s(1-s)]\leq1/[\delta(1-\delta)]$, giving the quadratic upper bound. Hence $(x-q)^2\geq2\delta(1-\delta)V$ and $\dot V\leq-2b\delta(1-\delta)V$. Integration gives $V(t)\leq V(0)e^{-2b\delta(1-\delta)t}$ while the trajectory remains in the interval.
\end{answers}

\section*{Chapter 7}
\begin{answers}
\item[7.3] The coefficients are $100/9,16/3,4,100/9$. Interchanging the two Bernoulli outcomes replaces $p$ by $1-p$ and preserves information.
\item[7.5] The exact values are approximately $0.151898$ and $0.100167$. The left-endpoint approximations are $1/6\approx0.166667$ and $0.1$, respectively.
\item[7.7] The arc coordinate is $\theta(t)=\pi/3+t\pi/3$, so $p(t)=\sin^2(\pi/6+t\pi/6)$. Its speed is $\pi/3$ throughout.
\item[7.9] Expanding $(dx+dy)^2/z$ gives the matrix. Its determinant is $1/(xy)+1/(xz)+1/(yz)=1/(xyz)$. The usual two-by-two inverse formula gives the stated inverse after using $x+y+z=1$.
\item[7.11] The mixed product is $-1/60+1/40=1/120$. The squared norms are $49/1200$ and $1/12$. Their product has square root $7/120$, giving the cosine $1/7$.
\item[7.13] The bound is $p(1-p)/N=3/1600=0.001875$. The empirical proportion has exactly that variance.
\item[7.15] The mean fitness is $1.9$, and variance is $0.2(1.1)^2+0.3(-0.9)^2+0.5(0.1)^2=0.49$. The metric calculation gives $0.22^2/0.2+(-0.27)^2/0.3+0.05^2/0.5=0.49$.
\item[7.18] Cancel the arbitrary $v$ to obtain $w=\Phi'/g$. For the given metric and potential, $w=x(1-x)(3-4x)$.
\end{answers}

\section*{Chapter 8}
\begin{answers}
\item[8.1] The differential sends $v_i$ to $v_i/\sqrt{p_i}$, giving the Fisher inner product by an ordinary dot product. Its dot product with $2\sqrt P$ is $2\sum_i v_i=0$.
\item[8.3] The coefficient is $B\approx0.948475$. The central angle is $0.322408$ radians, Fisher distance is $0.644817$, and chord length is $0.642028$. Evaluate the formulas before rounding to keep their small differences visible.
\item[8.5] The half-angle identity gives $\sin(\alpha/2)/\sin\alpha=1/[2\cos(\alpha/2)]$. Squaring the midpoint square-root coordinates gives \eqref{eq:fishermidpoint}, with total numerator $2+2B$.
\item[8.7] The separation is $\pi/6$, the Fisher distance is $\pi/3$, and the chord is $4\sin(\pi/12)=\sqrt6-\sqrt2$.
\item[8.9] Along the mixture the derivative is $R-P$. Substitute it into the Fisher speed and integrate. The geodesic distance is the infimum of admissible path lengths, giving the inequality.
\item[8.11] The square-root geodesic is $(2\cos(\pi t/2),2\sin(\pi t/2),0)$. The probability path is $(\cos^2(\pi t/2),\sin^2(\pi t/2),0)$, of constant Fisher speed $\pi$.
\item[8.13] The coefficient $B$ lies in $[0,1]$, so $2\arccos B\leq\pi$. Equality holds precisely for disjoint supports.
\item[8.17] A simultaneous permutation merely reorders the summands $\sqrt{p_ir_i}$, leaving their sum and the distance unchanged.
\end{answers}

\section*{Chapter 9}
\begin{answers}
\item[9.3] The flow is $x'=x(1-x)(2-5x)$. The equilibria are $0,2/5,1$, with the interior state attracting and the endpoints repelling from the interior.
\item[9.5] The derivative fitnesses are $(2,1,0)$ and their mean is $4/3$. The gradient is $(1/3,-1/9,-2/9)$. Its squared Fisher norm, and ascent rate, is $5/9$.
\item[9.7] With $S=\sum_i p_i^2$, the KL derivative is $1/n-S\leq0$, strictly below zero away from uniform. The divergence from a positive uniform reference tends to infinity at the boundary, so its sublevels are compact in the interior. The strict Lyapunov argument gives convergence.
\item[9.9] The state is approximately $(0.4288,0.3264,0.2448)$. Since the displacement factor is $e^{-t}$, a tenfold reduction takes $t=\ln10$.
\item[9.11] The ambient derivatives are $-\ln(p_i/q_i)-1$. Subtract their weighted mean to obtain the displayed vector field. Log-ratios of $p_i/q_i$ obey $r'=-r$; integrating and normalizing gives \eqref{eq:reverseklflow}.
\item[9.13] Normalize $(\sqrt{1/2},\sqrt{1/3},\sqrt{1/6})$. The resulting probabilities are approximately
\[
(0.4177,0.3411,0.2412).
\]
\item[9.15] A tangent vector has $v_C=-v_A-v_B$, so $\sum_i f_iv_i=(f_A-f_C)v_A+(f_B-f_C)v_B$. For the cyclic field the two cross derivatives are $3$ and $-3$, which disagree.
\item[9.16] Expand the numerator as $p_i(1+\eta f_i+O(\eta^2))$ and the denominator as $1+\eta\bar f+O(\eta^2)$. Multiplying by the reciprocal denominator gives $p_i+\eta p_i(f_i-\bar f)+O(\eta^2)$.
\end{answers}

\section*{Chapter 10}
\begin{answers}
\item[10.1] The recursion is $\E[Z^k]=(k-1)\E[Z^{k-2}]$ for $k\geq2$. Hence the even moments needed are 1 and 3. Odd moments vanish because their integrands are odd and integrable on $\R$.
\item[10.3] $G_{\mu\mu}=\E[Z^2]/\sigma^2=1/\sigma^2$, $G_{\mu\sigma}=\E[Z^3-Z]/\sigma^2=0$, and $G_{\sigma\sigma}=\E[(Z^2-1)^2]/\sigma^2=2/\sigma^2$.
\item[10.5] The norms are $1/3$, $\sqrt2/3$, and $1/\sqrt3$, respectively.
\item[10.7] The standard deviations are 3 and 6, so the distance is $\sqrt2\ln2$. A constant-speed path is $\mu(t)=2$, $\sigma(t)=3\cdot2^t$.
\item[10.9] The center is 0, radius is $\sqrt3$, distance is $\sqrt2\arcosh5$, and maximum height is $\sqrt3$. The height is the intermediate standard deviation.
\item[10.11] The directed values are $\ln2-1/4\approx0.443147$ and $2-\ln2\approx1.306853$. The Fisher distance is $\sqrt2\arcosh(11/8)\approx1.189381$.
\item[10.13] Compute $c,R$ from \eqref{eq:circlecenter} and endpoint angles $\varphi_j\in(0,\pi)$. Let $w_j=\ln\tan(\varphi_j/2)$ and $w(t)=(1-t)w_1+tw_2$. Then $\varphi(t)=2\arctan(e^{w(t)})$, $u=c+R\cos\varphi(t)$, $v=R\sin\varphi(t)$. Its speed is $\sqrt2|w_2-w_1|$.
\item[10.15] For every truncated portion of the curve, $L_F\geq\sqrt2|\ln(\sigma(t)/\sigma(0))|$. The right-hand side diverges as $\sigma(t)\to0$, regardless of mean motion.
\item[10.18] Completing the square gives the displayed overlap formula. For the stated pair it is $e^{-2}$, so the full-density spherical distance is $2\arccos(e^{-2})$, whereas the intrinsic Gaussian distance is $\sqrt2\arcosh5$. The latter is larger because it restricts the connecting path to Gaussian laws.
\end{answers}


\renewcommand{\bibname}{References and Further Reading}
\begin{thebibliography}{9}
\addcontentsline{toc}{chapter}{References and Further Reading}
\bibitem{course}
\emph{Introduction to Information Geometry}. Interactive readings, demonstrations, and investigations.
\url{https://bobovski66.github.io/IntroInfoGeom/}.

\bibitem{hofbauer}
J. Hofbauer and K. Sigmund, \emph{Evolutionary Games and Population Dynamics}. Cambridge University Press, 1998.
\url{https://doi.org/10.1017/CBO9781139173179}.
An extended treatment of replicator dynamics, stability, and evolutionary interaction models.

\bibitem{shannon}
C. E. Shannon, ``A mathematical theory of communication,'' \emph{Bell System Technical Journal} \textbf{27} (1948), 379--423 and 623--656.
\url{https://doi.org/10.1002/j.1538-7305.1948.tb01338.x}.
The foundational development of entropy and information in communication theory.

\bibitem{shahshahani}
S. Shahshahani, \emph{A New Mathematical Framework for the Study of Linkage and Selection}. Memoirs of the American Mathematical Society, no. 211, 1979.
The geometric treatment of selection dynamics associated with the simplex metric used here.

\bibitem{amari}
S.-I. Amari, ``Natural gradient works efficiently in learning,'' \emph{Neural Computation} \textbf{10}, no. 2 (1998), 251--276.
\url{https://doi.org/10.1162/089976698300017746}.
A primary source for natural-gradient methods in statistical learning.

\bibitem{nielsen}
F. Nielsen, ``An elementary introduction to information geometry,'' \emph{Entropy} \textbf{22}, no. 10 (2020), article 1100.
\url{https://doi.org/10.3390/e22101100}.
A broader introduction to the differential-geometric structures of statistical families.
\end{thebibliography}

\cleardoublepage\phantomsection\addcontentsline{toc}{chapter}{Index}\printindex
\end{document}
